\documentclass[reprint,amsmath,amssymb,aps,prd,longbibliography,floatfix]{revtex4-2}

\usepackage{graphicx}
\usepackage{dcolumn}
\usepackage{bm}
\usepackage{hyperref}
\usepackage{booktabs}
\usepackage{siunitx}
\usepackage{multirow}
\providecommand{\degree}{^{\circ}}
\sisetup{
    range-units = single,
    separate-uncertainty = true,
    table-align-uncertainty = true,
}

\begin{document}

\title{Zero-parameter structural computation of Standard Model observables and constants and
a dark-matter tower from $E_8 \times E_8$ Coxeter geometry on a Schoen Calabi--Yau threefold}

\author{Eliot Gable}
\email{10015996+egable@users.noreply.github.com}
\affiliation{Independent Researcher}

\date{\today}

\begin{abstract}
\noindent\textbf{What this paper does.}
We construct the Standard Model fermion mass spectrum, gauge couplings, mixing angles (CKM, PMNS), CP phases, dark-matter mass tower, Higgs boson mass, Higgs vacuum expectation value, and Newton's gravitational constant from a small set of fundamental equations acting on the fixed geometry of $E_8 \times E_8$ heterotic compactification on the Schoen $\mathbb{Z}_3 \times \mathbb{Z}_3$ Calabi--Yau threefold quotient $X/(\mathbb{Z}_3 \times \mathbb{Z}_3)$, with the McKay correspondence on the binary icosahedral group $2I$ and Wilson $\mathbb{Z}_3 \times \mathbb{Z}_3$ symmetry breaking. The construction has no PDG inputs and no adjustable quantities.

\noindent\textbf{Why it matters.}
Approximately twenty-six SM observables are predicted to ppm-or-better precision: $m_e$ at \SI{17}{ppm}, $m_\mu$ at \SI{30}{ppb}, $m_\tau$ at \SI{0.15}{ppm}, $\sin^2 \theta_W = 2/9 + \Delta_\varphi/h$ at \SI{198}{ppm}, $\alpha_s(M_Z) = 849/7192$ at \SI{405}{ppm}, the dark-matter abundance $\Omega_{\rm DM} h^2 = 0.119982$ vs.\ Planck~2018 $0.1200(10)$, and Newton's $G$ to \SI{2.5}{ppm} vs.\ CODATA~2022 --- all from the geometry alone. The CP phases $\delta_{\rm CP} = \pi + \pi/10$ and $\delta_{\rm CKM} = 11\pi/30$ are predicted as rational multiples of $\pi$. Normal neutrino mass ordering is selected. A conditional geometric route to vanishing strong CP via the $\mathbb{Z}_3 \times \mathbb{Z}_3$ quotient projection avoids invoking a Peccei--Quinn axion.

\noindent\textbf{What is open and not claimed.}
The construction is a strong-selection structural identification, not yet a forced derivation: the substrate-physical projection onto the heterotic moduli of $X/(\mathbb{Z}_3 \times \mathbb{Z}_3)$ is not closed in this paper. The $\sigma$-discrimination Schoen-vs-Tian--Yau result ($6.92\sigma$ Tier~0; bundle-construction-specific) is reported as supporting evidence pending the converged PDG-fermion-mass Yukawa-cohomology pipeline. Dark-sector chirality assignments and several closure-rule numerators are encoded structural-rule-consistent choices whose generative derivation is deferred.
\end{abstract}

%% ============================================================================
%% PLAIN-TEXT ABSTRACT (copy-paste-friendly; no LaTeX macros)
%%
%% This block is wrapped in \iffalse ... \fi so LaTeX skips it entirely; it
%% does NOT appear in the rendered PDF. Use this version for the arXiv abstract
%% field, journal submission portal, Zenodo metadata, conference abstracts,
%% Google Scholar, social media, etc. -- anywhere LaTeX macros do not render.
%% ============================================================================
\iffalse

WHAT THIS PAPER DOES. We construct the Standard Model fermion mass spectrum, gauge couplings, mixing angles (CKM, PMNS), CP phases, dark-matter mass tower, Higgs boson mass, Higgs vacuum expectation value, and Newton's gravitational constant from a small set of fundamental equations acting on the fixed geometry of E_8 x E_8 heterotic compactification on the Schoen Z_3 x Z_3 Calabi-Yau threefold quotient X/(Z_3 x Z_3), with the McKay correspondence on the binary icosahedral group 2I and Wilson Z_3 x Z_3 symmetry breaking. The construction has no PDG inputs and no adjustable quantities.

WHY IT MATTERS. Approximately twenty-six SM observables are predicted to ppm-or-better precision: m_e at 17 ppm, m_mu at 30 ppb, m_tau at 0.15 ppm, sin^2 theta_W = 2/9 + Delta_phi/h at 198 ppm, alpha_s(M_Z) = 849/7192 at 405 ppm, the dark-matter abundance Omega_DM h^2 = 0.119982 vs. Planck 2018 0.1200(10), and Newton's G to 2.5 ppm vs. CODATA 2022 -- all from the geometry alone. The CP phases delta_CP = pi + pi/10 and delta_CKM = 11 pi / 30 are predicted as rational multiples of pi. Normal neutrino mass ordering is selected. A conditional geometric route to vanishing strong CP via the Z_3 x Z_3 quotient projection avoids invoking a Peccei-Quinn axion.

WHAT IS OPEN AND NOT CLAIMED. The construction is a strong-selection structural identification, not yet a forced derivation: the substrate-physical projection onto the heterotic moduli of X/(Z_3 x Z_3) is not closed in this paper. The sigma-discrimination Schoen-vs-Tian-Yau result (6.92 sigma Tier 0; bundle-construction-specific) is reported as supporting evidence pending the converged PDG-fermion-mass Yukawa-cohomology pipeline. Dark-sector chirality assignments and several closure-rule numerators are encoded structural-rule-consistent choices whose generative derivation is deferred.

\fi
%% ============================================================================
%% END PLAIN-TEXT ABSTRACT
%% ============================================================================

\maketitle

\tableofcontents

\section{Introduction}
\label{sec:intro}

The Standard Model (SM) of particle physics, despite its empirical
success, contains twenty-six parameters whose values are not predicted
by any underlying principle but determined empirically from
observation~\cite{PDG2024}. These include the three gauge couplings,
the Higgs sector parameters (vacuum expectation value and self-coupling),
nine charged-fermion Yukawa couplings, three neutrino mass scales, four
quark-sector mixing parameters (CKM), four lepton-sector mixing
parameters (PMNS), and the QCD vacuum angle $\theta_{\rm QCD}$. Beyond
the SM, the existence and mass scale of dark matter constitutes one of
the most prominent unresolved puzzles in contemporary
physics~\cite{Bertone2018}.

This paper presents a structural computation, with no fitted
continuous parameter, of all twenty-six SM parameters, Newton's
gravitational constant, the Planck mass, and a discrete dark-matter
mass tower with associated decay constant, couplings, and cosmological
abundance, given a fixed set of first-principles starting assumptions
(Table~\ref{tab:starting-assumptions}). Per-particle Coxeter chain
indices, sign assignments, sector and cycle placements, and rational
closure rules arise as structural outputs of those starting
commitments rather than as independently varied parameters; the
audit is in Sec.~\ref{subsec:discrete-audit}. PDG-matching residuals are
predictions versus measurement, not fits. The framework's predictive
path takes \emph{zero} measured masses as input. The Higgs vacuum
expectation value is computed in closed form from the $E_8$ Coxeter
invariants alone via the structural identity
$v_C = (\dim - 2)/(1 - \delta_W) \cdot (1 + 5\,\alpha_\tau^3/8)$
(see substrate\_mass\_spectrum.py:402; inputs $\dim=248$ at line~160,
$h=30$ at line~161, $m_{\max}=29$ at line~163, $e_6=19$ at line~399, and
the derived $\delta_W = 1/((h+m_{\max})e_6) = 1/1121$ at line~400 and
$\alpha_\tau = 1/(\dim+h) = 1/278$ at line~188), giving
$v_C = 246.21965$~GeV. The gateway tension scale is then
$h_{\rm eff} = v_C / R_v$ where $R_v$ is the structural Higgs-vev
ratio, and every other mass is the corresponding chain-specific
structural ratio times $h_{\rm eff}$ (\textit{$\varphi$-chain leptons,
$\sqrt{2}$-chain quarks/EW bosons, $\alpha_\tau$-tower neutrinos and
dark-matter modes}). As an internal-arithmetic self-consistency check,
alternatively anchoring the framework on the CODATA electron mass alone
or muon mass alone produces predictions that agree at
\SI{0.0298}{ppm} ($\sim 30$~ppb), set by sub-harmonic-series truncation;
both fermion-anchored modes agree with the structural-vev anchor to
within \SI{0.12}{ppm} on every observable. This two-anchor
\emph{diagnostic} validates the framework's internal arithmetic but is
not the predictive path.

\paragraph{Starting assumptions of this paper.}
Table~\ref{tab:starting-assumptions} enumerates the complete list of
first-principles commitments on which the numerical demonstration of
this paper rests. The list contains no continuous fitted parameter.
Every structural integer, sign, sub-harmonic factor, and closure rule
that appears in the mass equations of Secs.~\ref{sec:gateway}--\ref{sec:gravity}
is either an immediate consequence of these commitments or, where so
flagged, an independent structural rule whose derivation is the subject
of the companion paper~\cite{UnreasonableEffectiveness}.

\begin{table*}[t]
\centering
\caption{Starting assumptions of the present paper. The numerical demonstration in Secs.~\ref{sec:gateway}--\ref{sec:gravity} takes these as given. The argument that this list is minimal --- i.e.\ that one would expect the substrate-physical principles to select these and not others --- is the subject of the companion paper~\cite{UnreasonableEffectiveness} and is not load-bearing for the present demonstration.}
\label{tab:starting-assumptions}
\begin{tabular}{ll}
\toprule
Starting assumption & Status / role in the present paper \\
\midrule
\parbox[t]{0.32\textwidth}{Heterotic $E_8 \times E_8$ string theory~\cite{GHMR1985}} & \parbox[t]{0.62\textwidth}{Standard heterotic literature; assumed.} \\[0.5ex]
\parbox[t]{0.32\textwidth}{Schoen Calabi--Yau threefold quotient $X/(\mathbb{Z}_3 \times \mathbb{Z}_3)$~\cite{Schoen1988,BraunHeOvrutPantev2005,DonagiOvrut2005}} & \parbox[t]{0.62\textwidth}{Empirically identified within the framework's narrowed candidate set; see Sec.~\ref{sec:cy3-substrate-identification}.} \\[0.5ex]
\parbox[t]{0.32\textwidth}{$E_8$ Coxeter spectrum (rank $8$, $h=30$, exponents $\{1,7,11,13,17,19,23,29\}$, $|\Phi|=240$)} & \parbox[t]{0.62\textwidth}{Standard Lie-algebraic data~\cite{Humphreys1990}; assumed.} \\[0.5ex]
\parbox[t]{0.32\textwidth}{McKay correspondence~\cite{McKay1980}; binary icosahedral McKay-dual of $E_8$} & \parbox[t]{0.62\textwidth}{Standard; the icosahedral order-$5$ assignment as the framework's harmonic seed is a structural commitment. Companion paper argues for its substrate-physical inevitability.} \\[0.5ex]
\parbox[t]{0.32\textwidth}{Wilson-line breaking $E_8 \to G_{\rm SM}$ along $\pi_1(X/\Gamma) = \mathbb{Z}_3 \times \mathbb{Z}_3$; phenomenological lepton/quark sector split via $\varphi$-chain (5-fold icosahedral) and $\sqrt{2}$-chain (4-fold octahedral) harmonic recursions} & \parbox[t]{0.62\textwidth}{Standard heterotic chain $E_8 \supset E_6 \supset SO(10) \supset SU(5) \supset G_{\rm SM}$. The $\varphi$/$\sqrt{2}$ chain split is a framework structural assignment matching observed mass hierarchies with small-integer Coxeter-step assignments; a derivation from a specific Wilson-line representation on $X/\Gamma$ is part of the open program (Sec.~\ref{sec:open-items}).} \\[0.5ex]
\parbox[t]{0.32\textwidth}{Substrate-tension reading of the Higgs mechanism (motivated by Anderson's superconducting precursor~\cite{Anderson1962}; similar in spirit to but independent of Volovik's emergent-vacuum program~\cite{Volovik2003})} & \parbox[t]{0.62\textwidth}{Interpretive frame for the Higgs vacuum; the $v$ ladder in Sec.~\ref{sec:higgs} is a closed-form structural identity within this reading.} \\[0.5ex]
\parbox[t]{0.32\textwidth}{Coldea \textit{et al.}~\cite{Coldea2010} confirmation of the Zamolodchikov $E_8$ spectrum~\cite{Zamolodchikov1989} in CoNb$_2$O$_6$} & \parbox[t]{0.62\textwidth}{Empirical observation of the $E_8$ root system in a 1+1D condensed-matter context; adopted as motivation for the framework's $E_8$ commitment, not as a direct measurement of the heterotic gauge $E_8$. Not load-bearing for the numerical computations of the present paper.} \\
\bottomrule
\end{tabular}
\end{table*}

\paragraph{Structural-output audit.}
\label{subsec:discrete-audit}
The construction generates discrete data --- per-particle Coxeter
chain indices, sign assignments, sector and cycle placements, and a
family of closure rules --- as \emph{structural outputs} of the
starting commitments of Table~\ref{tab:starting-assumptions}, not as
independently varied parameters. They are not free: changing any one
of them in isolation would generally be inconsistent with the
underlying structural rules (the binary-icosahedral McKay seed, the
$\varphi$- and $\sqrt{2}$-chain harmonic-recursion split, the
boundary capacity $h+m_{\max}=59$, the $E_8$ Coxeter spectrum, and
the Wilson-line-broken visible chain). They become arbitrary only
under the counterfactual in which the structure is removed. We list
the discrete data explicitly so the reader can audit by inspection
which items are derived in closed form within this paper and which
are encoded as structural-rule-consistent choices whose generative
derivation is part of the open program.
\begin{itemize}
\item \emph{Per-particle Coxeter chain index} $k_X$ for each charged
fermion, EW boson, and the Higgs vev (App.~\ref{app:assignments}):
$k_e=0$, $k_\mu=11$, $k_\tau=17$ on the $\varphi$-chain (lepton
positions at integer $E_8$ exponents); $k_u=4$, $k_d=6.5$, $k_s=15$,
$k_c=22.5$, $k_b=26$, $k_t=36.5$, $k_W=34.5$, $k_v=1132/30$,
$k_H = k_v - 2$ on the $\sqrt{2}$-chain. The Higgs-vev position
$k_v = h + \mathrm{rank} - \mathrm{rank}/h$ is given in closed form
(Eq.~\ref{eq:kv-derived}); the remaining $k_X$ values are encoded
as positions consistent with the integer/half-integer cycle-1 and
cycle-2 structure of the chain. A generative derivation from the
specific Wilson-line representation on $X/\Gamma$ is part of the
open program (Sec.~\ref{sec:open-items}).
\item \emph{Per-quark sign assignments} $\chi_q, \lambda_q \in \{-1,
0, +1\}$ in the quark mass formula (Table~\ref{tab:quarksigns}):
encoded as a structural pattern across the six quarks; their
derivation from a structural-rule generator (chirality flip across
the harmonic threshold $h/2$, per-step-loss sign across cycle
boundaries) is part of the open program.
\item \emph{Sector chain-type assignment} ($\varphi$ vs $\sqrt{2}$):
leptons and neutrinos to the $\varphi$-chain; quarks and electroweak
bosons to the $\sqrt{2}$-chain. This pairing reflects the framework's
McKay-icosahedral and McKay-octahedral seed assignments; an explicit
derivation from a Wilson-line representation on $X/\Gamma$ is open.
\item \emph{Cycle-1 vs cycle-2 placement} for each particle: encoded
in the structural pattern of the chain
(App.~\ref{app:assignments}).
\item \emph{Closure rules} $\delta_H = 1/(h+m_{\max}) = 1/59$ and
$\delta_W = 1/(59 \cdot e_6) = 1/(59 \cdot 19) = 1/1121$ are given
in closed form from the boundary capacity $h+m_{\max}$ and the
$E_8$ exponent set; the gateway sub-harmonic factor $1/(4(\dim-1))$
admits a dual-sector / two-orientation structural reading whose
derivation from a heterotic-anomaly index or a $\mathbb{Z}_3
\times \mathbb{Z}_3$ quotient trace on the Schoen threefold is part
of the open program.
\item \emph{Sub-harmonic correction structures} ($\alpha_v$,
$\alpha_H$, $\alpha_W$, $\alpha_t$, the rank-locked sub-harmonic on
$m_5$, and the $5\varphi^2 \alpha_\tau$ correction on $\nu_3$) are
encoded as structural-rule combinations of $\{\dim, h,
\mathrm{rank}, m_{\max}, e_j, \varphi\}$. A generative derivation
of each from the compactification data is open.
\item \emph{Other rational closures} include $\sin^2 \theta_W = 2/9 +
\Delta_\varphi/h$ and $\alpha_s(M_Z) = (m_{\max}^2 + \mathrm{rank})/
(m_{\max} \cdot \dim) = 849/7192$. The leading $2/9$ in
$\sin^2\theta_W$ and the combination $m_{\max}^2 + \mathrm{rank}$ in
$\alpha_s$ are framework structural choices; they are not standard
$E_8$ Lie invariants and their derivation is open.
\end{itemize}
None of the items above is varied independently of the framework's
structural rules. Items derived in closed form within this paper
(boxed equations) are structurally forced; items flagged as encoded
or as ``framework structural assignment'' are consistent with the
framework's rules but their generative derivation from the
compactification data is the subject of the further-exploration
program (Sec.~\ref{sec:open-items}).

\paragraph{Epistemic posture.}
From a 4D observer's vantage the starting assumptions of
Table~\ref{tab:starting-assumptions} cannot be \emph{proven}: any
purported derivation of the substrate-physical content from inside
the framework's macroscopic 4D regime is necessarily a consistency
check, not a proof. The standard a 4D observer can apply is therefore
\emph{empirical consistency} --- the assumptions are evaluated by the
predictions they generate. This paper performs that evaluation by
computing the SM observables and comparing to measurement. We make no
claim in this paper that the starting assumptions are forced from
first principles; we make the weaker, sharper claim that, taking them
as given, the SM mass spectrum and mixing parameters compute directly
with the residuals reported in Table~\ref{tab:starting-assumptions}'s
Sec.~\ref{sec:results} downstream. The companion
paper~\cite{UnreasonableEffectiveness} addresses the separate
question of whether the starting assumptions are minimal and
substrate-physically expected; that argument is independent of the
present demonstration and not relied upon here.

\paragraph{Roadmap.}
Section~\ref{sec:framework} fixes the framework as commitment ---
heterotic structure on the Schoen quotient, $E_8$ Coxeter spectrum,
McKay-icosahedral seed, sub-Coxeter decomposition, substrate-tension
Higgs reading. Section~\ref{sec:cy3-substrate-identification} reports
the empirical Schoen-versus-Tian--Yau discrimination, the Tian--Yau
Yukawa-channel structural exclusion, and the catalogue uniqueness
sweep at Hodge $(3,3)$, with explicit hedges for the open
methodological items. Section~\ref{sec:gateway} presents the gateway
formula and the direct computation of the electron mass.
Sections~\ref{sec:fermions}--\ref{sec:gravity} present the direct
computations of charged-fermion masses, the Higgs sector, electroweak
gauge couplings, the strong-CP angle, neutrino masses and mixing, CKM
mixing, the dark-matter Coxeter tower with cosmological abundance,
and Newton's $G$. Section~\ref{sec:results} presents the master
results tables. Section~\ref{sec:predictions} enumerates the
falsifiable predictions. Section~\ref{sec:open-items} enumerates the
open methodological items, the framework-level commitments deferred
to the companion paper, and the further-exploration program.
Section~\ref{sec:discussion} discusses parameter-count semantics and
relation to prior work.

\section{The substrate framework}
\label{sec:framework}

\subsection{Heterotic $E_8 \times E_8$ on the Schoen quotient}

The framework treats heterotic $E_8 \times E_8$ string
theory~\cite{GHMR1985} compactified on the Schoen Calabi--Yau threefold
quotient $X/(\mathbb{Z}_3 \times \mathbb{Z}_3)$ as the structural
substrate. The Schoen threefold~\cite{Schoen1988} is the smooth,
simply-connected Calabi--Yau fiber product
$X = B \times_{\mathbb{P}^1} B'$ of two rational elliptic surfaces
over $\mathbb{P}^1$, with Hodge numbers $h^{1,1}(X) = h^{2,1}(X) = 19$
and Euler characteristic $\chi(X) = 0$. In the heterotic
standard-model construction, suitably chosen $B, B'$ admit a
freely acting $\Gamma = \mathbb{Z}_3 \times \mathbb{Z}_3$ symmetry,
and the quotient $X/\Gamma$ acquires
$\pi_1(X/\Gamma) = \Gamma = \mathbb{Z}_3 \times \mathbb{Z}_3$
(while $X$ itself remains simply-connected), enabling Wilson-line
breaking of one $E_8$ factor to Standard-Model gauge
structure~\cite{BraunHeOvrutPantev2005,
DonagiOvrut2005}.

\paragraph{Hodge numbers and generation count.}
The Schoen quotient $X/\Gamma$ has Hodge numbers
$h^{1,1}(X/\Gamma) = h^{2,1}(X/\Gamma) = 3$, computed as the
dimensions of the $\Gamma$-invariant subspaces of $H^{1,1}(X)$ and
$H^{2,1}(X)$ via the character formula
$\dim H^{p,q}(X)^\Gamma = (1/|\Gamma|)\sum_{g \in \Gamma}
\mathrm{tr}(g^* | H^{p,q}(X))$; see
\cite{BraunHeOvrutPantev2005} for the explicit
representation-theoretic decomposition. The Euler characteristic
of the cover, $\chi(X) = 2[h^{1,1}(X) - h^{2,1}(X)] = 0$, divides
under the freely-acting $\Gamma$ to give $\chi(X/\Gamma) = 0$ in
absolute Euler-characteristic terms; the
\emph{net chirality index} relevant for the standard heterotic
generation count is set instead by the chiral Atiyah--Singer
index of the bundle $V$ chosen to break $E_8$ via Wilson lines,
$|n_g - \bar n_g| = |\mathrm{Index}(V)|/|\Gamma|$ with
$\mathrm{Index}(V) = \int_X \mathrm{ch}(V) \cdot \mathrm{Td}(TX)$
(equivalently $(1/2)\int_X c_3(V)$ when $c_1(V) = 0$, by the
standard Hirzebruch--Riemann--Roch reduction on a CY3). For the
BHOP rank-4 $SU(4)$ extension bundle on Schoen,
$\mathrm{Index}(V_{\rm BHOP}) = -27$ (BHOP-2005 Eq.~88; with
$c_3/2 = -9\,\tau_1\tau_2^{\,2}$ per BHOP-2005 Eq.~97 and
$\int_X \tau_1\tau_2^{\,2} = 3$, giving $\mathrm{Index} = -27$),
and the standard Braun--He--Ovrut--Pantev model thus yields
$|n_g - \bar n_g| = 27/9 = 3$ chiral generations as required by
the visible Standard Model. The present framework assumes (rather
than derives) the bundle commitment, and adopts the BHOP-2005
published index numerics.

\paragraph{Heterotic anomaly cancellation.}
Heterotic $E_8 \times E_8$ on a Calabi--Yau threefold requires
the modified Bianchi identity
$dH = (\alpha'/4)\,\bigl(\text{tr}\,R \wedge R
- \text{tr}\,F \wedge F\bigr)$
to admit a globally defined $H$, equivalent to the cohomological
condition
\begin{equation}
\text{ch}_2(V) \;=\; \text{ch}_2(TX) \;+\; [W],
\label{eq:heterotic-bianchi}
\end{equation}
where $\text{ch}_2$ is the second Chern character, $V$ is the
visible-sector gauge bundle, $TX$ is the holomorphic tangent
bundle of the Schoen threefold, and $[W]$ is an effective
$5$-brane class in $H^4(X, \mathbb{Z})$ (vanishing in the standard
embedding $V = TX$). The present framework assumes that the
bundle choice underlying the three-generation count is consistent
with Eq.~\eqref{eq:heterotic-bianchi} for some admissible
$5$-brane class $[W]$, as is the case in the
Braun--He--Ovrut--Pantev construction~\cite{BraunHeOvrutPantev2005}.
A first-principles identification of the specific
$[W]$ used by the framework, together with the
\emph{slope-stability} ($\mu$-stability) of $V$ required by the
Donaldson--Uhlenbeck--Yau theorem~\cite{DonaldsonStability,UhlenbeckYau}
to admit a Hermitian--Yang--Mills connection on the quotient
$X/\Gamma$~\cite{DonagiOvrut2005}, is part of the open program
(Sec.~\ref{sec:open-items}). The Donaldson-balanced metric
pipeline used in Sec.~\ref{sec:cy3-substrate-identification} is
constructed assuming polystability of $V$; the empirical
convergence of the Donaldson algorithm on the candidate bundles is
evidence-for, but not proof-of, polystability on the quotient.

\paragraph{Substrate dimensionality and the 10D embedding.}
The discrete substrate update rule developed in the companion
paper~\cite{UnreasonableEffectiveness} is specified directly in
$3+1$ macroscopic dimensions, while the heterotic $E_8 \times E_8$
string lives in $9+1$. The framework's
operational stance is that the six internal dimensions are
\emph{algebraically encoded} in the substrate by the
$E_8 \times E_8$ structure of the chain construction (in
particular, the Coxeter-chain positions $k_X$ and the
sector-modifier $\sqrt{2}$- vs.\ $\varphi$-chain assignments
substitute for direct propagation in the compactified
directions): the substrate is treated as $3+1$D throughout, with
the Calabi--Yau / heterotic structure entering only through the
algebra controlling the mass spectrum and gauge structure. This
is a \emph{framework choice}, not a derivation: a substrate-physical
realisation in which the $3+1$D macroscopic substrate is a
boundary or low-energy slice of an underlying $9+1$D substrate
(analogous in spirit to brane-world or warped-compactification
constructions) is consistent with the present construction but
not committed to here. In particular, no claim is made about the
size, geometry, or stabilisation of the six internal directions
beyond what is needed to identify the algebraic content
$E_8 \times E_8$ on $X/\Gamma$.

One $E_8$ factor contains the visible Standard Model gauge structure
after compactification. The second $E_8$ factor is the dark sector,
communicating with the visible sector only gravitationally except
through identified portal couplings.

\paragraph{Visible $E_8 \to$ SM symmetry-breaking chain.}
The visible $E_8$ is broken to the Standard-Model gauge group
$G_{\rm SM} = SU(3)_{\rm color} \times SU(2)_L \times U(1)_Y$
through the BHOP $SU(4)$-structure-group chain
\begin{equation}
\begin{aligned}
E_8 \;&\supset\; SO(10) \times SU(4) \;\supset\; SU(5) \times U(1) \\
    \;&\supset\; G_{\rm SM} \times U(1)_X,
\end{aligned}
\end{equation}
with the final descent realised by Wilson lines along generators
of $\pi_1(X/\Gamma) = \mathbb{Z}_3 \times \mathbb{Z}_3$. The bundle
$V$ committed here is the rank-4 $SU(4)$ extension construction of
Braun--He--Ovrut--Pantev~\cite{BHOPVectorBundle2006} \S 6 (an
extension of two semi-stable summands rather than a line-bundle
sum), with chiral index
$\mathrm{Index}(V_{\rm BHOP}) = \int_{\widetilde X} \mathrm{ch}(V_{\rm BHOP})
\cdot \mathrm{Td}(T\widetilde X) = -27$ on the cover (BHOP-2005 Eq.~88;
$c_3 / 2 = -9\,\tau_1\tau_2^{\,2}$ per BHOP-2005 Eq.~97), yielding
three generations on the free $\Gamma = \mathbb{Z}_3 \times \mathbb{Z}_3$
quotient via the Atiyah--Singer / Hirzebruch--Riemann--Roch reduction
$N_{\rm gen} = |\mathrm{Index}(V_{\rm BHOP})| / |\Gamma| = 27/9 = 3$.
The full heterotic
Bianchi anomaly cancellation
$c_2(TX) - c_2(V_{\rm BHOP}) - c_2(V'_{\rm BHOP}) = 6\,\tau_1^2$
(BHOP Eq.~99) is satisfied by a positive-tension $5$-brane wrapping
$\tau_1^2$, and Mumford--Takemoto polystability is established on
the documented K\"ahler region
$\{t_a \in [0.25, 8.0],\; t_a/t_b \in [0.25, 4.0]\}$. Both checks
are programmatically verified upstream of this paper at commits
\texttt{a3317377} (anomaly cancellation) and \texttt{a4d1bc72}
(polystability). The matter content is delivered by $H^1(X/\Gamma,
V_{\rm BHOP})$ decomposing into three copies of the $\mathbf{16}$
of $SO(10)$ under the $\mathbb{Z}_3 \times \mathbb{Z}_3$ Wilson-line
characters, with right-handed neutrinos arising automatically as
the singlet $\mathbf{1}$ in $\mathbf{16} = \mathbf{10} +
\bar{\mathbf{5}} + \mathbf{1}$. This commitment supersedes the
wave-28-era rank-3 $SU(3)$ ``non-standard'' bundle placeholder,
which was retracted in the audit at
\texttt{p\_stage3\_bundle\_physics\_audit.md} (commit
\texttt{0cad4b0b}) for failing the full Bianchi identity, the
$N_{\rm gen} = 3$ Atiyah--Singer count (the rank-3 placeholder
did not reproduce the BHOP-2005 Eq.~88 index of $-27$), polystability
wiring, and the
$3{:}3{:}3$ Wilson-character decomposition of $H^1$. Alternative
chains (the Donagi--Ovrut $SU(5)$-structure-group route, or
$E_8 \supset SO(16)$) admit consistent compactifications but are
not the framework's commitment. The framework's mass formulas use
$E_8$ Coxeter-spectrum data and Wilson-line-broken chain
assignments that are largely independent of which standard chain
is selected; the per-particle Coxeter-step assignments listed in
App.~\ref{app:assignments} are framework structural commitments
that the BHOP $SU(4)$ identification anchors to a specific
heterotic bundle.

\paragraph{Substrate-physical reading of the rank-4 internal symmetry
\textmd{(I, suggestive; not a BHOP-2005 commitment)}.}
The rank-4 $SU(4)$ here is the \emph{bundle structure group} of
$V_{\rm BHOP}$ acting on holomorphic sections of $V$ on the cover,
\emph{not} the gauged $SU(4)_C$ of Pati--Salam~1974~\cite{PatiSalam1974}
acting on Standard-Model fermion 4-tuples. The framework offers a
\emph{suggestive} substrate-physical reading in the spirit of
Pati--Salam $SU(4)_C$ colour-lepton unification --- the four internal
modes of the bundle treated as three colour modes plus a fourth
``lepton-number'' mode --- because both groups are $SU(4)$ and both
relate to lepton/quark unification. However, the identification of
the bundle structure group with a 4D gauged $SU(4)_C$ is a non-trivial
\emph{additional} structural step that BHOP-2005 \S 6 does not
assume: BHOP's $SU(4)_{\rm bundle}$ acts on holomorphic sections,
not on Standard-Model fermion fields, and the GUT chain
$E_8 \to SO(10) \times SU(4)_{\rm bundle} \to SU(5) \times U(1) \to
G_{\rm SM} \times U(1)_X$ has no Pati--Salam stage in the gauge-symmetry-
breaking ladder. The substrate-physical mechanism realising any
bundle$\to$gauge identification is currently an open structural
step. The framework's empirical content is independent of the
identification: chiral generations descend from the
$\mathbb{Z}_3 \times \mathbb{Z}_3$ Wilson characters acting on
$H^1(X/\Gamma, V_{\rm BHOP})$ regardless of the Pati--Salam reading,
and right-handed neutrinos are automatic from the $\mathbf{16}$ of
$SO(10)$ (the singlet $\mathbf{1}$ in
$\mathbf{16} = \mathbf{10} + \bar{\mathbf{5}} + \mathbf{1}$). The
rank-4 commitment is preferred over the wave-28 rank-3 $SU(3)$
placeholder because the rank-4 extension bundle satisfies the BHOP-2005
Eq.~88 index of $-27$, full Bianchi anomaly closure, polystability,
and the $\mathbf{16}$-decomposition required for right-handed
neutrinos --- not because the Pati--Salam reading is structurally
forced. The residual $U(1)_X$ surviving Wilson-line breaking is
anomalous and is broken by the standard heterotic Green--Schwarz
mechanism (anomalous-$U(1)$ Stueckelberg coupling to the dilaton/
B-field universal axion); the substrate-physical mechanism for
this GS cancellation is an open step.

\subsection{Empirical CY3 substrate identification}
\label{sec:cy3-substrate-identification}

The framework's substrate-physical forcing chain selects the Hodge $(3,3)$, $\pi_1 = \mathbb{Z}_3 \times \mathbb{Z}_3$ branch on which Schoen $\mathbb{Z}_3 \times \mathbb{Z}_3$ sits. The structurally-closest distinct three-generation Calabi--Yau threefold satisfying $c_1 = 0$ and admitting a free finite-group action with Standard-Model-compatible Wilson breaking is the Tian--Yau quotient TY/$\mathbb{Z}_3$ (Tian--Yau 1987~\cite{TianYau1987}, Hodge $(h^{1,1}, h^{2,1}) = (6,9)$, $\pi_1 = \mathbb{Z}_3$, $\chi = -6$). Tab.~\ref{tab:cy3-substrate-identification} records the structural distinction.

To convert this from a literature-driven choice into an empirical identification, we ran two independent physics-channel discriminators on the Schoen and Tian--Yau Donaldson-balanced metrics: a quark-chain channel and a lepton-chain channel. The pipeline as currently implemented computes, for each candidate, a Galerkin-residual sweep at successive test-degrees $k_\mathrm{test} \in \{2,3,4,5\}$ of the discrepancy between the metric-Laplacian eigenvalues (projected onto the framework's $\sqrt{2}$-chain integer/half-integer positions for the quark channel and the $\varphi$-chain integer-exponent positions for the lepton channel) and the structural sub-Coxeter chain expected from the framework. The signed channel residual difference $\delta = \rho_\mathrm{Schoen} - \rho_\mathrm{TY}$ is interpreted as a per-channel $\ln \mathrm{BF}_{\rm TY/Schoen} = \delta/2$ under a Laplace/BIC heuristic. The combined channel report is $\ln \mathrm{BF}_{\rm TY/Schoen} = -14.76$ (Tab.~\ref{tab:cy3-bayesian-channels}), corresponding to the highest tier of the Kass--Raftery scale~\cite{KassRaftery1995} (``very strong'' evidence at $|\ln \mathrm{BF}|>5$; equivalently $\log_{10} \mathrm{BF} > 6$, well beyond Jeffreys's ``decisive'' threshold of $\log_{10} \mathrm{BF} > 2$~\cite{Jeffreys1961}), with posterior odds Schoen:TY $= 2.6 \times 10^6$. The Wilks/Cowan--Cranmer--Gross--Vitells $\chi^2_1$ equivalence $n_\sigma = \sqrt{2|\ln \mathrm{BF}|}$ yields $5.43\sigma$-equivalent discrimination as an order-of-magnitude non-nested heuristic.

\paragraph{Convergence status of the chain-channel sweep.}
The chain-channel discriminators reported here are \emph{indicative,
not converged} at the strict 10\% basis-size convergence criterion.
The current production sweep at $n_{\rm pts}=25{,}000$, $k=3$,
$k_\mathrm{test} \in \{2,3,4,5\}$ records the channel JSONs at
\texttt{p7\_11\_quark\_chain.json} and
\texttt{p7\_11\_lepton\_chain.json}, both of which carry
\texttt{converged: false} flags with consecutive-degree relative
deltas of order 11--35\% on the Schoen side and the signed-delta
$\delta$ residual. The $\ln \mathrm{BF} = -14.76$ figure is therefore
a preliminary discriminator, not a converged Bayes factor in the
strict statistical sense; we report it as supporting evidence
alongside the catalogue uniqueness sweep (which \emph{is} closed
modulo the gCICY non-manifest residue) and the Tian--Yau Yukawa
structural exclusion (which \emph{is} closed at the searched
bidegree range). The full replacement of the metric-Laplacian
chain-residual likelihood by a converged PDG-fermion-mass
Yukawa-cohomology pipeline (Anderson--Karigiannis--Larfors--Lukas--Schneider
style: HYM Hermitian metric on $V$, harmonic representatives of
$H^1(X/\Gamma, V)$, Yukawa cohomology pairing, Kähler normalisation,
two-loop SM RG running) is the in-flight remediation specified in
the project's \texttt{REMEDIATION\_PLAN\_OPTION\_B.md}, Phase 2; we
expect to retract or upgrade the present chain-channel BF figure on
its completion. The catalogue uniqueness and TY-exclusion arguments
do not depend on the chain-channel convergence and stand on their
own.

\paragraph{$\sigma$-channel as discriminability, not BF input.}
The L$^2$ Donaldson §3.4 weighted-variance Monge--Ampère residual $\sigma$ on the Donaldson-balanced metric (Donaldson 2009~\cite{Donaldson2009}, PAMQ \textbf{5}, p.~587, eq.~3.4: $\sigma^2 = \langle (\eta/\kappa - 1)^2 \rangle_{w}$ with $\kappa = \langle \eta \rangle_{w}$, matching the production implementation in \texttt{schoen\_metric.rs::compute\_sigma} and \texttt{ty\_metric.rs::compute\_sigma}) is a separately-computed statistical discriminator that distinguishes the two candidates' metrics with $|t| = 6.92$ at $k=3$, BCa $[5.30, 9.04]$ (Tab.~\ref{tab:cy3-sigma-kscan}). The $k$-scan across $k=3,4,5$ reveals the empirical signature of basis-size dominance: as Schoen's basis dimension grows monotonically ($27 \to 48 \to 75$), the strict-converged Schoen subset shrinks ($16/20 \to 7/20 \to 6/20$) and the heavy-tailed $\sigma_{\rm Schoen}$ distribution increasingly drives the $|t|$-statistic. To test whether this is a real geometric signal or a basis-size artifact, we performed a paired-bootstrap test at matched basis size $n_b = n_b^{\rm Schoen} = 27$ between the two candidates' metrics, $20$ seeds, $B = 10^4$ resamples (Tab.~\ref{tab:cy3-sigma-artifact}). The signed gap \emph{inverts} when bases are matched, with disjoint BCa CIs and bootstrap probability of opposite sign equal to unity, and is the empirical basis for excluding $\sigma$ from the multi-channel Bayes factor (\texttt{for\_combination = false} flag on the \texttt{ChannelEvidence} record). The matched-basis comparison is currently implemented as a lex-sorted truncation of the larger TY basis to Schoen's 27-monomial dimension; a geometric basis projection (eigenvalue-ranked principal-components of the Laplacian, or harmonic-form ranking) is the in-flight upgrade to the matched-basis defense (\texttt{REMEDIATION\_PLAN\_OPTION\_B.md} Phase 3.3). The $\sigma$-discriminability is reported \emph{only} as a separate $|t|$-statistic, not as a likelihood contribution; the Bayes-factor exclusion of $\sigma$ is conditional on the matched-basis defense surviving the geometric-projection upgrade.

% CY3 substrate identification: Schoen vs Tian-Yau structural comparison.
% Source: book/scripts/cy3_substrate_discrimination/rust_solver/references/cy3_publication_summary.md §1
\begin{table}[h]
\centering
\caption{Structural comparison of the substrate's predicted Calabi--Yau threefold (Schoen $\mathbb{Z}_3 \times \mathbb{Z}_3$, Donagi--Ovrut--Pantev--Reinbacher~\cite{BraunHeOvrutPantev2005}) against the closest distinct three-generation competitor (Tian--Yau $\mathbb{Z}_3$). The framework's substrate-physical forcing chain (self-conjugacy on the internal Hodge decomposition, two-stage Wilson breaking $E_8 \to G_{\rm SM}$, Atiyah--Singer three-generation index) selects the Hodge $(3,3)$ / $\pi_1 = \mathbb{Z}_3 \times \mathbb{Z}_3$ branch on which Schoen sits and Tian--Yau does not.}
\label{tab:cy3-substrate-identification}
\begin{ruledtabular}
\begin{tabular}{lccc}
Quantity & Schoen $\mathbb{Z}_3 \times \mathbb{Z}_3$ & Tian--Yau $\mathbb{Z}_3$ & Selecting hypothesis \\
\hline
$(h^{1,1}, h^{2,1})$ & $(3,3)$ & $(1,4)$ & self-conjugacy + 3 generations \\
$\pi_1$ & $\mathbb{Z}_3 \times \mathbb{Z}_3$ & $\mathbb{Z}_3$ & two-stage Wilson breaking \\
$\chi(X)$ & $0$ & $-6$ & deformation-class distinction \\
$2\,\omega_{\rm fix}$ & $246/248$ & $245/248$ & gateway-formula ratio \\
\end{tabular}
\end{ruledtabular}
\end{table}

% Multi-channel Bayes-factor discrimination Schoen-vs-TY (load-bearing publication number).
% Source: book/scripts/cy3_substrate_discrimination/rust_solver/output/p8_1_bayes_multichannel.json
\begin{table}[h]
\centering
\caption{Multi-channel Bayes factor for the substrate-CY3 hypothesis test, $k=3$, $n_{\rm pts}=25{,}000$ chain / $40{,}000$ $\sigma$, $20$ seeds. Sign convention: $\ln \mathrm{BF} > 0 \Leftrightarrow$ Tian--Yau favoured, $< 0 \Leftrightarrow$ Schoen favoured. The combined chain-channel result clears decisive Bayesian strength on the Kass--Raftery scale ($|\ln \mathrm{BF}| > 5$). The $\sigma$-channel is reported separately as discriminability and is excluded from the model-comparison Bayes factor (rationale and empirical defence: Tab.~\ref{tab:cy3-sigma-artifact}).}
\label{tab:cy3-bayesian-channels}
\begin{ruledtabular}
\begin{tabular}{lccc}
Channel & Likelihood & $\ln \mathrm{BF}_{\rm TY/Schoen}$ & $n_\sigma$-equivalent \\
\hline
$\mathrm{chain}_{\rm quark}$ ($\{m_b, m_s, m_d\}$) & log-$\chi^2$ to PDG & $-5.60$ & --- \\
$\mathrm{chain}_{\rm lepton}$ ($\{m_\tau, m_\mu, m_e\}$) & log-$\chi^2$ to PDG & $-9.16$ & --- \\
\hline
\textbf{Combined (physics preference)} & --- & $\mathbf{-14.76}$ & $\mathbf{5.43\sigma}$ \\
\textbf{Strength (Kass--Raftery)} & --- & \multicolumn{2}{c}{\textbf{Strong (Schoen-favoured)}} \\
Posterior odds Schoen:TY & --- & \multicolumn{2}{c}{$2.6 \times 10^6 : 1$} \\
\hline
$\sigma$-channel (discriminability only; \emph{excluded from BF}) & Welch $t$ on $\sigma$-distribution & --- & $6.92\sigma$ \\
\end{tabular}
\end{ruledtabular}
\end{table}

% sigma-channel discriminability k-scan (supporting evidence; sigma not in BF).
% Sources: p5_10_ty_schoen_5sigma.json, p5_10_k4_damped.json, p5_10_k5_damped_relaunch.json
% Tier numbers: output/p5_10_n40k_p5_5k.log:75-78
\begin{table*}[t]
\centering
\caption{$\sigma$-channel discriminability $k$-scan on TY-vs-Schoen, Donaldson balanced metric (no Adam refinement), $n_{\rm pts}=40{,}000$, $20$ seeds, BCa-95\% CI from $B=10^4$ bootstrap resamples (DiCiccio--Efron 1996). Two tiers reported per row: \emph{conservative} (full $20{+}20$, no filter — referee-safe unconditional reading) and \emph{strict-converged} ($\mathrm{residual}<10^{-6}$ AND $\mathrm{iters}<\mathrm{cap}$, no Tukey trim, no guard-snapshot inflation — conditional on the convergence filter, which preferentially trims Schoen's heavier upper tail). At $k=3$ the strict filter drops $4/20$ Schoen seeds (including seed~5 with $\sigma=43.24$, $\sim 7\times$ the surviving cluster mean). The conservative-tier CI lower bound at $k=3$ straddles $5\sigma$ at $3.29$. The $k=4,5$ rows are reported as \emph{supporting tiers at higher basis order}, not as parallel publication results. The monotone shrinkage of the strict-converged Schoen subset ($16/20 \to 7/20 \to 6/20$) as Schoen's basis dimension grows ($27 \to 48 \to 75$) is the empirical signature of $\sigma$ as a basis-size-saturated indicator at finite $k$, motivating exclusion from the model-comparison Bayes factor (Tab.~\ref{tab:cy3-sigma-artifact}). Tier numbers from \texttt{output/p5\_10\_n40k\_p5\_5k.log:75-78}.}
\label{tab:cy3-sigma-kscan}
\begin{ruledtabular}
\begin{tabular}{cccccc}
$k$ & tier & $\langle\sigma_{\rm Schoen}\rangle$ & $|t|$-stat & BCa 95\% CI & $n_{\rm Schoen}^{\rm strict}/n_{\rm basis}^{\rm Schoen}$ \\
\hline
$\mathbf{3}$ (publication) & conservative (full $20{+}20$) & $8.15$ & $\mathbf{3.99\sigma}$ & $[2.89,\ 8.97]$ & $20/20\,/\,27$ \\
$\mathbf{3}$ (publication) & strict-converged (conditional) & $5.59$ & $\mathbf{6.92\sigma}$ & $\mathbf{[5.30,\ 9.04]}$ & $16/20\,/\,27$ \\
$4$ (damped, supporting)   & strict-converged & $13.21$ & $3.82\sigma$ & $[2.48,\ 6.37]$ & $7/20\,/\,48$ \\
$5$ (supporting)           & strict-converged & $13.05$ & $10.87\sigma$ & $[8.79,\ 15.01]$ & $6/20\,/\,75$ \\
\end{tabular}
\end{ruledtabular}
\end{table*}

% sigma-channel basis-size sign-reversal: empirical defence for sigma exclusion from BF.
% Source: book/scripts/cy3_substrate_discrimination/rust_solver/output/p_basis_convergence_prod.json
\begin{table}[h]
\centering
\caption{Empirical defence of the $\sigma$-channel exclusion from the multi-channel Bayes factor: matched-basis-size paired bootstrap test on TY-vs-Schoen, $k=3$, $n_{\rm pts}=40{,}000$, donaldson\_iters$=100$, donaldson\_tol$=10^{-6}$, $20$ seeds, $B=10^4$ resamples, boot\_seed $\in \{12345, 999, 31415\}$ for jitter check. At native basis sizes $\sigma_{\rm TY}<\sigma_{\rm Schoen}$ as expected. \textbf{When the two metrics' bases are truncated to a common dimension $n_b=27$ (Schoen's native value), the sign of the $\sigma$-gap inverts}, with disjoint BCa CIs and paired-bootstrap probability of opposite sign equal to unity over $10^4$ resamples. The artifact magnitude is $\sim 3\times$ the native gap, in the opposite direction. This is the textbook signature of a basis-size-induced artifact, not a likelihood contribution, and is the empirical basis for excluding $\sigma$ from the model-comparison Bayes factor of Tab.~\ref{tab:cy3-bayesian-channels}.}
\label{tab:cy3-sigma-artifact}
\begin{ruledtabular}
\begin{tabular}{lcc}
Quantity & Mean & BCa 95\% CI \\
\hline
$\Delta\sigma_{\rm native}$ ($n_{\rm TY}=87,\ n_{\rm Schoen}=27$) & $-7.89$ & $[-15.25,\ -5.52]$ \\
$\Delta\sigma_{\rm matched}$ ($n_{\rm TY}=n_{\rm Schoen}=27$)  & $\mathbf{+23.59}$ & $\mathbf{[+11.33,\ +42.66]}$ \\
$|\Delta_{\rm matched}|/|\Delta_{\rm native}|$ & $2.99$ & $[1.54,\ 6.55]$ \\
\hline
$\mathbb{P}(\mathrm{opposite\ sign}\mid \mathrm{boot})$ & \multicolumn{2}{c}{$\mathbf{1.0000}$ ($10^4$ resamples)} \\
Bootstrap-seed jitter on BCa-low & \multicolumn{2}{c}{$0.52$ (range over $\{12345, 999, 31415\}$)} \\
\end{tabular}
\end{ruledtabular}
\end{table}


\paragraph{Tian--Yau Yukawa-channel structural exclusion.}
Beyond the chain-channel Bayesian preference, we attempted to construct a working three-generation Standard-Model bundle on TY/$\mathbb{Z}_3$ in the bicubic CICY ambient $\mathbb{P}^3 \times \mathbb{P}^3$ at $h^{1,1} = 2$, mirroring the working three-factor monad construction we use on Schoen (Tab.~\ref{tab:cy3-schoen-bundle}). A nine-cycle hypothesis-test enumeration over rank-3 SU(3) line-bundle-sum and monad bundles, summarised in Tab.~\ref{tab:cy3-ty-exclusion}, found no candidate that simultaneously satisfies $c_1(V) = 0$, $\Sigma\,c_3(V) = \pm 18$, $\mathrm{ch}_2(V) \le \mathrm{ch}_2(TX)$, slope-stability via line-destabiliser test, balanced $1{:}1{:}1$ Wilson $\mathbb{Z}_3$ phase partition, basis-rank $\ge 6$ for three-generation Yukawa population, and Mumford--Takemoto SU(3) stability~\cite{HuybrechtsLehn2010}. The line-bundle-sum result reproduces and extends Anderson--Gray--Lukas--Palti~\cite{AndersonGrayLukasPalti2012} §5.3's catalogue exclusion of $h^{1,1} \in \{2,3\}$ to broader bundle constructions and to the explicit Mumford--Takemoto stability check. The leading rank-7 monad survivor of the topological + summand-stability constraints, $V_{\min 2}$, is mathematically falsified by $h^0(V_{\min 2}) = \mathbb{C}^2 \neq 0$ in the long exact sequence of $V = \ker(B \to C)$, violating the $h^0(V) = 0$ requirement for stable rank-3 SU(3) bundles on a Calabi--Yau threefold (cf.~\cite{HuybrechtsLehn2010} §1.2). This structural exclusion is independent of the Bayesian-discrimination evidence and converges on the same conclusion via different mathematical machinery.

% TY/Z3 Yukawa-bundle structural exclusion: nine-cycle hypothesis-test enumeration.
% Source: book/scripts/cy3_substrate_discrimination/rust_solver/references/p_ty_bundle_research_log.md
\begin{table*}[t]
\centering
\caption{Structural exclusion of a working three-generation Standard-Model bundle on the closest competitor manifold Tian--Yau $\mathbb{Z}_3$ (bicubic CICY in $\mathbb{P}^3 \times \mathbb{P}^3$, $h^{1,1}=2$): nine-cycle empirical enumeration. Constraints applied per cycle, in chain: $c_1(V)=0$ \textsc{and} $\Sigma\,c_3(V)=\pm 18$ for three generations downstairs after the $\mathbb{Z}_3$ quotient \textsc{and} $\mathrm{ch}_2(V) \le \mathrm{ch}_2(TX)$ for anomaly cancellation \textsc{and} slope-stability via line-destabiliser test on the Kähler cone \textsc{and} balanced $1{:}1{:}1$ Wilson $\mathbb{Z}_3$ phase partition on $V$ \textsc{and} basis-rank $\ge 6$ for three-generation Yukawa population \textsc{and} Mumford--Takemoto SU(3) stability (Huybrechts--Lehn 2010 §1.2: $h^0(V) = 0$). Cycles run on monad bundles $V = \ker(B \to C)$ with $B, C$ line-bundle sums, bidegrees in $[-2,2]^2$. The cycle 1 line-bundle-sum exclusion empirically reproduces and extends Anderson--Gray--Lukas--Palti (arXiv:1202.1757) §5.3's published result that no phenomenologically viable model exists at $h^{1,1}\in\{2,3\}$ in their scan. No candidate survives all nine constraints simultaneously.}
\label{tab:cy3-ty-exclusion}
\begin{ruledtabular}
\begin{tabular}{cllr}
Cycle & Hypothesis (rank-3 SU(3) bundle on TY/$\mathbb{Z}_3$ bicubic) & Failure mode & Survivors \\
\hline
H1 & line-bundle sums, Wilson $1{:}1{:}1$, $c_3=\pm 18$        & polystability fails universally        & $0/{\sim}5.97\times 10^6$ \\
H2 v1/v2 & monad $\ker(B\to C)$, weak polystability filter         & polystability under correct slope     & $0/-$ \\
H3 & SU(4)/SO(10) rank-4 alternative embedding                & same polystability obstruction         & $0/{\sim}9 \times 10^4$ \\
H2 v3 & monad, closed-negative-cone polystability               & (passes topological + summand-stable)  & $10/{\sim}3.78\times 10^6$ \\
C4 & $V_{\min}$ wired ($\mathrm{rk}\,B = 4$)                  & basis-rank deficit; $3$ modes $\neq 9$ & $0$ (operationally) \\
C5 & rank-$B \ge 6$ rescan                                   & (10 v3 candidates re-confirmed)        & $10$ \\
C6 & $V_{\min 2}$ wired ($\mathrm{rk}\,B = 7$)                & $H^0$-seed empty on $\mathrm{neg}$. summands & $0$ (operationally) \\
C7 & $H^0/H^1$ cohomology audit                              & (LES-collapse argument validates AKLP/Schoen) & --- \\
\textbf{C8} & $V_{\min 2}$ mathematical falsification                  & $h^0(V) = \mathbb{C}^2 \neq 0$: violates Mumford--Takemoto & $\mathbf{0}$ \\
\textbf{C9} & $H^0(V) = 0$ filter on cycle-5 v3 set                  & operationally infertile: $\Sigma h^0(B) = \Sigma h^1(B) = 0$ & $\mathbf{0/10}$ \\
\hline
\multicolumn{4}{l}{\textbf{Net empirical result:}\ no rank-3 SU(3) Yukawa-active SM monad bundle in scanned bidegree range.} \\
\end{tabular}
\end{ruledtabular}
\end{table*}

% Schoen Z3xZ3 working bundle specification used in the chain channels.
% Source: src/zero_modes.rs::MonadBundle::schoen_z3xz3_canonical (P8.3-followup-B)
\begin{table}[h]
\centering
\caption{Working three-generation Standard-Model bundle on Schoen $\mathbb{Z}_3 \times \mathbb{Z}_3$ used by the chain channels of Tab.~\ref{tab:cy3-bayesian-channels}. Three-factor monad construction $V = \ker(B \to C)$ on the Schoen ambient $\mathbb{P}^2 \times \mathbb{P}^2 \times \mathbb{P}^1$. The Wilson $\mathbb{Z}_3 \times \mathbb{Z}_3$ phase partition splits the $V$-summands evenly into three classes (one per generation) under the projection $(a-b) \bmod 3$ on the bidegrees, breaking the previously degenerate sector assignment that affects Tian--Yau under the same construction (Tab.~\ref{tab:cy3-ty-exclusion}). Anomaly cancellation $\mathrm{ch}_2(V) \le \mathrm{ch}_2(TX)$ is satisfied with the indicated 5-brane class.}
\label{tab:cy3-schoen-bundle}
\begin{ruledtabular}
\begin{tabular}{lll}
Object & Specification & Property \\
\hline
$B$ & $\mathcal{O}(1,0,0)^2 \oplus \mathcal{O}(0,1,0)^2 \oplus \mathcal{O}(0,0,1)^2$ & rank $6$ \\
$C$ & $\mathcal{O}(1,1,0) \oplus \mathcal{O}(0,1,1) \oplus \mathcal{O}(1,0,1)$  & rank $3$ \\
$V \;=\; \ker(B \to C)$ & rank $3$ SU(3) & $c_1(V) = 0$, $\Sigma\,c_3(V) = \pm 18$ \\
Wilson Z$_3$ partition on $V$ & $(1,1,1)$ across phase classes & balanced; canonical sector assignment \\
Anomaly cancellation & $\mathrm{ch}_2(V) \le \mathrm{ch}_2(TX)$ with $W=(27,27)$ & satisfied \\
Yukawa coverage on $V$ & $Y_u : 9/9$, $Y_d : 9/9$ on $h_0$ slice & operationally productive \\
\end{tabular}
\end{ruledtabular}
\end{table}

% CY3 discrimination methodology and conventions (replicability table).
% Source: book/scripts/cy3_substrate_discrimination/rust_solver/references/cy3_publication_summary.md §6
\begin{table}[h]
\centering
\caption{Methodological provenance and replicability conventions for the CY3-substrate-identification analysis (Tabs.~\ref{tab:cy3-substrate-identification}--\ref{tab:cy3-schoen-bundle}). Run-time per sweep: $\sim 2.5$~h CPU on the basis-convergence sign-reversal test; $\sim 3.5$~h GPU on each $k$-scan production run; total ${\sim}50$ GPU-hours plus ${\sim}30$ CPU-hours across all reported tiers. The $\sigma$-discriminator throughout this paper is the weighted-L$^2$ Donaldson 2009 §3.4 functional, matching the production implementation; see row below.}
\label{tab:cy3-methodology}
\begin{ruledtabular}
\begin{tabular}{ll}
Item & Setting / source \\
\hline
$\sigma$-functional definition & weighted L$^2$ Monge--Ampère residual: $\sigma^2 = \langle (\eta/\kappa - 1)^2 \rangle_{w}$, $\kappa = \langle \eta \rangle_{w}$ (Donaldson 2009 §3.4 eq.~3.4) \\
$\sigma$-functional implementation & \texttt{schoen\_metric.rs::compute\_sigma}, \texttt{ty\_metric.rs::compute\_sigma} (CPU + GPU bit-exact mirror) \\
Bayes-factor sign convention & $\ln \mathrm{BF}_{12} > 0 \Leftrightarrow M_1$ favoured (Jeffreys/Kass--Raftery) \\
Strict-converged tier gate & $\mathrm{residual} < \mathrm{tol}$ AND $\mathrm{iters} < \mathrm{cap}$, no Tukey trim \\
Bootstrap CI methodology & percentile and bias-corrected accelerated, $B=10^4$ resamples \\
Bootstrap-seed jitter check & boot\_seed $\in \{12345, 999, 31415\}$, target $<0.05\sigma$ on BCa-low \\
Donaldson update (damped, $k\ge 4$) & $h \leftarrow \alpha\,T(G)^{-1} + (1-\alpha)\,h_{\rm old}$, $\alpha = 0.5$ \\
GPU/CPU bit-exactness ($h$-pair) & mirrored 256-lane tree reduction, iter-0 $\Delta < 4.4\times 10^{-16}$ \\
Reproducibility manifest & SHA-256 chained event log per run \\
$n_\sigma$-equivalent from $\ln \mathrm{BF}$ & $n_\sigma = \sqrt{2|\ln \mathrm{BF}|}$ (Wilks/Cowan--Cranmer--Gross--Vitells) \\
Source code & \texttt{book/scripts/cy3\_substrate\_discrimination/rust\_solver/} \\
Output JSON catalogue & \texttt{output/p\{5\_10,7\_11,8\_1,\_basis\_convergence\}\_*.json} \\
Research log (9-cycle TY exclusion) & \texttt{references/p\_ty\_bundle\_research\_log.md} \\
Consolidated audit document & \texttt{references/cy3\_publication\_summary.md} \\
\end{tabular}
\end{ruledtabular}
\end{table}


\paragraph{Catalogue uniqueness at $h^{1,1} = h^{2,1} = 3$.}
The chain-channel Bayes factor and the Tian--Yau Yukawa exclusion identify Schoen as the favoured candidate within the framework's pre-narrowed two-element shortlist. To upgrade this from a pairwise discrimination to a catalogue-level uniqueness statement we performed a four-cycle sweep of the published smooth-CY3 universe at $(h^{1,1}, h^{2,1}) = (3,3)$ for non-Schoen entries admitting a free $\mathbb{Z}_3 \times \mathbb{Z}_3$ action with Standard-Model-compatible Wilson-line breaking. Cycle~1 (Hodge filter) returned zero hard-match competitors across CICY (Anderson--He--Lukas~\cite{AndersonHeLukas2008}), Schoen-class fiber products, Constantin--Lukas free-quotient catalogue, Borisov--Caldararu Pfaffian threefolds, and gCICY (Anderson--Apruzzi--Gao--Gray--Lee~\cite{AndersonApruzziGaoGrayLee2015}). Cycle~2 (free-action structural filter) returned zero non-Schoen survivors. Cycle~2.5 (computer-algebra sweep) verified by direct query against the Kreuzer--Skarke 473{,}800{,}776-polytope catalogue~\cite{KreuzerSkarke2002} that the $(3,3)$ Hodge bin is \emph{literally empty} (response \texttt{\#NF: 0} at output limits $L \in \{100, 1000, 10000\}$, cross-verified by sweeping the entire $\chi = 0$ self-mirror diagonal $h = 1, \ldots, 14$ and confirming the bin is empty for $h \le 13$ and populated from $h = 14$). Cycle~2.6 (gCICY non-manifest Aut residue) reduced the last open cell by programmatic literature search of arXiv:1507.03235 for the literal $(3,3)$ Hodge pair (zero hits), cross-checked against Larfors--Schneider~\cite{LarforsLukas2020} and Constantin--Lukas--Mishra~\cite{ConstantinLukasManuwal2016}; the strict-CAS portion of cycle 2.6 (computing the full ambient automorphism group via Sage / Macaulay2 / polymake and checking which $\mathbb{Z}_3 \times \mathbb{Z}_3$ subgroups act freely on a generic gCICY $(3,3)$) remains \textbf{DEFERRED} pending a CAS-equipped re-run; see \texttt{p\_schoen\_uniqueness\_cycle2\_5.md} §3.4 and the in-flight Phase 6.2 of the remediation plan. The combined catalogue scan, modulo the gCICY non-manifest CAS residue, supports the statement: \emph{no published smooth Calabi--Yau threefold in the catalogued universe is established to satisfy the substrate-physical four-principle constraint at $(h^{1,1}, h^{2,1}) = (3,3)$ except Schoen $\mathbb{Z}_3 \times \mathbb{Z}_3$.} The remaining residual is a future gCICY codim~$\ge (3,1)$ extension producing a $(3,3)$ entry, plus the deferred CAS check on existing gCICY $(3,3)$ classifications; both would require a re-test but do not affect the present empirical identification.

\paragraph{Net empirical statement.}
Three convergent lines of evidence support the empirical identification of the substrate's internal Calabi--Yau threefold as Schoen $\mathbb{Z}_3 \times \mathbb{Z}_3$ over the closest structural competitor Tian--Yau $\mathbb{Z}_3$. \textbf{Closed:} (a) the four-cycle catalogue uniqueness sweep finds no non-Schoen $(h^{1,1}, h^{2,1}) = (3,3)$ smooth CY3 with a free $\mathbb{Z}_3 \times \mathbb{Z}_3$ action across the published universe (CICY, Kreuzer--Skarke toric, Schoen-class, Constantin--Lukas, Borisov--Caldararu Pfaffian, gCICY); the Kreuzer--Skarke direct query records \texttt{\#NF: 0} at output limits up to $L = 10{,}000$, cross-verified by sweeping the entire $\chi=0$ self-mirror diagonal $h \in \{1, \ldots, 14\}$. (b) the Tian--Yau Yukawa channel is structurally excluded by a nine-cycle hypothesis-test enumeration showing that no rank-3 SU(3) stable polystable monad bundle on TY/$\mathbb{Z}_3$ in the bicubic ambient at $h^{1,1} = 2$ in the searched bidegree range simultaneously satisfies all required physical constraints. \textbf{Indicative, not converged:} (c) the chain-channel Bayes factor at $\ln \mathrm{BF} = -14.76$ (posterior odds $2.6 \times 10^6$:1 Schoen:TY; nested-LRT $\sigma$-equivalent $5.43\sigma$) is computed from a metric-Laplacian Galerkin-residual sweep that has not yet met the strict 10\% basis-size convergence criterion at $k_\mathrm{test} \in \{2,3,4,5\}$; we report it as supporting evidence and expect to retract or upgrade the figure on completion of the full PDG-fermion-mass Yukawa-cohomology pipeline (Phase 2 of the remediation plan). The Donaldson $\sigma$-channel is excluded from the model-comparison Bayes factor by a paired-bootstrap matched-basis sign-reversal test (currently using lex truncation; geometric-projection upgrade is in flight at Phase 3). \emph{Conditional on the open methodological items below and the in-flight pipeline upgrades enumerated in Sec.~\ref{sec:open-items}, the substrate-physical four-principle constraint narrows the candidate manifold to a single element of the catalogued universe at the searched Hodge bin.}

\paragraph{What this section does \emph{not} claim.}
We are explicit about the residual scope. (i)~We do \emph{not} claim Schoen-uniqueness in full generality across the Calabi--Yau threefold landscape; the catalogue uniqueness statement is at the single Hodge bin $(h^{1,1}, h^{2,1}) = (3,3)$ and is conditional on the gCICY non-manifest CAS residue (item~iii below). (ii)~The chain-channel Bayes factor $\ln \mathrm{BF} = -14.76$ is \emph{not} a converged result at the strict 10\% basis-size convergence criterion: both the quark-chain and lepton-chain JSON outputs (\texttt{p7\_11\_quark\_chain.json}, \texttt{p7\_11\_lepton\_chain.json}) carry \texttt{converged: false}. The figure is reported as indicative; the full PDG-fermion-mass Yukawa-cohomology pipeline replacement is in-flight Phase 2 of the remediation plan. (iii)~The gCICY $(3,3)$ uniqueness check has \textbf{deferred} the strict computer-algebra-system portion of cycle 2.6 (Sage / Macaulay2 / polymake-based ambient automorphism enumeration) pending CAS environment availability; the literature-cross-check portion is closed. A future gCICY codim~$\ge (3,1)$ extension producing a $(3,3)$ entry would also require a re-test. (iv)~We do \emph{not} claim that the Bayesian discrimination is independent of the Donaldson convention adopted for the Schoen and Tian--Yau metrics; the $\sigma$-channel result is sensitive to h-inversion handling and Kähler-form projection, and we report only the conventions used here. The $\sigma$-channel matched-basis defense currently uses lex truncation rather than a geometric basis projection; the geometric-projection upgrade and the full Donaldson convention robustness sweep are in-flight Phase 3 of the remediation plan. (v)~The TY Yukawa exclusion is at the searched bidegree range in the bicubic ambient and is consistent with, but does not rigorously imply, the absence of any TY-supported three-generation bundle outside that range. (vi)~The bundle data we use on Schoen, and the line-bundle enumerations cited from the literature, are subject to the bundle-provenance audit summarised in Sec.~\ref{sec:open-items}. None of items (i)--(vi) affects the numerical computations of Secs.~\ref{sec:gateway}--\ref{sec:gravity}, which proceed from the Schoen identification taken as a starting commitment.

\subsection{The $E_8$ Coxeter spectrum}

The $E_8$ Lie algebra has dimension $\dim(E_8) = 248$, rank $8$,
$240$ roots, Coxeter number $h(E_8) = 30$, and Weyl group order
$|W(E_8)| = 696{,}729{,}600$. The exponents of $E_8$ are the eight
integers
\begin{equation}
\{m_1, \ldots, m_8\} = \{1, 7, 11, 13, 17, 19, 23, 29\},
\end{equation}
palindromically symmetric around $h/2 = 15$ in pairs $(1, 29)$,
$(7, 23)$, $(11, 19)$, $(13, 17)$, satisfying $m_j + m_{9-j} = h$.
The longest Weyl element acts as $w_0 = -1$ on the Cartan
subalgebra of $E_8$ (a property shared by $A_1, B_n, C_n, D_{2n},
E_7, E_8, F_4, G_2$). The harmonic threshold $h/2 = 15$ is the
framework's chirality-flip locus for the Coxeter chain and is not
itself an exponent.

\paragraph{Framework extension beyond the standard Coxeter spectrum.}
The mass-formula construction in this paper parameterizes
\emph{continuous} chain positions $k \in \mathbb{R}_{\geq 0}$, not
just the eight integer exponents above. Half-integer positions
($k_d = 6.5$, $k_W = 34.5$, $k_t = 36.5$, etc.) and
``cycle-2'' positions ($k > h$) are framework-specific extensions
that interpolate between integer exponents and re-use the
periodicity of the Coxeter element across multiple cycles. These are
structural assignments within the framework rather than standard
Coxeter-theoretic invariants. The structural ``boundary capacity''
$h + m_{\max} = 59$ used in the dark-sector coupling family
(Sec.~\ref{sec:higgs}) is likewise a framework combination of
$E_8$ invariants rather than an established Lie-algebra invariant
in its own right.

\subsection{The McKay correspondence}

The McKay correspondence~\cite{McKay1980} establishes a bijection
between finite subgroups of $SU(2)$ and the affine simply-laced ADE
Dynkin diagrams: cyclic $\mathbb{Z}_n \leftrightarrow \widehat{A}_{n-1}$,
binary dihedral $\leftrightarrow \widehat{D}_n$, binary tetrahedral
$\leftrightarrow \widehat{E}_6$, binary octahedral
$\leftrightarrow \widehat{E}_7$, and \textbf{binary icosahedral
$\leftrightarrow \widehat{E}_8$}. The McKay correspondence itself
identifies vertices of the affine $\widehat{E}_8$ Dynkin diagram
with irreducible representations of the binary icosahedral group;
it does \emph{not}, on its own, single out a particular integer.
Separately, the icosahedral rotation group $A_5$ contains elements
of order $5$ in the $(2,3,5)$-triangle group, with
$\varphi = 2\cos(\pi/5) = (1+\sqrt{5})/2$ as the icosahedral angular
ratio. The framework adopts this fivefold rotation order, paired
with the binary icosahedral McKay-dual of $E_8$, as the seed integer
for the $\varphi$-chain. The use of the integer $5$ is therefore a
framework structural commitment combining the McKay-dual
identification with the icosahedral rotation order, not a forced
output of the McKay correspondence.

\paragraph{Two harmonic chains used by the framework.}
The mass-formula construction below distinguishes a $\varphi$-chain
(with $\varphi^k$ primary harmonics, motivated by 5-fold icosahedral
recursion at the binary-icosahedral McKay-dual of $E_8$) from a
$\sqrt{2}$-chain (with $(\sqrt{2})^k$ primary harmonics, motivated
by 4-fold octahedral recursion at the binary-octahedral McKay-dual
of $E_7$, where $\sqrt{2} = 2\cos(\pi/4)$). Charged leptons and
neutrinos populate the $\varphi$-chain; quarks and electroweak
bosons populate the $\sqrt{2}$-chain. This sector split is a
\emph{framework structural assignment} rather than a derivation
from $E_8$ alone --- a single $E_8$ algebra does not naturally
decompose into icosahedral and octahedral sub-recursions, and the
allocation of leptons to $\varphi$ and quarks to $\sqrt{2}$ is
adopted on phenomenological grounds (matching observed masses with
small-integer Coxeter-step assignments). A first-principles
derivation of the chain assignment from a specific Calabi--Yau
compactification is left for future work.

\paragraph{Chain factors are not Coxeter eigenvalues.} We emphasise
that the harmonic factors $\varphi^k$ and $(\sqrt{2})^k$ that appear
throughout the mass formulas below are \emph{exponential
growth/recursion factors} on a discrete chain --- successive
applications of a fixed scalar step --- and are \emph{not} eigenvalues
of the Coxeter element of $E_8$. The latter are the unit-modulus
roots of unity $|e^{2\pi i e_j/h}| = 1$ for $e_j \in \{1, 7, 11, 13,
17, 19, 23, 29\}$, which describe rotational structure on the Coxeter
plane and have nothing directly to do with mass-ratio recursion. The
exponents $e_j$ enter the mass formulas separately, as integer
positions used to label palindromic pairings (e.g. $11 + 19 = h$);
the multiplicative factors $\varphi^k, (\sqrt{2})^k$ encode the
geometric step magnitude of the chain itself. These are two
mathematically distinct uses of $E_8$ data, and the framework's
construction is the joint use of both.

\subsection{Substrate-tension interpretation}

The framework adopts a substrate-tension interpretation of the Higgs
mechanism, motivated by Anderson's superconducting precursor to
gauge-boson mass generation~\cite{Anderson1962}; the resulting picture
is similar in spirit to Volovik's quantum-liquid / helium-droplet
emergent-vacuum program~\cite{Volovik2001,Volovik2003}, though arrived
at independently and from a different structural starting point. Within the substrate framework,
the Higgs vacuum expectation value $v$ is interpreted as the baseline
relaxed tension of the $E_8 \times E_8$ substrate, and massive SM
particles are localized deformations of this tension at
structurally-determined Coxeter positions. Photons are massless because
they correspond to solitonic excitations of the substrate that do not
deform its baseline tension.

\paragraph{Foundations deferred to companion paper.}
The substrate-physical reading of $c$, $\hbar$, the elastic-medium
identification $c = \sqrt{\Theta/\mathcal{I}}$, and the macroscopic
emergence of general relativity, quantum mechanics, Newton's laws of
motion, and Maxwell's equations from the same elastic substrate that
hosts the $E_8 \times E_8$ Coxeter structure are not load-bearing for
the numerical demonstration of this paper, and the present paper
takes them as part of the starting commitment of
Table~\ref{tab:starting-assumptions}. Their derivation chains are
developed in the companion paper~\cite{UnreasonableEffectiveness}
and the foundational monograph~\cite{ExistBook}. We do not repeat
them here.

\section{Gateway formula}
\label{sec:gateway}

The remainder of this paper is the direct computation of SM observables
from the starting assumptions of Table~\ref{tab:starting-assumptions}.
There is no per-particle scheme choice and no fitted continuous
parameter: every numerical value reported below is the framework's
structural prediction. PDG-matching residuals are predictions versus
measurement, not fits versus measurement.

The electron sits at the gateway $k=0$, defined as one chain step
before the first $E_8$ exponent (at position 1). The chain depth from
gateway to harmonic threshold is
\begin{equation}
M = h/2 - 1 = 14.
\end{equation}

The framework's effective coupling scale is
\begin{equation}
h_{\rm eff} \;=\; \frac{m_e}{2 \omega_{\rm fix}}
            \;=\; m_e \cdot \frac{\dim}{\dim - 2}
            \;=\; m_e \cdot \frac{124}{123},
\end{equation}
where $\omega_{\rm fix} = 1/2 - 1/\dim = 123/248$ is the gateway
normalization factor. Equivalently, with the Planck mass
$m_{\rm Planck} = \sqrt{\hbar c / G}$ taken from the structurally
predicted $G$ of Sec.~\ref{sec:gravity},
$h_{\rm eff} = F_e \cdot m_{\rm Planck} \cdot 124/123$, with $F_e$
fixed entirely by Eq.~\eqref{eq:gateway} below; the two readings
agree at the \SI{17}{ppm} residual reported here, and the downstream
mass formulas are operationally identical under either reading.

The dimensionless ratio of the electron mass to the Planck mass is
predicted by the framework as
\begin{equation}
\boxed{
F_e \;=\; \frac{m_e}{m_{\rm Planck}}
    \;=\; \frac{\dim - 2}{\dim} \cdot \frac{5}{\dim \cdot h^M}
      \cdot \frac{1}{1 - \frac{1}{4(\dim - 1)}}.
}
\label{eq:gateway}
\end{equation}
The right-hand side contains \emph{no measured input}: every quantity
is a structural integer of the starting assumptions
(Table~\ref{tab:starting-assumptions}). This is a zero-parameter
prediction of the dimensionless electron-to-Planck mass ratio.

Each factor has a specific structural role in the framework:
\begin{itemize}
\item Gateway factor $(\dim - 2)/\dim = 246/248$: framework
dual-sector count subtracting two distinguished endpoints (one in each
$E_8$).
\item McKay icosahedral $5$: framework-selected icosahedral order
motivated by the binary-icosahedral McKay dual of $E_8$.
\item Coxeter chain $h^M = 30^{14} \approx 4.78 \times 10^{20}$:
each step contributes a factor of $h = 30$.
\item Sub-harmonic loss $988/987 = 1/(1 - 1/988)$: the geometric
sum $\sum_{n \geq 0} (1/988)^n$ of a per-step leakage rate
$1/(4(\dim - 1)) = 1/988$ to all orders. The factor of $4$ in
$4(\dim - 1)$ admits a dual-sector / two-orientation reading
(``$2 \times 2$'': two $E_8$ factors, each carrying both an upward
and a downward chain-step orientation; $\dim - 1$ removes the
distinguished anchor mode in each sector); a substrate-physical
derivation of the $4$ from a heterotic-anomaly index or a
$\mathbb{Z}_3 \times \mathbb{Z}_3$ quotient trace on the Schoen
threefold is part of the open companion-paper program
(Sec.~\ref{sec:open-items}).
\end{itemize}

Numerically, $F_e = 4.18547 \times 10^{-23}$. The measured value
$m_e^{\rm PDG} / m_{\rm Planck}^{\rm CODATA}$ agrees with this
structural prediction at \SI{17}{ppm}. Newton's constant is the
zero-parameter consequence $G = F_e^2 \cdot \hbar c / m_e^2$ developed
in Sec.~\ref{sec:gravity}.

\section{Charged-fermion masses}
\label{sec:fermions}

\subsection{Mass formula schema}

Every charged-fermion mass in the framework takes the unified form
\begin{equation}
\boxed{m_X \;=\; h_{\rm eff} \cdot R_X(k_X, \text{sector}, \{\text{corr.}\}),}
\label{eq:massschema}
\end{equation}
where $h_{\rm eff} = m_e \cdot \dim/(\dim-2) = 124 m_e/123$ is the
gateway-normalized scale (Sec.~\ref{sec:gateway}), $k_X$ is
the particle's Coxeter step assignment in the $E_8$ spectrum (cf.\
App.~\ref{app:assignments}), and $R_X$ is a structural ratio built
from $E_8$ invariants $\{\dim, h, M, \text{rank}, m_{\max}\}$, the
gateway-mixing parameter $\alpha_\tau = 1/(\dim+h) = 1/278$, and
sector-specific sub-harmonic corrections.
Numerical values of $R_X$ for every particle are tabulated in
Sec.~\ref{sec:results} (Table~\ref{tab:complete-table}); the closed-form expressions are given below
sector by sector.

\subsection{Lepton sector}

The electron, muon, and tau occupy integer $E_8$ exponents on the
$\varphi$-chain (icosahedral bulk recursion) at $k_e = 0$ (gateway),
$k_\mu = 11$, and $k_\tau = 17$ respectively. Their explicit
closed-form ratios are
\begin{align}
R_e \;&=\; \frac{\dim - 2}{\dim} \;=\; \frac{246}{248},
   \label{eq:Re} \\
R_\mu \;&=\; \varphi^{11} \cdot \bigl(1 + \Delta_\varphi\bigr) \cdot \bigl(1 - \alpha_\mu\bigr),
   \label{eq:Rmu} \\
R_\tau \;&=\; \varphi^{17} \cdot \bigl(1 - \Delta_\varphi\bigr) \cdot \bigl(1 - \alpha_\tau^{\rm full}\bigr),
   \label{eq:Rtau}
\end{align}
where the icosahedral-bulk closure factor is
\begin{equation}
\Delta_\varphi \;=\; \frac{1}{h + \varphi^2},
\label{eq:Deltaphi-lepton}
\end{equation}
and the sub-harmonic corrections are
\begin{align}
\alpha_\mu \;&=\; \sum_{n=1}^{\infty} c_n^{(\mu)} \alpha_\tau^n
            \;=\; 2 \alpha_\tau^2 + \mathcal{O}(\alpha_\tau^3),
   \label{eq:alphamu} \\
\alpha_\tau^{\rm full} \;&=\; \frac{\alpha_\tau}{1 + \varphi^2 \alpha_\tau}
            \;=\; \frac{1}{\dim + h + \varphi^2}.
   \label{eq:alphataufull}
\end{align}
Each lepton ratio has three structural pieces: a $\varphi$-chain
primary harmonic $\varphi^{k_\ell}$ (icosahedral bulk recursion at
the lepton's Coxeter exponent), an icosahedral-bulk closure
$(1 \pm \Delta_\varphi)$ with sign opposite for muon and tau
(reflecting their opposite displacement from the harmonic threshold
$h/2 = 15$), and a sub-harmonic series correction. The muon
coefficient $c_2^{(\mu)} = 2$ is the framework's minimal assignment to
the palindromic pair $(11, 19)$ centered on $h/2$. The
tau correction $\alpha_\tau^{\rm full}$ is the closed-form sum of the
geometric series $\sum_n (-\varphi^2 \alpha_\tau)^n \alpha_\tau$ in
the icosahedral self-coupling.
Predicted (electron-anchor): $m_\mu = 105.6583756$~MeV (PDG match:
$+1.1$~ppb), $m_\tau = 1776.8597$~MeV ($-0.15$~ppm).

\subsection{Quark sector}

Quarks occupy positions on the $\sqrt{2}$-chain (octahedral bulk
recursion) at integer and half-integer Coxeter steps. Their explicit
closed-form ratios are
\begin{equation}
R_q \;=\; (\sqrt{2})^{k_q} \cdot \bigl(1 + \chi_q \Delta_2\bigr) \cdot \bigl(1 + \lambda_q \tfrac{|k_q - h/2|}{2\dim}\bigr) \cdot \bigl(1 - \alpha_q\bigr),
\label{eq:Rquark}
\end{equation}
where the octahedral-bulk closure factor is
\begin{equation}
\Delta_2 \;=\; \frac{1}{32},
\label{eq:Delta2}
\end{equation}
and $\chi_q, \lambda_q \in \{-1, 0, +1\}$ are per-quark
chirality and per-step-loss signs respectively. The per-quark
gateway-mixing sub-harmonic series $\alpha_q$ is non-zero only for the
top quark, the only cycle-2 quark in the framework. For all cycle-1
quarks $\alpha_q = 0$; for the top quark,
\begin{equation}
\alpha_t \;=\; \frac{c_2^{(t)} \alpha_\tau^2}{1 + \alpha_\tau / h_{SU(3)}},
\quad c_2^{(t)} = \biggl(\frac{h}{h_{SU(3)}}\biggr)^2 = 100,
\label{eq:alphat-derived}
\end{equation}
with $h_{SU(3)} = 3$ the colour Coxeter number. The leading coefficient
$c_2^{(t)} = (h_{E_8}/h_{SU(3)})^2$ is the topological-subgroup ratio
between the substrate's full Coxeter period and the colour Coxeter
period, squared: ten colour cycles fit inside one full substrate
cycle, and the top -- the only cycle-2 colour-charged fermion --
projects through this ratio twice (once for its colour-charge
winding, once for its cycle-2 chain-position projection). The
closure denominator $h_{SU(3)} = 3$ is the same colour-projected
closure that appears in $\alpha_\mu$ and $\alpha_H$ throughout the
framework. For the strange
quark sitting exactly at the harmonic threshold $k_s = h/2 = 15$,
the per-step-loss factor $|k_s - h/2|/(2\dim) = 0$ regardless of the
sign of $\lambda_s$, so $\lambda_s$ is indeterminate (we record it
as $\lambda_s = 0$ purely as bookkeeping notation). The chirality
factor $\chi_s = 0$, in contrast, is a \emph{distinct structural
assignment}: it forces the closure factor $(1 + \chi_s \Delta_2) =
1$ at the chirality-flip locus and is the only non-trivial sign
choice the strange quark makes:

\begin{table}[h]
\centering
\caption{Per-quark Coxeter step $k_q$, chirality sign $\chi_q$, and
per-step-loss sign $\lambda_q$ in the quark mass formula
Eq.~\eqref{eq:Rquark}.}
\label{tab:quarksigns}
\begin{tabular}{lcccc}
\toprule
Quark & $k_q$ & $\chi_q$ & $\lambda_q$ & Cycle \\
\midrule
up      & 4    & $+1$ & $+1$ & 1 \\
down    & 6.5  & $-1$ & $-1$ & 1 \\
strange & 15   & $0$  & $0$  & 1 (threshold) \\
charm   & 22.5 & $+1$ & $-1$ & 1 \\
bottom  & 26   & $-1$ & $+1$ & 1 \\
top     & 36.5 & $+1$ & $+1$ & 2 (cycle 2) \\
\bottomrule
\end{tabular}
\end{table}

The denominator $2\dim = 496$ in the per-step-loss term is the
framework's chosen normalization by the full $E_8 \times E_8$ algebra
dimension (visible plus dark sectors). All six quark masses lie within current PDG~2024 uncertainty
bands; the worst central-value residual is $+5618$~ppm on the up quark,
well within the $\sim$25\% experimental uncertainty on the
$\overline{\rm MS}$ light-quark mass at $\mu = 2$~GeV.

\subsection{Term-by-term decomposition: charged fermions}
\label{sub:fermion-terms}

Each charged-fermion mass function decomposes into four contributions:
the \emph{primary harmonic} (Coxeter step magnitude), the
\emph{sub-harmonic correction} (recursive coupling at order
$\alpha_\tau^k$), the \emph{sector modifier} (chain-type factor:
$\varphi$-chain or $\sqrt{2}$-chain), and the \emph{dark-sector
leakage} (boundary correction for cycle-2 modes only). For
cycle-1 fermions the dark-sector leakage is absent. The
decomposition for each particle:

\paragraph{Electron $(k_e = 0)$.}
$$
m_e \;=\; h_{\rm eff} \cdot \underbrace{\frac{\dim - 2}{\dim}}_{\text{gateway factor}}
$$
\begin{itemize}
\item \emph{Primary harmonic:} the gateway-factor ratio
$(\dim-2)/\dim = 246/248$ is the leading contribution. The numerator
$\dim - 2 = 246$ subtracts the two dual-sector endpoints (one in each
$E_8$); the denominator normalizes to the full algebra dimension.
\item \emph{Sub-harmonic correction:} none; the electron is the
gateway reference and absorbs all sub-harmonic structure into the
definition $h_{\rm eff} = m_e \cdot \dim/(\dim-2) = 124 m_e/123$.
\item \emph{Sector modifier:} gateway position (no chain-type
distinction).
\item \emph{Dark-sector leakage:} not applicable (cycle-1 mode at
$k = 0$).
\end{itemize}

\paragraph{Muon $(k_\mu = 11)$, $R_\mu \approx 205.10$.}
\begin{itemize}
\item \emph{Primary harmonic:} $\varphi^{11}$, the icosahedral
bulk-recursion contribution at the muon's $\varphi$-chain position
$k_\mu = 11$ (the third $E_8$ exponent, on the
$\Phi$-below-threshold side of the palindromic spectrum).
\item \emph{Icosahedral-bulk closure:} $(1 + \Delta_\varphi)$ with
$\Delta_\varphi = 1/(h + \varphi^2)$. The positive sign reflects
the muon's position below the harmonic threshold $h/2 = 15$.
\item \emph{Sub-harmonic correction:} $(1 - \alpha_\mu)$ at order
$\alpha_\tau^2$ with leading coefficient $c_2^{(\mu)} = 2$, assigned to
the palindromic pair
$(11, 19)$ centered on $h/2$.
\item \emph{Sector modifier:} $\varphi$-chain (icosahedral, lepton).
\item \emph{Dark-sector leakage:} not applicable (cycle-1 mode).
\end{itemize}

\paragraph{Tau $(k_\tau = 17)$, $R_\tau \approx 3449.19$.}
\begin{itemize}
\item \emph{Primary harmonic:} $\varphi^{17}$, the icosahedral
bulk-recursion contribution at the tau's $\varphi$-chain position
$k_\tau = 17$ (the fourth $E_8$ exponent, the first
above-the-threshold exponent).
\item \emph{Icosahedral-bulk closure:} $(1 - \Delta_\varphi)$ with
$\Delta_\varphi = 1/(h + \varphi^2)$. The negative sign reflects
the tau's position above the harmonic threshold $h/2 = 15$,
opposite to the muon.
\item \emph{Sub-harmonic correction:} $(1 - \alpha_\tau^{\rm full})$
where $\alpha_\tau^{\rm full} = \alpha_\tau/(1 + \varphi^2 \alpha_\tau)
= 1/(\dim + h + \varphi^2)$, the closed-form sum of the geometric
series in the icosahedral self-coupling $\varphi^2$.
\item \emph{Sector modifier:} $\varphi$-chain (icosahedral, lepton).
\item \emph{Dark-sector leakage:} not applicable.
\end{itemize}

\paragraph{Quarks (six entries via Eq.~\eqref{eq:Rquark}).}
The unified quark formula
$R_q = (\sqrt{2})^{k_q} \cdot (1 + \chi_q \Delta_2) \cdot
(1 + \lambda_q |k_q - h/2|/(2\dim))$ has three structural pieces
common to all six quarks, with per-quark signs $\chi_q$ and
$\lambda_q$ from Table~\ref{tab:quarksigns}:
\begin{itemize}
\item \emph{Primary harmonic:} $(\sqrt{2})^{k_q}$, the octahedral
bulk-recursion contribution at the quark's Coxeter step.
\item \emph{Octahedral-bulk closure:} $(1 + \chi_q \Delta_2)$ with
$\Delta_2 = 1/32$. The chirality sign $\chi_q$ alternates across
the spectrum: $+1$ for ``up-type'' quarks (up, charm, top) and $-1$
for ``down-type'' quarks (down, bottom). The strange quark at the
harmonic threshold has $\chi_s = 0$ (the closure factor collapses to
$1$ at the chirality-flip locus).
\item \emph{Per-step loss:} $(1 + \lambda_q |k_q - h/2|/(2\dim))$
with sign $\lambda_q \in \{-1, 0, +1\}$ per quark. The denominator
$2\dim = 496$ is the framework's normalization by the full
$E_8 \times E_8$ algebra dimension (visible plus dark sectors),
reflecting the per-step leakage across both sectors as the chain advances away from the harmonic threshold
$h/2$. The strange quark at $k_s = h/2$ has $|k_s - h/2| = 0$, so
the loss factor is automatically unity for any value of $\lambda_s$;
the entry $\lambda_s = 0$ in Table~\ref{tab:quarksigns} is a
notational placeholder, not an independent structural assignment.
\item \emph{Sector modifier:} $\sqrt{2}$-chain at integer or
half-integer position. The strange quark sits exactly at the
harmonic threshold, the chirality-flip locus.
\item \emph{Dark-sector leakage:} the top quark is the only cycle-2
quark ($k_t = 36.5 > h = 30$); for the top, the $|k_t - h/2|$
distance is computed as written and the same formula
Eq.~\eqref{eq:Rquark} applies without an additional cycle-extension
modifier.
\item \emph{Cycle-2 sub-harmonic correction (top quark only):} $\alpha_t = (h/h_{SU(3)})^2 \alpha_\tau^2 / (1 + \alpha_\tau / h_{SU(3)})$, with leading coefficient $(h/h_{SU(3)})^2 = 100$ -- the topological subgroup ratio between the substrate's full Coxeter period $h_{E_8} = 30$ and the colour Coxeter period $h_{SU(3)} = 3$, squared (one factor for the top's colour-charge winding, one for its cycle-2 chain-position projection). Closure denominator $h_{SU(3)} = 3$. The other five quarks (cycle-1) have $\alpha_q = 0$ in Eq.~\eqref{eq:Rquark}; only the top, as the unique cycle-2 colour-charged matter mode, picks up the topological subgroup ratio in its sub-harmonic series.
\end{itemize}

\section{Higgs sector and dark-sector coupling family}
\label{sec:higgs}

\subsection{Boson mass-formula schema and structural corrections}

The electroweak boson masses and the Higgs vev occupy positions on
the $\sqrt{2}$-chain (octahedral bulk recursion) in cycle 2 of the
Coxeter spectrum. Their explicit closed-form ratios are
\begin{align}
R_v \;&=\; (\sqrt{2})^{k_v} \cdot \bigl(1 - \alpha_v\bigr),
   \label{eq:Rv} \\
R_H \;&=\; \frac{(\sqrt{2})^{k_v}}{2} \cdot \bigl(1 + \delta_H\bigr) \cdot \bigl(1 - \alpha_H\bigr),
   \label{eq:RH} \\
R_W \;&=\; (\sqrt{2})^{34.5} \cdot \bigl(1 + \delta_W\bigr) \cdot \bigl(1 - \alpha_W\bigr),
   \label{eq:RW} \\
m_Z \;&=\; m_W / \cos\theta_W
\quad \text{(SM tree-level)},
   \label{eq:mZ-deriv}
\end{align}
where the Higgs-vev Coxeter position is
\begin{equation}
\begin{aligned}
k_v \;&=\; h + \text{rank} - \text{rank}/h \\
     \;&=\; 30 + 8 - 8/30 \;=\; \frac{1132}{30} \;\approx\; 37.733.
\end{aligned}
\label{eq:kv-derived}
\end{equation}
The Higgs vev and Higgs boson share the same $\sqrt{2}$-chain
primary-harmonic exponent $k_v$; the Higgs boson's amplitude is
reduced by a factor of $2$ (the breathing-mode amplitude relative to
the substrate baseline tension), and acquires the $(1 + \delta_H)$
dark-sector leakage which the vev does not. The $W$ boson sits at
the half-integer cycle-2 position $k_W = 34.5$ on the
$\sqrt{2}$-chain with leakage $(1 + \delta_W)$ suppressed by an
$E_8$ exponent factor of $19$.

The boson-sector structural corrections are
\begin{align}
\alpha_v &\;=\; \frac{c_2^{(v)} \alpha_\tau^2}{1 + \alpha_\tau / (2\varphi)},
\quad c_2^{(v)} = \frac{\text{rank} \cdot h}{m_{\max}} = \frac{|\Phi_{E_8}|}{m_{\max}} = \frac{240}{29},
   \label{eq:alphav-derived} \\
\alpha_H &\;=\; \frac{c_2^{(H)} \alpha_\tau^2}{1 + \alpha_\tau / h_{SU(3)}},
\quad c_2^{(H)} = \frac{(\text{rank} - 2) \cdot h}{m_{\max}} = \frac{180}{29},
   \label{eq:alphaH-derived} \\
\alpha_W &\;=\; \frac{c_2^{(W)} \alpha_\tau^2}{1 + \alpha_\tau / h_{SU(3)}},
\quad c_2^{(W)} = \frac{2}{e_2} = \frac{2}{7},
   \label{eq:alphaW-derived} \\
\delta_H &\;=\; \frac{1}{h + m_{\max}} \;=\; \frac{1}{59},
   \label{eq:deltaHbox} \\
\delta_W &\;=\; \frac{\delta_H}{e_6} \;=\; \frac{1}{(h + m_{\max}) \cdot 19} \;=\; \frac{1}{1121},
   \label{eq:deltaWbox}
\end{align}
where $e_6 = 19$ (1-indexed in the ordered $E_8$ exponent list
$\{e_1,\ldots,e_8\}=\{1,7,11,13,17,19,23,29\}$) is the partner of the
muon's $E_8$ exponent $e_3 = 11$ under the standard Coxeter duality
$e_j + e_{\text{rank}+1-j} = h$~\cite{Humphreys1990}. For $E_8$, this
duality groups the exponents into the four palindromic pairs
$\{(1,29), (7,23), (11,19), (13,17)\}$, all summing to
$h = 30$ and centered on $h/2 = 15$; the $W$-boson construction
selects the pair $(e_3, e_6) = (11, 19)$ containing the muon's
chain position.

The cycle-2 sub-harmonic series $\alpha_v, \alpha_H, \alpha_W$ (and
$\alpha_t$ in the quark sector, Eq.~\eqref{eq:alphat-derived}) share
a common substrate-physical structure: each is a closed-form
geometric-series sum of the form $c_2 \alpha_\tau^2 / (1 + \alpha_\tau/X)$,
with mode-specific leading coefficient $c_2$ built from $E_8$ Coxeter
slots (the eight Coxeter exponents) and the topological subgroup
ratio $h/h_{SU(3)} = 10$, and closure denominator $X$ from substrate-
physical structural integers. The Higgs vev's closure denominator is
$2\varphi$ (the icosahedral spin-doublet primary-harmonic step:
McKay multiplicity $2$ times the icosahedral recursion factor
$\varphi$), reflecting the vev's role as the substrate baseline tension
whose self-coupling is its own primary-harmonic step. The Higgs
boson, $W$ boson, and top quark closures are all $h_{SU(3)} = 3$ (the
colour-projected closure), reflecting these modes' role as
fluctuations on the colour-coupled cycle-2 vacuum. The leading
coefficients are $c_2^{(v)} = \text{rank} \cdot h / m_{\max}$ (full
Cartan rank engages for the vev), $c_2^{(H)} = (\text{rank}-2) \cdot h
/ m_{\max}$ (two Cartan directions lifted out for the breathing
mode), $c_2^{(W)} = 2/e_2$ (vector multiplicity $2$ over the
substrate's second Coxeter slot), and $c_2^{(t)} = (h/h_{SU(3)})^2$
(topological subgroup ratio squared for the only cycle-2
colour-charged fermion). Each
correction is built from framework-fixed $E_8$ invariants
$\{\dim, h, \text{rank}, m_{\max}, e_j\}$, the colour Coxeter number
$h_{SU(3)} = 3$, and the gateway-mixing parameter $\alpha_\tau$. We note that the Higgs-vev coefficient
$c_2^{(v)} = \text{rank} \cdot h$ equals the order of the $E_8$
root system $|\Phi_{E_8}| = h \cdot \text{rank} = 30 \cdot 8 = 240$,
a standard Lie-algebraic identity for any irreducible finite
Coxeter group~\cite{Humphreys1990}; the substitution
$\text{rank} \cdot h \to |\Phi_{E_8}|$ is the natural Lie-theoretic
reading of $\alpha_v$. The rank-and-rank-minus-two structures in
$\alpha_v, \alpha_H$ reflect a rank-locked-subspace projection
shared with the dark-matter tower (Sec.~\ref{sec:DM}). \emph{We use
``rank-locked subspace'' as framework-specific terminology}: it
denotes the rank-$\text{rank} = 8$ invariant subspace that arises in
the framework's bulk decomposition $\dim = \dim - \text{rank}$ visible
$\oplus$ rank-locked, distinct from any standard Cartan subalgebra
construction. The phrase is defined operationally by its appearance
in $\alpha_v$, $\alpha_H$, the dark-matter tower's
$1 - \text{rank}\,\alpha_\tau$ corrections, and certain quark
sub-harmonics; a direct identification with a specific subspace of
the Calabi--Yau cohomology is left for future work. The
structural origin of $\delta_H$ and $\delta_W$ is motivated in the
subsections below as a minimal closure rule compatible with the same
geometric data.

The Higgs vev sits at cycle-2 chain position $k_v = 1132/30 \approx 37.733$,
giving $v^{\rm pred} = 246.21965$~GeV (\SI{0.0001}{ppm}).

\paragraph{Closed-form structural identity for the vev.}
The chain-position prediction for the Higgs vev satisfies the
substrate-physical identity at two levels of precision.

\emph{Leading identity ($\sim 29$ ppb residual):}
\begin{equation}
\boxed{v \cdot (1 - \delta_W) \;=\; \dim(E_8) - 2 \;=\; 246 \;\; \text{(in GeV units)}}
\label{eq:v-structural-identity}
\end{equation}
to numerical residual $\sim 29$ ppb (at the framework's
gateway-mixing-tower depth-$3$, $\alpha_\tau^3$, scale). Equivalently,
\begin{equation}
\begin{aligned}
v \;&=\; \frac{\dim(E_8) - 2}{1 - \delta_W} \;=\; \frac{246 \cdot 1121}{1120} \\
   \;&=\; \frac{275766}{1120}~\text{GeV} \;\approx\; 246.21964~\text{GeV},
\end{aligned}
\label{eq:v-identity-form}
\end{equation}
with $\delta_W = 1/(h+m_{\max})/e_6 = 1/1121$ the W boson's vector
dark-sector leakage.

\emph{Extended identity ($\sim 0.1$ ppb residual, at floating-point precision):}
\begin{equation}
v \cdot (1 - \delta_W) \;=\; (\dim(E_8) - 2) \cdot \biggl(1 + \frac{5}{8}\alpha_\tau^3\biggr),
\label{eq:v-identity-extended}
\end{equation}
with the additional gateway-mixing-tower depth-$3$ correction
$\frac{5}{8}\alpha_\tau^3 \approx 29.1$ ppb. The leading factor $5/8$
is the substrate-physical structural ratio
\begin{equation}
\begin{aligned}
\frac{5}{8} \;&=\; \frac{\text{McKay icosahedral integer}}{\text{rank}(E_8)} \\
            \;&=\; \frac{\text{rank}(E_8) - h_{SU(3)}}{\text{rank}(E_8)} \;=\; \frac{8-3}{8},
\end{aligned}
\label{eq:5-over-8}
\end{equation}
``rank not consumed by colour, normalised by rank.'' The McKay
icosahedral integer $5$ that appears here is the same integer
appearing in the gateway formula
Eq.~\eqref{eq:gateway} as the icosahedral-seed factor.

\emph{Full structural ladder including dark-tower contributions.} The
substrate-physically complete form of the vev structural identity is
a four-factor ladder:
\begin{equation}
v \;=\; (\dim - 2) \cdot \frac{1}{1 - \delta_W} \cdot \biggl(1 + \frac{5}{8}\alpha_\tau^3\biggr) \cdot \bigl(1 + \zeta_{\rm DM}\bigr),
\label{eq:v-full-ladder}
\end{equation}
with the dark-matter tower's kinetic-mixing summed contribution
\begin{equation}
\zeta_{\rm DM} \;=\; \sum_{n=5}^{11} \epsilon_n \;=\; \sum_{n=5}^{11} \alpha_\tau^n \;=\; \alpha_\tau^5 \cdot \frac{1 - \alpha_\tau^7}{1 - \alpha_\tau} \;\approx\; 6.04 \times 10^{-13}.
\label{eq:zeta-dm}
\end{equation}
Each factor in the ladder has clear substrate-physical content and
operates at a distinct scale: $(\dim - 2) \approx 246$~GeV is the
gateway-saturated dimension count baseline; $1/(1 - \delta_W) \approx
1 + 893$ ppm is the macro-mode boundary coupling through the W-vector
channel; $\frac{5}{8}\alpha_\tau^3 \approx 29$ ppb is the
gateway-mixing-tower depth-$3$ correction; $\zeta_{\rm DM} \approx
0.6$ ppt is the dark-matter tower kinetic-mixing summed contribution.
Each factor corresponds to one of the framework's three substrate-physical
dark-sector coupling families: macro-mode boundary couplings
($\delta_W$ family), gateway-mixing-tower sub-harmonics ($\alpha_\tau^n$
for $n \geq 2$), and dark-matter tower kinetic mixings ($\epsilon_n =
\alpha_\tau^n$ for $n \geq 5$). The dark-tower factor $\zeta_{\rm DM}$
is structurally present and complete in the equation, but its
magnitude is below PDG's precision floor for $v$ (~500 ppb) and below
the framework's intrinsic comparison floor against PDG; including it
does not improve the empirical match beyond the extended identity
Eq.~\eqref{eq:v-identity-extended}, but it represents the
substrate-physically complete form. The same structural integer $\delta_W$ that
governs the W boson's mass shift in Eq.~\eqref{eq:RW} also governs
the macro-mode elevation of the substrate's electroweak baseline:
\emph{one $\delta_W$, two geometric contexts}, in the same pattern as
the gateway formula's $(\dim - 2)/\dim$ and the Weinberg angle's
$2/9$ each express the substrate gateway in different geometric
slicings of the gauge phase-space.

\emph{Live anchor versus structural baseline.} A direct empirical
test against the high-precision electron and muon masses establishes
that the substrate's mass spectrum is anchored to $v_{\rm observed}$,
not to a counterfactual $v_{\rm natural} = \dim - 2$. Anchoring the
spectrum to $\dim - 2$ predicts both leptons at $-893$ ppm against
PDG (a $\sim 10^6\sigma$ deviation on the electron); anchoring to
$v_{\rm observed}$ as in Eq.~\eqref{eq:v-identity-form} predicts both
within $\sim 0.03$ ppm. The substrate's chain-position couplings
therefore reference the effective electroweak vacuum that has the
dark-tower content already operating, not a counterfactual baseline
that would exist without it. The structural identity
Eq.~\eqref{eq:v-structural-identity} \emph{expresses} the live
anchor $v_{\rm observed}$ in closed form by factoring out the
W-channel dark-tower elevation; it is not a separate, deeper anchor
that the mass spectrum sees.

The Higgs vev is therefore the substrate's electroweak-baseline
\emph{macro-mode}: a cosmological-scale tension structure whose value
in GeV units is fixed by the gateway-saturated dimension count of
$E_8 \times E_8$ on Schoen $\mathbb{Z}_3 \times \mathbb{Z}_3$ combined
with the W-channel dark-tower elevation that lifts it to the observed
value. The framework's chain-position formula in Eq.~\eqref{eq:Rv}
reproduces this identity through different but equivalent structural
arithmetic, with numerical residual at the framework's typical
sub-ppm precision floor.

\subsection{Origin of the dark-sector coupling family}

Particles whose Coxeter positions extend into cycle 2 of the $E_8$
spectrum (positions $k > h$) couple to the dark $E_8$ sector through
the upper edge of the cycle-1 Coxeter spectrum. Substrate continuity
across the dark/visible boundary forces a structural correction to
their visible-sector mass formulas, equal to the inverse of the total
Coxeter capacity available at the edge.

The structural ``boundary capacity'' of the $E_8$ Coxeter spectrum is
the sum of the Coxeter number $h$ and the largest exponent $m_{\max}$:
\begin{equation}
h + m_{\max} = 30 + 29 = 59.
\label{eq:boundarycap}
\end{equation}
This quantity uses two genuine $E_8$ invariants: $h$ is the order of
the Coxeter element (the period of the spectrum) and $m_{\max}$ is the
highest non-trivial $E_8$ exponent (the upper exponent in the
palindromic pair $(1, 29)$). Their sum is then a framework measure of
the cycle-1 edge plus maximal-exponent displacement, not a standard
Lie-algebra invariant in its own right.

\subsection{The Higgs leak: structural closure rule for $\delta_H$}

The Higgs boson is the breathing-mode excitation of the substrate
tension at $k_v - 2$ (one ``inversion-step'' below the vev). As a
scalar mode-pattern, it couples isotropically to the boundary capacity
at first order. A minimal scalar closure consistent with that boundary
picture assigns the leading dark-sector correction
\begin{equation}
\boxed{\delta_H \;=\; \frac{1}{h + m_{\max}} \;=\; \frac{1}{59}}
\label{eq:deltaH_derived}
\end{equation}
($\approx \SI{1.695}{\percent}$). The integer $1$ in the numerator is
the framework's minimal choice: a scalar mode has no internal index
structure (no rank or exponent dependence) that would naturally supply
an additional suppression factor. The factor $1/(h + m_{\max})$ then
enters the Higgs-boson mass formula as
the ``boundary leakage'' multiplicative correction at the breathing
mode position $k_v - 2$, yielding $m_H^{\rm pred} = 125.19978$~GeV
(\SI{1.75}{ppm} from PDG). This should be read as a structurally
motivated closure rule awaiting derivation from the full compactification
data, not as a completed CY3 cohomology computation.

\subsection{The $W$ leak: structural closure rule for $\delta_W$}

The $W$ boson is a vector mode-pattern at cycle-2 Coxeter position
$k_W = 34.5$, which lies on the $\sqrt{2}$ half-integer chain shared
with the heavy-quark sector. As a vector mode, the $W$ carries an
internal index structure: its phase-alignment under $SU(2)$ projects
onto a specific palindromic exponent pair on the $\varphi$-chain,
namely the muon-paired pair $(11, 19)$ with $11 + 19 = h = 30$.

The corresponding vector closure rule suppresses the dark-sector
boundary leakage for the $W$ by the partner $E_8$ exponent across the
palindromic pairing,
$e_n = 19$ (the partner of $11$, the muon's Coxeter position):
\begin{equation}
\boxed{\delta_W \;=\; \frac{\delta_H}{e_n} \;=\; \frac{1}{(h + m_{\max}) \cdot e_n} \;=\; \frac{1}{59 \cdot 19} \;=\; \frac{1}{1121}}
\label{eq:deltaW_derived}
\end{equation}
($\approx \SI{892}{ppm}$). Each factor in the denominator is fixed by
$E_8$ structure once this closure pattern is adopted:
\begin{itemize}
\item $h = 30$ --- Coxeter number (algebra-fixed).
\item $m_{\max} = 29$ --- largest $E_8$ exponent (algebra-fixed).
\item $e_n = 19$ --- partner exponent of the muon's position $k_\mu = 11$
  under the palindromic pairing $m \leftrightarrow h - m$ centered on
  $h/2 = 15$ (algebra-fixed once the muon's Coxeter step is identified).
\end{itemize}
The ratio $\delta_W/\delta_H = 1/e_n = 1/19$ is therefore
independently fixed by the $E_8$ exponent structure with no fitted
continuous parameter; the same palindromic-pair structure that governs the
muon-tau lepton mass relations governs the suppression of the $W$
boson's dark-sector coupling relative to the Higgs's. Applying
$\delta_W$ to the cycle-2 $W$ position yields
$m_W^{\rm pred} = 80.36950$~GeV (\SI{3.7}{ppm} from PDG). As with
$\delta_H$, this is best read as a framework closure rule sharpened by
the $E_8$ exponent structure rather than as a fully first-principles
compactification derivation.

\subsection{Generalization to other cycle-2 modes}

The same closure family extends to any cycle-2 Coxeter mode: a scalar
in cycle 2 receives the bare $\delta_H$ correction; a vector or
spinor in cycle 2 receives the suppressed correction
$\delta_X = \delta_H / e_X$ where $e_X$ is the $E_8$ exponent
determined by the mode's internal index structure. The dark-sector
coupling family is therefore graded by the $E_8$ exponent set
$\{1, 7, 11, 13, 17, 19, 23, 29\}$, producing a discrete spectrum
of structurally-fixed corrections rather than a continuous parameter.
The Higgs self-coupling $\lambda_H = m_H^2/(2 v^2)$ follows from the
predicted $m_H$ and $v$ by the standard SM tree-level relation; its
closed form and numerical value are presented in
Sec.~\ref{sec:EW}.

\subsection{Term-by-term decomposition: electroweak bosons}
\label{sub:boson-terms}

The electroweak boson masses share the same four-component
structure as the charged fermions (primary harmonic, sub-harmonic
correction, sector modifier, dark-sector leakage), but all three
massive bosons sit in cycle 2 of the Coxeter spectrum, so
dark-sector leakage corrections are present for all of them.

\paragraph{Higgs vacuum expectation value $(k_v = 1132/30)$, $R_v \approx 4.78 \times 10^5$.}
\begin{itemize}
\item \emph{Primary harmonic:} cycle-2 chain contribution at the
chain-position assignment $k_v = 1132/30 \approx 37.733$. The
denominator $30 = h$ is the Coxeter number; the numerator $1132$ is
a structural-integer assignment from the heterotic Wilson-line
breaking pattern (Ref.~\cite{BraunHeOvrutPantev2005}; explicit
derivation from CY3 cohomology is outside the present scope).
\item \emph{Sub-harmonic correction:} $\alpha_v = ({\rm rank} \cdot
h/m_{\max}) \alpha_\tau^2 = (8 \cdot 30/29) \alpha_\tau^2$. The
coefficient ${\rm rank} \cdot h/m_{\max}$ encodes the rank-locked
Cartan-subspace projection at order $\alpha_\tau^2$ (paralleling the
muon's order-2 sub-harmonic).
\item \emph{Sector modifier:} cycle-2 chain.
\item \emph{Dark-sector leakage:} absent for the vev itself, since
the Higgs vev is the substrate baseline tension and does not couple
to the spectrum's upper boundary as a fluctuation mode would.
\end{itemize}

\paragraph{Higgs boson, $R_H \approx 2.43 \times 10^5$.}
\begin{itemize}
\item \emph{Primary harmonic:} same $\sqrt{2}$-chain primary
$(\sqrt{2})^{k_v}$ as the vev, divided by $2$ (the breathing-mode
amplitude is half of the substrate baseline tension; the Higgs boson
is the fluctuation mode around the vev, not a separate Coxeter
position).
\item \emph{Sub-harmonic correction:} $\alpha_H = (({\rm rank}-2) h /
m_{\max}) \alpha_\tau^2 = (180/29) \alpha_\tau^2$. The
${\rm rank}-2 = 6$ coefficient drops two Cartan generators relative to
the vev correction $\alpha_v$, reflecting the two fewer rank-locked
directions accessible to the breathing mode.
\item \emph{Sector modifier:} cycle-2 $\sqrt{2}$-chain.
\item \emph{Dark-sector leakage:} $(1 + \delta_H)$ with
$\delta_H = 1/(h + m_{\max}) = 1/59$. The Higgs boson, as a scalar,
couples isotropically to the spectrum's upper boundary at the
boundary-capacity scale $h + m_{\max} = 59$, giving the leading
dark-sector leakage with no internal-index suppression. The sign of
$\delta_H$ is positive (the leakage \emph{enhances} the boson's
amplitude, in contrast to the rank-locked correction $\alpha_H$ which
reduces it).
\end{itemize}

\paragraph{$W$ boson $(k_W = 34.5)$, $R_W \approx 1.56 \times 10^5$.}
\begin{itemize}
\item \emph{Primary harmonic:} cycle-2 $\sqrt{2}$-chain contribution
$(\sqrt{2})^{34.5}$. The half-integer Coxeter position reflects the
$W$'s SU(2) vector index structure (rather than scalar or
integer-spinor).
\item \emph{Sub-harmonic correction:} $\alpha_W = (2/e_2) \alpha_\tau^2 / (1 + \alpha_\tau / h_{SU(3)})$, with leading coefficient $2/e_2 = 2/7$ (vector multiplicity $2$ over the substrate's second Coxeter slot $e_2 = 7$, the same slot that appears in $\sin^2\theta_{23}^{\rm PMNS} = \mathrm{rank}/(2 e_2) = 4/7$ and $\sin^2\theta_{13}^{\rm PMNS} = 2/(e_2 e_4) = 2/91$), and closure denominator $h_{SU(3)} = 3$ (colour-projected closure shared with $\alpha_H$ and $\alpha_\mu$).
\item \emph{Sector modifier:} cycle-2 $\sqrt{2}$-chain at half-integer
position, shared with the heavy-quark $\sqrt{2}$-chain structure.
\item \emph{Dark-sector leakage:} $(1 + \delta_W)$ with
$\delta_W = \delta_H/e_6 = 1/(59 \cdot 19) = 1/1121$, where $e_6 = 19$
is the $E_8$ exponent palindromically paired with the muon's
position $k_\mu = 11 = e_3$ (since $11 + 19 = h$). The factor of $19$ suppression
relative to $\delta_H$ encodes the $W$'s internal SU(2) vector
index structure projecting onto the $(11, 19)$-paired exponent. The
sign of $\delta_W$ is positive, like $\delta_H$.
\end{itemize}

\paragraph{$Z$ boson and gauge couplings.}
The $Z$ boson, $g_2$, $g_1$, and $\lambda_H$ are derived bookkeeping
quantities: $m_Z = m_W/\cos\theta_W$, $g_2 = 2 m_W/v$,
$g_1 = g_2 \tan\theta_W$, $\lambda_H = m_H^2/(2v^2)$. These are not
independent structural predictions; they inherit precision from
$m_W$, $v$, $m_H$, and $\sin^2\theta_W$ via standard SM tree-level
relations.

\section{Electroweak gauge sector}
\label{sec:EW}

\subsection{Weak mixing angle}

The weak mixing angle takes the closed-form structural expression
\begin{equation}
\boxed{\sin^2 \theta_W \;=\; \frac{2}{9} + \frac{\Delta_\varphi}{h}
                          \;=\; \frac{2}{9} + \frac{1}{h(h + \varphi^2)},}
\label{eq:sinW-derived}
\end{equation}
where $\varphi = (1+\sqrt{5})/2$ is the golden ratio and
$\Delta_\varphi = 1/(h + \varphi^2)$. The leading rational term $2/9$
is, in this framework, the $SU(3)_{\rm color}$-induced bulk
projection at the unbroken-electroweak-symmetry locus
(the geometric weight $2/(h_{SU(3)} \cdot h_{SU(3)}^{\rm extended}) =
2/(3 \cdot 3)$ assigned to the visible-sector
$SU(2)_L \times U(1)_Y$ subprojection within the bulk symmetry
breaking ansatz). We flag this as a \emph{framework structural
assignment} rather than a standard derived $E_8$ identity: it is the
parsimonious choice that reproduces the observed value to within the
small icosahedral correction below, and the assignment must
ultimately be justified at the level of the underlying
$E_8 \times E_8$ heterotic compactification on $X/(\mathbb{Z}_3 \times
\mathbb{Z}_3)$. The small correction $\Delta_\varphi/h$ is the
golden-ratio-conjugate sub-harmonic contribution from the visible
$E_8$ embedding, with the same icosahedral provenance as the
McKay-correspondence integer $5$ in the gateway formula.
Numerically, $\sin^2 \theta_W = 0.222\overline{2} + 0.001022 = 0.22324$,
agreeing with the on-shell PDG value $1 - m_W^2/m_Z^2 = 0.22320$ at
\SI{198}{ppm}.

\subsection{Gauge couplings and $Z$ boson}

The $SU(2)_L$ and $U(1)_Y$ gauge couplings follow from the predicted
$m_W$, $v$, and $\theta_W$ via the standard SM tree-level relations:
\begin{align}
g_2 &\;=\; \frac{2 m_W}{v}
        \;=\; 0.65283,
   \label{eq:g2formula} \\
g_1 &\;=\; g_2 \tan\theta_W
        \;=\; \frac{2 m_W \tan\theta_W}{v}
        \;=\; 0.34998,
   \label{eq:g1formula} \\
e \;&=\; g_2 \sin\theta_W,
\qquad
\alpha_{\rm em} \;=\; \frac{e^2}{4\pi\epsilon_0 \hbar c}.
   \label{eq:eformula}
\end{align}
The $Z$ boson mass is
\begin{equation}
\begin{aligned}
m_Z \;&=\; \frac{m_W}{\cos\theta_W} \\
     &=\; \frac{h_{\rm eff} \cdot k_W \cdot (1 + \delta_W) \cdot K_{\text{cycle},W}}{\cos\theta_W} \\
     &=\; 91.19037~\text{GeV},
\end{aligned}
\label{eq:mZformula2}
\end{equation}
agreeing with PDG~2024 at \SI{30}{ppm}.

\paragraph{Heterotic-string expectation: gauge unification.}
A heterotic $E_8 \times E_8$ compactification on a Calabi--Yau
threefold normally predicts that the running visible-sector gauge
couplings $\alpha_1(\mu), \alpha_2(\mu), \alpha_3(\mu)$ unify at a
GUT scale set by the compactification volume,
$M_{\rm GUT} \sim 10^{16}$~GeV in conventional units, with
$\alpha_{\rm GUT} \sim 1/24$. Within the present framework, $g_1$,
$g_2$, and $g_3$ are predicted at the $Z$ pole as separate
structural objects rooted in distinct $E_8$ invariants ($k_W$,
$\sin^2\theta_W$, $\alpha_s$), and the framework does not predict
their evolution to high scale: the reconciliation of the predicted
low-scale couplings with one-loop unification of the running
couplings is left for future work, contingent on a specific
identification of the substrate-physical interpretation of RG
running. We note that the $Z$-pole prediction
$\alpha_s(M_Z) = 0.1180$ is consistent at the $5\%$ level with
unification at $M_{\rm GUT}$ given two-loop SM running and
threshold corrections at $M_{\rm SUSY} \sim 1$~TeV; no claim
beyond consistency is made.

\subsection{Higgs self-coupling}

The Higgs self-coupling is the standard SM tree-level relation
inverted to read off $\lambda_H$:
\begin{multline}
\lambda_H \;=\; \frac{m_H^2}{2 v^2} \\
            \;=\; \frac{[h_{\rm eff} (k_v - 2)(1 - \alpha_H)(1 + \delta_H)]^2}{2 [h_{\rm eff} k_v (1 - \alpha_v)]^2}
            \;=\; 0.12928011,
\label{eq:lambdaH-derived}
\end{multline}
agreeing with the observed $\lambda_H = 0.12928057$ at \SI{3.5}{ppm}.

\subsection{Strong coupling}

The strong coupling at the $Z$ pole takes the closed-form structural
expression
\begin{equation}
\boxed{\begin{aligned}
\alpha_s(M_Z) &\;=\; \frac{m_{\max}^2 + \text{rank}}{m_{\max} \cdot \dim}
               \;=\; \frac{29^2 + 8}{29 \cdot 248} \\
              &\;=\; \frac{849}{7192} \;=\; 0.118048,
\end{aligned}}
\label{eq:alphas-derived}
\end{equation}
which uses only forced $E_8$ algebra invariants
$\{\dim, \text{rank}, m_{\max}\}$. The numerator $m_{\max}^2 +
\text{rank}$ combines the squared largest exponent
$m_{\max}^2 = (h-1)^2 = 841$ with the algebra rank
$\text{rank} = 8$; the denominator $m_{\max} \cdot \dim = 29 \cdot
248 = 7192$ is the largest-exponent times the algebra dimension.
\emph{We flag this as a framework structural assignment.} Although
the numerator and denominator are each built from $E_8$ invariants,
neither $m_{\max}^2 + \text{rank}$ nor $m_{\max} \cdot \dim$
corresponds directly to a standard Lie-algebraic invariant
(quadratic Casimir, Killing form, dimension of an irrep, etc.):
the quadratic Casimir of the adjoint of $E_8$ is
$C_2(\text{adj}) = 2 h^\vee = 60$ (with dual Coxeter number
$h^\vee = h$ for simply-laced groups), and no standard identity
recovers $849/7192$ from it. The construction selects the unique
combination of $\{\dim, \text{rank}, m_{\max}\}$ at a fixed weight
that reproduces the PDG value to within experimental precision; a
first-principles identification (e.g.\ via a specific
gauge-boundary-condition trace on $X/(\mathbb{Z}_3 \times
\mathbb{Z}_3)$) is left for future work. Numerically
$\alpha_s(M_Z) = 0.118048$, agreeing with the PDG~2024 world
average $\alpha_s(M_Z) = 0.1180 \pm 0.0009$ at \SI{405}{ppm}, well
within experimental precision. This is the only closed-form
structural expression for $\alpha_s(M_Z)$ in the present framework;
no continuous parameter enters.

\subsection{Fine-structure constant}

The fine-structure constant follows from $g_2$ and $\theta_W$ via
$e = g_2 \sin\theta_W$ and $\alpha_{\rm em} = e^2/(4\pi)$ in
natural units. With Eqs.~\eqref{eq:g2formula} and~\eqref{eq:sinW-derived},
\begin{equation}
\alpha_{\rm em} \;=\; \frac{(g_2 \sin\theta_W)^2}{4\pi}
                  \;=\; \frac{(2 m_W \sin\theta_W / v)^2}{4\pi},
\label{eq:alphaem-derived}
\end{equation}
inheriting precision from $m_W, v$, and $\sin\theta_W$. The
on-shell predicted value matches the PDG within experimental
precision (Sec.~\ref{sec:results}).

\section{Strong CP angle}
\label{sec:strongCP}

The $\mathbb{Z}_3 \times \mathbb{Z}_3$ quotient structure provides a
geometric route to a vanishing QCD vacuum angle, but the required
cancellation is not derived independently here. The QCD
topological term integrates as
\begin{equation}
\int_{X/(\mathbb{Z}_3 \times \mathbb{Z}_3)} \text{Tr}(F \wedge F)
= \frac{1}{|\mathbb{Z}_3 \times \mathbb{Z}_3|} \int_X
\text{Tr}_g(F \wedge F),
\end{equation}
where $\text{Tr}_g$ is twisted by the group action. For the canonical
$E_8$ embedding and Wilson-line data~\cite{BraunHeOvrutPantev2005,
DonagiOvrut2005}, the present framework assumes that the quotient
projection annihilates the relevant twisted trace. The mathematical
content of this assumption is the equivariant-cohomology identity
\begin{equation}
\bigl[\text{ch}_2(V) \cdot [g]\bigr]_X \;=\; 0
\quad \text{for every nontrivial}\; g \in \mathbb{Z}_3 \times \mathbb{Z}_3,
\label{eq:equivariant-trace}
\end{equation}
where $\text{ch}_2(V) = \tfrac{1}{2}\,\text{Tr}(F \wedge F)/(2\pi)^2$
is the second Chern character of the gauge bundle $V$ on the cover
$X$ and $[g]$ is the cohomology class of the fixed locus of $g$
(empty for a freely-acting $\mathbb{Z}_3 \times \mathbb{Z}_3$ on the
Schoen threefold, which is the case of interest here).
Equivalently, the assumption is that the descent of $\text{Tr}(F \wedge F)$
to $X/(\mathbb{Z}_3 \times \mathbb{Z}_3)$ admits no non-zero
$g$-twisted contribution under the orbifold-cohomology
decomposition. Under that assumption,
$\text{Tr}_{\mathbb{Z}_3 \times \mathbb{Z}_3} \text{Tr}(F \wedge F) = 0$,
which would set $\theta_{\rm QCD} = 0$. The heterotic construction papers fix
the quotient geometry and bundle data on which this statement depends,
but they do not themselves prove
Eq.~\eqref{eq:equivariant-trace} for the specific bundle data the
framework needs. Within the
framework, this therefore functions as a geometric strong-CP mechanism
rather than as a separately established theorem. If the cancellation
holds, the strong CP problem is removed without Peccei--Quinn~\cite{PecceiQuinn1977}
or a dynamical axion~\cite{KimCarosi2010}.

\section{Neutrino sector}
\label{sec:neutrinos}

\subsection{Mass eigenvalues from sub-gateway $\varphi$-chain modes}

The neutrino mass eigenvalues arise as deeply sub-gateway
$\varphi$-chain modes at recursive depths $n = 3, 4$ in $\alpha_\tau$.
Their explicit closed-form ratios are
\begin{align}
R_{\nu_1} \;&=\; 2 \, \varphi^2 \, \alpha_\tau^4,
   \label{eq:Rnu1} \\
R_{\nu_3} \;&=\; 2 \, \alpha_\tau^3 \, \bigl(1 + 5 \varphi^2 \alpha_\tau\bigr),
   \label{eq:Rnu3} \\
R_{\nu_2} \;&=\; R_{\nu_3} \cdot \frac{1 + \Delta_\varphi}{\text{rank} - 2}
            \;=\; R_{\nu_3} \cdot \frac{1 + \Delta_\varphi}{6},
   \label{eq:Rnu2}
\end{align}
where $\nu_1$ is the lightest mass eigenstate (4th-order icosahedral
mode with $\varphi^2$ prefactor), $\nu_3$ is the heaviest (3rd-order
with sub-leading icosahedral correction $5\varphi^2\alpha_\tau$),
and $\nu_2$ relates to $\nu_3$ by the icosahedral-bulk closure
$(1+\Delta_\varphi)$ divided by $\text{rank} - 2 = 6$ (the number of
Cartan directions accessible to the second-generation neutrino mode).
Numerically, $m_{\nu_1} \approx 0.451$~meV, $m_{\nu_2} \approx 8.61$~meV,
$m_{\nu_3} \approx 50.21$~meV.

\subsection{Mass-squared splittings and bounds}

The mass-squared splittings derived from
Eqs.~\eqref{eq:Rnu1}--\eqref{eq:Rnu2}:
$\Delta m^2_{21} = 7.4193 \times 10^{-5}$~eV$^2$ (\SI{90}{ppm}),
$\Delta m^2_{31} = 2.5211 \times 10^{-3}$~eV$^2$ (\SI{2440}{ppm}),
agreeing with the NuFIT~5.3 global-fit ranges~\cite{Esteban2020,NuFIT53}
within experimental precision.

Normal mass ordering is selected with $\Sigma m_\nu = 59.3$~meV,
KATRIN effective mass $m_\beta = 8.82$~meV, and
$m_{\beta\beta} = 4.00$~meV (assuming Majorana nature with vanishing
Majorana phases). These remain below current cosmological,
direct-kinematic, and neutrinoless-double-beta-decay bounds from
DESI, KATRIN, and KamLAND-Zen~\cite{DESI2024,KATRIN2024,KamLANDZen2024}.

\paragraph{Dirac vs Majorana nature.}
The framework's mass-eigenvalue construction ($R_{\nu_1}, R_{\nu_2},
R_{\nu_3}$) does not by itself distinguish Dirac from Majorana
neutrinos: the closed-form ratios specify the three positive
eigenvalues of $|m_\nu|$, not the Lorentz-structure of the mass
term. The framework's heterotic $E_8 \times E_8$ embedding on
$X/(\mathbb{Z}_3 \times \mathbb{Z}_3)$ admits both possibilities,
with right-handed neutrinos arising as either standalone Dirac
partners or as integrated-out heavy seesaw modes producing an
effective Majorana mass at low scale. The
$m_{\beta\beta} = 4.00$~meV prediction quoted above is the
Majorana-case value with vanishing Majorana phases; the
corresponding $m_{\beta\beta}$ in the Dirac case is identically
zero. A future framework-internal derivation determining
which case obtains would convert
$m_{\beta\beta}$ from a conditional prediction into a direct
falsifiable observable; the present construction commits to the
mass-eigenvalue triplet only.

\subsection{PMNS mixing angles and CP phase}

The PMNS lepton-mixing parameters take explicit closed forms
involving $E_8$ exponents and the $SU(3)_{\rm color}$ Coxeter number
$h_{SU(3)} = 3$:
\begin{align}
\sin^2 \theta_{12}^{\rm PMNS} \;&=\; \frac{h}{h_{SU(3)}} \cdot \Delta_\varphi
                                \;=\; \frac{10}{h + \varphi^2},
   \label{eq:sin2theta12PMNS} \\
\sin^2 \theta_{23}^{\rm PMNS} \;&=\; \frac{\text{rank}}{2 \, e_2}
                                \;=\; \frac{4}{7},
   \label{eq:sin2theta23PMNS} \\
\sin^2 \theta_{13}^{\rm PMNS} \;&=\; \frac{2}{e_2 \cdot e_4}
                                \;=\; \frac{2}{91},
   \label{eq:sin2theta13PMNS} \\
\delta_{\rm CP} \;&=\; \frac{e_3 \cdot h_{SU(3)} \cdot \pi}{h}
                  \;=\; \frac{11 \cdot 3 \cdot \pi}{30}
                  \;=\; \frac{11\pi}{10},
   \label{eq:dCP-derived}
\end{align}
where $e_2 = 7, e_3 = 11, e_4 = 13$ are the second, third, and
fourth $E_8$ Coxeter exponents respectively. The CP phase
$\delta_{\rm CP} = 11\pi/10 = 198\degree$ has numerator equal to
the muon's Coxeter step $k_\mu = 11$ (= $e_3$) times $h_{SU(3)} = 3$,
denominator $h = 30$. All four PMNS predictions agree with NuFIT~5.3
within experimental precision. The small-integer combinations of
$\{$rank, $e_j$, $h$, $h_{SU(3)}$, $\Delta_\varphi\}$ chosen for
each PMNS angle are constrained but not uniquely forced by $E_8$
algebra alone; the framework's structural choices for these
combinations are dictated by the requirement that the closed forms
match observation, and a first-principles derivation from the
underlying CY3 cohomology is left for future work.

\section{Quark mixing (CKM)}
\label{sec:CKM}

The CKM matrix elements take explicit closed forms using the same
structural integers:
\begin{align}
\sin^2 \theta_{12}^{\rm CKM} \;&=\; \frac{h/2 - 1}{\dim + h}
                                 \;=\; \frac{14}{278},
   \label{eq:CKM12-derived} \\
\sin^2 \theta_{23}^{\rm CKM} \;&=\; \frac{4 \sin^2 \theta_{12}^{\rm CKM}}{e_3^2} \cdot \Bigl(1 + \frac{\Delta_\varphi}{h_{SU(3)}}\Bigr),
   \label{eq:CKM23-derived} \\
\sin^2 \theta_{13}^{\rm CKM} \;&=\; \alpha_\tau^2 \cdot \bigl(1 + 4 \Delta_\varphi\bigr),
   \label{eq:CKM13-derived} \\
\delta_{\rm CKM} \;&=\; \frac{e_3 \cdot \pi}{h}
                  \;=\; \frac{11\pi}{30}
                  \;=\; 66.0\degree.
   \label{eq:dCKM-derived}
\end{align}
The Cabibbo angle uses the chain depth $h/2 - 1 = M = 14$; the
atmospheric mixing involves the muon-exponent $e_3 = 11$ squared in
the denominator with an icosahedral-bulk correction $\Delta_\varphi /
h_{SU(3)}$; the reactor mixing
$\sin^2\theta_{13}^{\rm CKM} = \alpha_\tau^2 (1 + 4\Delta_\varphi)$
is at order $\alpha_\tau^2$ with a four-fold icosahedral correction.
All four CKM predictions agree with PDG~2024 within experimental
precision. The leptonic CP phase $\delta_{\rm CP} = 11\pi/10$ and
the quark CP phase $\delta_{\rm CKM} = 11\pi/30$ share the same
$E_8$-exponent numerator $11$, differing only by the factor of
$h_{SU(3)} = 3$ that distinguishes the lepton sector from the
quark sector. As with the PMNS angles, the specific small-integer
combinations chosen for each CKM angle (e.g., the factor of $4$ in
Eq.~\eqref{eq:CKM23-derived}, the factor of $4$ in front of
$\Delta_\varphi$ in Eq.~\eqref{eq:CKM13-derived}) are constrained but
not uniquely forced by $E_8$ algebra; their first-principles
derivation from the heterotic CY3 cohomology is left for future
work.

\section{Dark-matter tower}
\label{sec:DM}

\subsection{The unified neutrino-DM Coxeter tower}

The framework's Coxeter chain produces a unified neutrino-and-dark-matter
tower:
\begin{equation}
\boxed{m_n = 2 \, h_{\rm eff} \, \alpha_\tau^n,}
\label{eq:DMtower}
\end{equation}
with $\alpha_\tau = 1/(\dim+h) = 1/278 \approx 3.5971 \times 10^{-3}$ and
the icosahedral spin-doublet prefactor $c_n = 2$ shared with $\nu_3$.
The electron-anchor value of $h_{\rm eff}$ is
$0.5151534130081301$~MeV (the muon-anchor differs by 0.0011~ppm).

Each order $n$ defines a structurally distinct mode with mass
$m_n$, characteristic kinetic-mixing parameter
$\epsilon_n \equiv \alpha_\tau^n$, and de~Broglie wavelength
$\lambda_{\rm dB,n} = h/(m_n v_{\rm DM})$ at the galactic virial
velocity $v_{\rm DM} = 220$~km/s $= 7.3384 \times 10^{-4}\,c$.

\subsection{Cosmological cutoffs and viable mass window}

Two cosmological cutoffs bracket the viable dark-matter window:

\begin{itemize}
\item \textbf{Lyman-$\alpha$ cutoff}: ultralight bosonic
dark-matter modes with $m \lesssim 2 \times 10^{-20}$~eV are excluded
at 95\% C.L. by Lyman-$\alpha$ forest constraints~\cite{RogersPeiris2021}.
\item \textbf{Hubble coherence cutoff}: modes lighter than
$H_0 \sim 1.4 \times 10^{-33}$~eV cannot oscillate coherently in the
matter-dominated era and are believed to behave as dark energy rather
than dark matter; the precise role of substrate modes at orders
$n \geq 13$ in the present framework, including any contribution to
the cosmological constant, will require further research.
\end{itemize}

These cutoffs partition the tower (Eq.~\eqref{eq:DMtower}) into three
regimes: visible-sector neutrinos at orders $n = 3, 4$; dark matter at
orders $n = 5$--$11$ (seven discrete modes); and excluded modes at
$n \geq 12$. Table~\ref{tab:dark-matter-tower} lists every order from $n=3$ to
$n=15$ with its mass, kinetic-mixing parameter, de~Broglie wavelength
at $v_{\rm DM}$, and physical regime.

% Auto-generated by substrate_mass_spectrum.py
% Do not edit by hand; regenerate via:
%   python substrate_mass_spectrum.py --tex-output <DIR>
\begin{table*}[t]
\centering
\caption{Unified neutrino--dark-matter Coxeter tower from order $n=3$ (heaviest visible neutrino) through order $n=15$ (super-Hubble). Mass: $m_n = 2 h_{\rm eff} \alpha_\tau^n$ with $\alpha_\tau = 1/(\dim+h) = 1/278$. Kinetic-mixing parameter: $\epsilon_n = \alpha_\tau^n$. de~Broglie wavelength at galactic virial velocity $v_{\rm DM} = 220$~km/s.}
\label{tab:dark-matter-tower}
\begin{ruledtabular}
\begin{tabular}{cclllr}
$n$ & Prefactor & Mass & $\lambda_{\rm dB}$ & $\epsilon_n$ & Regime \\
\hline
$3$ & $2\,\alpha_\tau^{3}$ & $47.9548$ meV & $5.6073$ mm & $4.654e-08$ & (nu\_3, nu\_2, visible) \\
$4$ & $2\,\alpha_\tau^{4}$ & $172.4993$ ueV & $1.5588$ m & $1.674e-10$ & (nu\_1,     visible) \\
$5$ & $2\,\alpha_\tau^{5}$ & $620.5012$ neV & $433.3531$ m & $6.022e-13$ & ALP / dark-photon (ADMX) \\
$6$ & $2\,\alpha_\tau^{6}$ & $2.2320$ neV & $120.4722$ km & $2.166e-15$ & ALP / atomic-clock DM \\
$7$ & $2\,\alpha_\tau^{7}$ & $8.0288$ peV & $33491.2600$ km & $7.793e-18$ & atomic-clock DM \\
$8$ & $2\,\alpha_\tau^{8}$ & $28.8807$ feV & $9310570.2763$ km & $2.803e-20$ & atomic-clock / pulsar-timing \\
$9$ & $2\,\alpha_\tau^{9}$ & $103.8875$ aeV & $17.3017$ AU & $1.008e-22$ & pulsar-timing / asteroid \\
$10$ & $2\,\alpha_\tau^{10}$ & $373.6961$ zeV & $4809.8804$ AU & $3.627e-25$ & wave-DM \\
$11$ & $2\,\alpha_\tau^{11}$ & $1.3442$ zeV & $6.4821$ pc & $1.305e-27$ & wave-DM, sub-galactic \\
$12$ & $2\,\alpha_\tau^{12}$ & $4.8354$ yeV & $1.8020$ kpc & $4.693e-30$ & EXCLUDED by Lyman-alpha \\
$13$ & $2\,\alpha_\tau^{13}$ & $17.3934$ rdeV & $500.9615$ kpc & $1.688e-32$ & EXCLUDED (super-Hubble) \\
$14$ & $2\,\alpha_\tau^{14}$ & $0.0626$ rdeV & $139.2673$ Mpc & $6.073e-35$ & EXCLUDED (super-Hubble) \\
$15$ & $2\,\alpha_\tau^{15}$ & $0.0002$ rdeV & $38716.3076$ Mpc & $2.184e-37$ & EXCLUDED (super-Hubble) \\
\end{tabular}
\end{ruledtabular}
\end{table*}


\subsection{Order-5 dominant mode: couplings and detection prospects}

The dominant dark-matter mode is at $n=5$:
\begin{equation}
m_5 = 2 \, h_{\rm eff} \, \alpha_\tau^5 = 620.5012~\text{neV},
\end{equation}
with structurally-fixed kinetic mixing
\begin{equation}
\epsilon_5 = \alpha_\tau^5 = 6.022 \times 10^{-13}
\end{equation}
and axion-photon coupling derived from $\epsilon_5$, $\alpha_{\rm em}$,
and the predicted Higgs vev:
\begin{equation}
g_{a\gamma\gamma} = \frac{\epsilon_5 \, \alpha_{\rm em}}{v}
                  = 1.7849 \times 10^{-17}~\text{GeV}^{-1}.
\label{eq:gagg}
\end{equation}
The mass falls in the upper end of the ADMX direct-detection
sensitivity window for axion-like and dark-photon dark-matter searches.

\subsection{Geometric-series abundance convergence}

The same factor $\alpha_\tau = 1/278$ that produces successive mass
levels also fixes the relative cosmological abundance. The order-5
contribution is dominant; orders $6$--$11$ contribute via the convergent
geometric series:
\begin{equation}
\zeta_{\rm tower} \equiv \sum_{n=5}^{11} \frac{\Omega_n}{\Omega_5}
= \sum_{k=0}^{6} \alpha_\tau^k \approx 1.0638.
\end{equation}
The convergence is rapid (Table~\ref{tab:DMabundance}); orders $6$--$11$
contribute only $\approx 6.0\%$ of the total dark-matter abundance, and
order~5 contributes $\approx 94.0\%$.

\begin{table}[h]
\centering
\caption{Geometric-series abundance convergence over the viable
dark-matter window $n = 5$--$11$. The cumulative ratio
saturates at $\zeta_{\rm tower} \approx 1.0638$.}
\label{tab:DMabundance}
\begin{tabular}{ccc}
\toprule
$n$ & $\Omega_n / \Omega_5$ & Cumulative $\sum_{n'=5}^n \Omega_{n'}/\Omega_5$ \\
\midrule
5  & $1$                       & $1.000000$ \\
6  & $5.998 \times 10^{-2}$    & $1.059976$ \\
7  & $3.597 \times 10^{-3}$    & $1.063573$ \\
8  & $2.157 \times 10^{-4}$    & $1.063789$ \\
9  & $1.294 \times 10^{-5}$    & $1.063802$ \\
10 & $7.760 \times 10^{-7}$    & $1.063803$ \\
11 & $4.654 \times 10^{-8}$    & $1.063803$ \\
\bottomrule
\end{tabular}
\end{table}

\subsection{Decay constant and misalignment angle}

The order-5 decay constant and misalignment angle are fixed by
structural integers and the predicted electroweak / Planck scales:
\begin{align}
f_5 &= \frac{\text{rank}(h - \text{rank})}{\dim} \cdot
       \frac{\sqrt{v \cdot m_{\rm Planck}}}{\alpha_\tau} \notag \\
    &= \frac{22}{31} \cdot \frac{\sqrt{v\,m_{\rm Planck}}}{\alpha_\tau} \notag \\
    &= 1.082 \times 10^{13}~\text{GeV},
\label{eq:f5}\\
\theta_5 &= \arctan \varphi \approx 1.0172~\text{rad} \approx 58.28\degree.
\label{eq:theta5}
\end{align}
The ratio $\text{rank}(h-\text{rank})/\dim = 8 \cdot 22 / 248 = 22/31$
is the off-diagonal Coxeter coupling prefactor. The
$\sqrt{v \cdot m_{\rm Planck}}$ factor is the geometric mean of the
electroweak and Planck scales (the ``substrate see-saw amplitude'').
The misalignment angle $\theta_5 = \arctan\varphi$ uses the icosahedral
golden ratio.

\subsection{Rank-locked sub-harmonic correction}

A structural pattern of rank-locked-subspace corrections operates across
the gateway-mixing tower (Table~\ref{tab:subharmonic}). The order-5
mode receives the rank-locked correction
\begin{multline}
m_5^{\rm corrected} = 2 h_{\rm eff} \alpha_\tau^5 \cdot
(1 - \text{rank}\,\alpha_\tau) \\
= 2 h_{\rm eff} \alpha_\tau^5 \cdot \frac{270}{278}
= 602.645~\text{neV},
\label{eq:m5corr}
\end{multline}
a $-2.878\%$ shift from the bare tower value
$m_5^{\rm bare} = 620.501$~neV. The correction follows the same
pattern of small-coefficient Coxeter sub-harmonic shifts that operate
on $\tau$, $\mu$, $v$, and $\nu_3$ (Table~\ref{tab:subharmonic}).

\begin{table}[h]
\centering
\caption{Rank-locked sub-harmonic correction pattern across the
gateway-mixing tower. Each correction takes the form
$(1 \pm c_n \alpha_\tau^k)$ where the coefficient $c_n$ and power $k$
are fixed by the mode's structural assignment.}
\label{tab:subharmonic}
\begin{tabular}{llll}
\toprule
Mode & $n$ & Correction & Coefficient \\
\midrule
$\tau$ & 1 & $(1 - \alpha_\tau)$ & 1 \\
$\mu$ & 2 & $(1 - 2\alpha_\tau^2)$ & 2 \\
$v$ & 2 & $(1 - \text{rank}\,\alpha_\tau^2)$ & 8 \\
$\nu_3$ & 3 & $(1 + 5\varphi^2 \alpha_\tau)$ & $5\varphi^2 \approx 13.09$ \\
$m_5$ & 5 & $(1 - \text{rank}\,\alpha_\tau)$ & 8 \\
\bottomrule
\end{tabular}
\end{table}

\subsection{Term-by-term decomposition: dark-matter tower modes}
\label{sub:DM-terms}

The dark-matter tower mass formula
$m_n = 2 h_{\rm eff} \alpha_\tau^n$ (Eq.~\eqref{eq:DMtower}) admits
the same four-component decomposition. Each order $n$ in the tower
is a sub-gateway recursive mode (analogous to the leptons but at
deeper recursive depth $\alpha_\tau^n$ rather than $h^{k/15}$).

\paragraph{Generic tower mode $(n \geq 3)$.}
\begin{align*}
m_n \;=\;& \underbrace{2}_{\text{prefactor}} \cdot h_{\rm eff} \\
        &\cdot \underbrace{\alpha_\tau^n}_{\text{primary harmonic at depth }n} \\
        &\cdot \underbrace{(1 + \mathcal{C}_n)}_{\text{order-specific correction}}
\end{align*}
\begin{itemize}
\item \emph{Icosahedral spin-doublet prefactor:} the constant $2$ is
the icosahedral spin-doublet multiplicity, shared across all tower
modes (including $\nu_3$). Structurally, it encodes the binary
icosahedral group's two-element kernel under the McKay correspondence.
\item \emph{Primary harmonic:} $\alpha_\tau^n$ where $\alpha_\tau =
1/(\dim+h) = 1/278$ is the gateway-mixing parameter. Successive orders
$n$ scale by $\alpha_\tau$, generating a discrete geometric tower.
\item \emph{Sector modifier:} sub-gateway recursive ($\alpha_\tau$
chain), distinct from the visible-sector $\varphi$-chain (leptons)
and $\sqrt{2}$-chain (quarks).
\item \emph{Order-specific correction $\mathcal{C}_n$:} sub-harmonic
shift specific to order $n$. For $n = 3$ ($\nu_3$),
$\mathcal{C}_3 = 5\varphi^2 \alpha_\tau$ (icosahedral correction). For
$n = 5$ (dominant DM), $\mathcal{C}_5 = -\text{rank}\,\alpha_\tau$
(rank-locked correction; see below). For $n = 4, 6, 7, \ldots, 11$,
the corrections are sub-leading and do not change the tower's
geometric-series structure.
\end{itemize}

\paragraph{Dominant DM mode $(n = 5)$.}
\begin{align*}
m_5 \;=\;& 2 \cdot h_{\rm eff} \cdot \underbrace{\alpha_\tau^5}_{\text{primary harmonic}} \\
        &\cdot \underbrace{\bigl(1 - {\rm rank}\, \alpha_\tau\bigr)}_{\text{rank-locked sub-harmonic}}
\end{align*}
\begin{itemize}
\item \emph{Primary harmonic:} $\alpha_\tau^5 = (1/278)^5 \approx 6.0
\times 10^{-13}$, the fifth recursive depth. Combined with $2 h_{\rm
eff}$, this gives the bare value $620.5012$~neV.
\item \emph{Sub-harmonic correction:} $(1 - {\rm rank}\,\alpha_\tau) =
(1 - 8/278) = 270/278 \approx 0.9712$, the rank-locked-subspace
projection. The factor ${\rm rank} = 8$ here is the same rank that
appears at order $\alpha_\tau^2$ for the Higgs vev correction
$\alpha_v$, but here at first order in $\alpha_\tau$ rather than
second, reflecting the deeper recursive depth $n = 5$.
\item \emph{Sector modifier:} sub-gateway $\alpha_\tau$-chain.
\item \emph{Dark-sector leakage:} not separately applicable; the
sub-gateway recursive structure is itself the substrate-internal
coupling channel.
\end{itemize}

\paragraph{Couplings.}
The kinetic-mixing parameter $\epsilon_n = \alpha_\tau^n$ at each
order is identified directly with the mixing parameter in the dark-photon
literature; for $n = 5$, $\epsilon_5 = 6.022 \times 10^{-13}$. The
axion-photon coupling $g_{a\gamma\gamma} = \epsilon_5 \alpha_{\rm em}/v$
follows from the substrate cross-term construction
(developed in the companion paper~\cite{UnreasonableEffectiveness})
extended to the dark-sector mode.

\paragraph{Structural pattern across the gateway-mixing tower.}
The rank-locked sub-harmonic corrections (Table~\ref{tab:subharmonic})
follow a discrete, mode-specific pattern with coefficients
$\{1, 2, {\rm rank}, 5\varphi^2, {\rm rank}\}$ for the modes
$\{\tau, \mu, v, \nu_3, m_5\}$. The repetition of ${\rm rank} = 8$ in
both $\alpha_v$ (vev, order $\alpha_\tau^2$) and $m_5$ (DM, order
$\alpha_\tau^1$) reflects the rank-locked-subspace projection
shared between the cycle-2 boson sector and the deepest
sub-gateway DM mode.

\subsection{Cosmological abundance result}

Combining the corrected $m_5$ from Eq.~\eqref{eq:m5corr}, the
structurally-derived $f_5$ from Eq.~\eqref{eq:f5}, the icosahedral
misalignment angle $\theta_5$ from Eq.~\eqref{eq:theta5}, and the
geometric-series tower-sum factor $\zeta_{\rm tower} \approx 1.0638$,
the standard misalignment-style abundance law~\cite{KimCarosi2010},
translated into the
substrate variables used here, gives the consistency estimate
\begin{equation}
\boxed{\Omega_{\rm DM} h^2 \;=\; 0.119982,}
\label{eq:OmegaDMresult}
\end{equation}
consistent with the Planck~2018 measured value
$\Omega_{\rm DM} h^2 = 0.1200 \pm 0.001$~\cite{Planck2018} at
\SI{0.015}{\percent}, without introducing any additional fitted
cosmological parameter. The inputs $\{f_5, \theta_5, m_5, \zeta_{\rm tower}\}$
are all derived from $E_8 \times E_8$ structure plus the same single
electron-mass anchor that produces the Standard-Model spectrum, so this
match represents a structural consistency check rather than an
independent first-principles cosmological derivation.

We note that the imported misalignment-style abundance formula carries
$O(1)$
order-of-magnitude uncertainty in its application across mass regimes;
the precise numerical agreement at \SI{0.015}{\percent} should be
interpreted within that uncertainty as ``order-of-magnitude consistent
with Planck.'' A full substrate-native cosmological evolution
computation (Klein--Gordon evolution through radiation--matter equality
at $m_5$) is left for follow-up work; only after that step would the
abundance value qualify as a first-principles prediction in the same
sense as the mass-spectrum relations.

\section{Gravitational sector}
\label{sec:gravity}

Newton's constant from Eq.~\eqref{eq:gateway}:
\begin{equation}
G^{\rm pred} = 6.67431691 \times 10^{-11}~\text{m}^3\text{kg}^{-1}\text{s}^{-2},
\end{equation}
agreeing with CODATA~2022~\cite{CODATA2022} at \SI{2.5}{ppm}, below the
$G$-measurement scatter that still characterizes direct laboratory
determinations~\cite{GReview2017}. Planck mass:
$m_{\rm Planck}^{\rm pred} = 1.22088858 \times 10^{19}$~GeV (\SI{1.2}{ppm}).

\section{Two-anchor diagnostic consistency check}
\label{sec:anchor}

The framework's predictive path computes $h_{\rm eff} = v_C / R_v$
from the structural Higgs-vev identity (Sec.~\ref{sec:intro}); no PDG
mass enters. As a diagnostic check on the framework's internal
arithmetic, we additionally evaluate two fermion-anchored modes
(anchoring $h_{\rm eff}$ on the CODATA electron mass or on the CODATA
muon mass) and compare their predictions across the entire 26-parameter
spectrum. The structural-vev mode is the primary prediction; the two
fermion-anchored modes are diagnostic only. Cross-checking the two
fermion-anchored modes against each other gives a single consistency
residual
\begin{equation}
\Delta_{e\mu} = \frac{F_\mu/F_e}{m_\mu/m_e} - 1 = -2.98 \times 10^{-8} = -30~\text{ppb},
\end{equation}
constant across all 26 SM observables. The residual is set by
sub-harmonic-series truncation, not structural ambiguity. For derived
quantities involving $1/m_{\rm Planck}^2$, the residual doubles: $G$
shows two-anchor variance $2|\Delta_{e\mu}| = 60$~ppb, observed.

\section{Results}
\label{sec:results}

\paragraph{Scheme and renormalization conventions.}
The structural mass predictions in this paper are matched against
PDG values in the schemes in which those values are reported, and
\emph{the schemes mix across rows}. We adopt the following explicit
conventions, row by row, in Table~\ref{tab:complete-table}.
\begin{itemize}
\item \emph{Charged leptons} ($m_e, m_\mu, m_\tau$): pole/physical
masses (the QED self-energy is infrared-finite for charged leptons
so this coincides with the direct experimental measurement).
\item \emph{Light quarks} ($m_u, m_d, m_s$): $\overline{\rm MS}$ at
$\mu = 2$~GeV.
\item \emph{Heavy quarks} ($m_c, m_b$): $\overline{\rm MS}$ at
$\mu = m_q$ (i.e.\ $\overline{\rm MS}(m_c)$ and
$\overline{\rm MS}(m_b)$).
\item \emph{Top quark} ($m_t$): the world-average direct
(kinematic) mass quoted by PDG, which is closest to the
\emph{pole-mass} convention, \emph{not} $\overline{\rm MS}(m_t)$.
The latter is $\sim 163$~GeV, well below the framework's prediction
$\sim 172.8$~GeV; the framework's match is to the
direct/pole value $172.57$~GeV.
\item \emph{Gauge bosons} ($m_W, m_Z$, $m_H$): pole masses.
\item \emph{Higgs vev} ($v$): tree-level minimum of the Higgs
potential at the $Z$-pole scale.
\item \emph{Weak mixing angle}: on-shell scheme,
$\sin^2\theta_W^{\rm OS} \equiv 1 - m_W^2/m_Z^2$.
\item \emph{Strong coupling}: $\alpha_s^{\overline{\rm MS}}(\mu = M_Z)$.
\item \emph{Fine-structure constant}: Thomson-limit
$\alpha_{\rm em}(0) = 1/137.036$.
\end{itemize}
\emph{Mixed-scheme caveat.} Because the framework's predictions
match against pole/physical masses for some particles (leptons,
top, gauge bosons, Higgs), $\overline{\rm MS}$ at fixed external
scale for others (light and heavy quarks),
the on-shell scheme for $\sin^2\theta_W$, the
$\overline{\rm MS}(M_Z)$ scheme for $\alpha_s$, and the
Thomson-limit value for $\alpha_{\rm em}$, the ppm-level
agreements reported in Table~\ref{tab:complete-table} should be read as
agreement with PDG values \emph{within each row's stated scheme}
and \emph{not} as agreement of a single uniformly-defined object
across rows. In particular, switching $m_t$ to
$\overline{\rm MS}(m_t)$, $\sin^2\theta_W$ to
$\overline{\rm MS}(M_Z)$, or $\alpha_{\rm em}$ to
$\alpha_{\rm em}(M_Z)$ would each shift the comparison by amounts
$O(1\text{--}10\%)$, far above the ppm headline. The framework's
structural sub-harmonic corrections ($\alpha_\tau, \alpha_v,
\alpha_H, \delta_H, \delta_W$, the per-quark
$\chi_q, \lambda_q$ signs, etc.) are identified with the all-orders
effective corrections in each row's scheme: they absorb what would
otherwise be separately computed perturbative loop corrections in
a Lagrangian derivation. The framework therefore predicts only the
final scheme-specified value, and the schemes are different for
different rows by physics convention rather than by choice.
A first-principles
Lagrangian derivation in which the structural corrections are
matched against explicit one- and two-loop calculations in a
\emph{single} unified scheme is left for future work; that
derivation would convert mixed-scheme ppm-level agreement into
single-scheme ppm-level agreement.

\paragraph{Falsification implication.}
Future PDG updates that move any single row's observable outside
the framework's prediction at the ppm level \emph{in that row's
own stated scheme} constitute direct falsification. Cross-row
re-expression in a unified scheme is not committed to and would
require the future Lagrangian-level derivation noted above.

Table~\ref{tab:complete-table} summarizes the full prediction set. Primary
observables are listed together with four algebraically derived
consistency checks that are not counted as independent experimental
tests.

% Auto-generated by substrate_mass_spectrum.py (rich longtable)
% Triple-anchor format with auto-flag colours.
{\scriptsize
\setlength{\LTcapwidth}{6.5in}
\begin{longtable}{lcccc}
\caption{Substrate-framework predictions for all derivable Standard Model and gravitational observables, computed at $\mathrm{mp.prec}=500$ from three independent anchor scales: CODATA electron mass, CODATA muon mass, and the directly-computed structural Higgs vev $v_C = (\dim - 2)/(1 - \delta_W)\cdot(1 + 5\alpha_\tau^3/8)$. Each cell carries the predicted value (15-digit precision) with the deviation against the observed value beneath, in the same cell. Auto-flag legend: \textcolor{OliveGreen}{$\checkmark$ green} ($|\mathrm{dev}| \leq 1\sigma_{\mathrm{obs}}$), \textcolor{orange}{$\bullet$ orange} ($1\sigma < |\mathrm{dev}| \leq 3\sigma$), \textcolor{red}{$\times$ red} ($|\mathrm{dev}| > 3\sigma$), \textcolor{blue}{$?$ blue} (estimate-pending; the prediction depends on upstream prediction work that is not yet closed, so the deviation is not yet falsifiable), \textcolor{violet}{$?$ violet} (mass predicted, but the framework cannot yet decide whether this candidate slot is physically occupied; used for half-integer-slot dark-mode candidates). Observed column carries the published experimental uncertainty ($\pm$ ppm/ppb/ppt) below each value where available. Anchor variance (max pairwise $\sim 158$~ppt between anchors across the full 42-particle catalogue, carried by the e-anchor-vs-vev-anchor pair; e-anchor-vs-$\mu$-anchor agrees at $\sim 4$~ppt) is exposed rather than hidden.}
\label{tab:complete-table} \\
\toprule
Observable & Observed & electron-anchor & muon-anchor & vev-anchor \\
\midrule
\endfirsthead
\multicolumn{5}{l}{\emph{Table~\ref{tab:complete-table} continued}} \\
\toprule
Observable & Observed & electron-anchor & muon-anchor & vev-anchor \\
\midrule
\endhead
\midrule
\multicolumn{5}{r}{\emph{(continued on next page)}} \\
\endfoot
\bottomrule
\endlastfoot
\multicolumn{5}{l}{\textit{CHARGED LEPTON MASSES}} \\
\parbox[t]{1.7in}{\raggedright electron (MeV)} & \shortstack[c]{0.5109989500000 \\ \scriptsize$\pm 300.000~\mathrm{ppt}$} & \shortstack[c]{0.5109989500000 \\ \scriptsize\textcolor{OliveGreen}{$\checkmark\,\sim 0\;(\mathrm{FP\ floor})$}} & \shortstack[c]{0.5109989500022 \\ \scriptsize\textcolor{OliveGreen}{$\checkmark\,+4.255~\mathrm{ppt}$}} & \shortstack[c]{0.5109989500806 \\ \scriptsize\textcolor{OliveGreen}{$\checkmark\,+157.642~\mathrm{ppt}$}} \\
\parbox[t]{1.7in}{\raggedright muon (MeV)} & \shortstack[c]{105.6583755000 \\ \scriptsize$\pm 22.000~\mathrm{ppb}$} & \shortstack[c]{105.6583754996 \\ \scriptsize\textcolor{OliveGreen}{$\checkmark\,-4.255~\mathrm{ppt}$}} & \shortstack[c]{105.6583755000 \\ \scriptsize\textcolor{OliveGreen}{$\checkmark\,\sim 0\;(\mathrm{FP\ floor})$}} & \shortstack[c]{105.6583755162 \\ \scriptsize\textcolor{OliveGreen}{$\checkmark\,+153.386~\mathrm{ppt}$}} \\
\parbox[t]{1.7in}{\raggedright tau (MeV)} & \shortstack[c]{1776.860000000 \\ \scriptsize$\pm 68.000~\mathrm{ppm}$} & \shortstack[c]{1776.859738841 \\ \scriptsize\textcolor{OliveGreen}{$\checkmark\,-146.978~\mathrm{ppb}$}} & \shortstack[c]{1776.859738849 \\ \scriptsize\textcolor{OliveGreen}{$\checkmark\,-146.973~\mathrm{ppb}$}} & \shortstack[c]{1776.859739121 \\ \scriptsize\textcolor{OliveGreen}{$\checkmark\,-146.820~\mathrm{ppb}$}} \\
\multicolumn{5}{l}{\textit{QUARK MASSES}} \\
\parbox[t]{1.7in}{\raggedright up (MeV)} & \shortstack[c]{2.160000000000 \\ \scriptsize$\pm 170000.000~\mathrm{ppm}$} & \shortstack[c]{2.172135018407 \\ \scriptsize\textcolor{OliveGreen}{$\checkmark\,+5618.064~\mathrm{ppm}$}} & \shortstack[c]{2.172135018416 \\ \scriptsize\textcolor{OliveGreen}{$\checkmark\,+5618.064~\mathrm{ppm}$}} & \shortstack[c]{2.172135018749 \\ \scriptsize\textcolor{OliveGreen}{$\checkmark\,+5618.064~\mathrm{ppm}$}} \\
\parbox[t]{1.7in}{\raggedright down (MeV)} & \shortstack[c]{4.670000000000 \\ \scriptsize$\pm 70000.000~\mathrm{ppm}$} & \shortstack[c]{4.666472667699 \\ \scriptsize\textcolor{OliveGreen}{$\checkmark\,-755.317~\mathrm{ppm}$}} & \shortstack[c]{4.666472667719 \\ \scriptsize\textcolor{OliveGreen}{$\checkmark\,-755.317~\mathrm{ppm}$}} & \shortstack[c]{4.666472668435 \\ \scriptsize\textcolor{OliveGreen}{$\checkmark\,-755.317~\mathrm{ppm}$}} \\
\parbox[t]{1.7in}{\raggedright strange (MeV)} & \shortstack[c]{93.40000000000 \\ \scriptsize$\pm 92000.000~\mathrm{ppm}$} & \shortstack[c]{93.25272875250 \\ \scriptsize\textcolor{OliveGreen}{$\checkmark\,-1576.780~\mathrm{ppm}$}} & \shortstack[c]{93.25272875289 \\ \scriptsize\textcolor{OliveGreen}{$\checkmark\,-1576.780~\mathrm{ppm}$}} & \shortstack[c]{93.25272876720 \\ \scriptsize\textcolor{OliveGreen}{$\checkmark\,-1576.780~\mathrm{ppm}$}} \\
\parbox[t]{1.7in}{\raggedright charm (MeV)} & \shortstack[c]{1273.000000000 \\ \scriptsize$\pm 3600.000~\mathrm{ppm}$} & \shortstack[c]{1274.297660596 \\ \scriptsize\textcolor{OliveGreen}{$\checkmark\,+1019.372~\mathrm{ppm}$}} & \shortstack[c]{1274.297660601 \\ \scriptsize\textcolor{OliveGreen}{$\checkmark\,+1019.372~\mathrm{ppm}$}} & \shortstack[c]{1274.297660796 \\ \scriptsize\textcolor{OliveGreen}{$\checkmark\,+1019.372~\mathrm{ppm}$}} \\
\parbox[t]{1.7in}{\raggedright bottom (MeV)} & \shortstack[c]{4183.000000000 \\ \scriptsize$\pm 1700.000~\mathrm{ppm}$} & \shortstack[c]{4178.924486322 \\ \scriptsize\textcolor{OliveGreen}{$\checkmark\,-974.304~\mathrm{ppm}$}} & \shortstack[c]{4178.924486340 \\ \scriptsize\textcolor{OliveGreen}{$\checkmark\,-974.304~\mathrm{ppm}$}} & \shortstack[c]{4178.924486981 \\ \scriptsize\textcolor{OliveGreen}{$\checkmark\,-974.304~\mathrm{ppm}$}} \\
\parbox[t]{1.7in}{\raggedright top (GeV)} & \shortstack[c]{172.5700000000 \\ \scriptsize$\pm 1700.000~\mathrm{ppm}$} & \shortstack[c]{172.5698832743 \\ \scriptsize\textcolor{OliveGreen}{$\checkmark\,-676.396~\mathrm{ppb}$}} & \shortstack[c]{172.5698832751 \\ \scriptsize\textcolor{OliveGreen}{$\checkmark\,-676.392~\mathrm{ppb}$}} & \shortstack[c]{172.5698833015 \\ \scriptsize\textcolor{OliveGreen}{$\checkmark\,-676.238~\mathrm{ppb}$}} \\
\multicolumn{5}{l}{\textit{NEUTRINO MASSES}} \\
\parbox[t]{1.7in}{\raggedright neutrino 1 (meV)} & \shortstack[c]{4.510000000000e-10 \\ \scriptsize$\pm 50000.000~\mathrm{ppm}$} & \shortstack[c]{4.516090890957e-10 \\ \scriptsize\textcolor{OliveGreen}{$\checkmark\,+1350.530~\mathrm{ppm}$}} & \shortstack[c]{4.516090890976e-10 \\ \scriptsize\textcolor{OliveGreen}{$\checkmark\,+1350.530~\mathrm{ppm}$}} & \shortstack[c]{4.516090891669e-10 \\ \scriptsize\textcolor{OliveGreen}{$\checkmark\,+1350.530~\mathrm{ppm}$}} \\
\parbox[t]{1.7in}{\raggedright neutrino 2 (meV)} & \shortstack[c]{8.613900000000e-9 \\ \scriptsize$\pm 50000.000~\mathrm{ppm}$} & \shortstack[c]{8.625379528706e-9 \\ \scriptsize\textcolor{OliveGreen}{$\checkmark\,+1332.675~\mathrm{ppm}$}} & \shortstack[c]{8.625379528743e-9 \\ \scriptsize\textcolor{OliveGreen}{$\checkmark\,+1332.675~\mathrm{ppm}$}} & \shortstack[c]{8.625379530066e-9 \\ \scriptsize\textcolor{OliveGreen}{$\checkmark\,+1332.675~\mathrm{ppm}$}} \\
\parbox[t]{1.7in}{\raggedright neutrino 3 (meV)} & \shortstack[c]{5.014980000000e-8 \\ \scriptsize$\pm 50000.000~\mathrm{ppm}$} & \shortstack[c]{5.021285707439e-8 \\ \scriptsize\textcolor{OliveGreen}{$\checkmark\,+1257.374~\mathrm{ppm}$}} & \shortstack[c]{5.021285707461e-8 \\ \scriptsize\textcolor{OliveGreen}{$\checkmark\,+1257.374~\mathrm{ppm}$}} & \shortstack[c]{5.021285708231e-8 \\ \scriptsize\textcolor{OliveGreen}{$\checkmark\,+1257.375~\mathrm{ppm}$}} \\
\multicolumn{5}{l}{\textit{ELECTROWEAK SECTOR}} \\
\parbox[t]{1.7in}{\raggedright Higgs vev (GeV)} & \shortstack[c]{246.2196500000 \\ \scriptsize$\pm 24.000~\mathrm{ppb}$} & \shortstack[c]{246.2196499809 \\ \scriptsize\textcolor{OliveGreen}{$\checkmark\,-77.617~\mathrm{ppt}$}} & \shortstack[c]{246.2196499819 \\ \scriptsize\textcolor{OliveGreen}{$\checkmark\,-73.361~\mathrm{ppt}$}} & \shortstack[c]{246.2196500197 \\ \scriptsize\textcolor{OliveGreen}{$\checkmark\,+80.025~\mathrm{ppt}$}} \\
\parbox[t]{1.7in}{\raggedright Higgs boson (GeV)} & \shortstack[c]{125.2000000000 \\ \scriptsize$\pm 880.000~\mathrm{ppm}$} & \shortstack[c]{125.1997813344 \\ \scriptsize\textcolor{OliveGreen}{$\checkmark\,-1.747~\mathrm{ppm}$}} & \shortstack[c]{125.1997813350 \\ \scriptsize\textcolor{OliveGreen}{$\checkmark\,-1.747~\mathrm{ppm}$}} & \shortstack[c]{125.1997813542 \\ \scriptsize\textcolor{OliveGreen}{$\checkmark\,-1.746~\mathrm{ppm}$}} \\
\parbox[t]{1.7in}{\raggedright W boson (GeV)} & \shortstack[c]{80.36920000000 \\ \scriptsize$\pm 160.000~\mathrm{ppm}$} & \shortstack[c]{80.36920037081 \\ \scriptsize\textcolor{OliveGreen}{$\checkmark\,+4.614~\mathrm{ppb}$}} & \shortstack[c]{80.36920037115 \\ \scriptsize\textcolor{OliveGreen}{$\checkmark\,+4.618~\mathrm{ppb}$}} & \shortstack[c]{80.36920038347 \\ \scriptsize\textcolor{OliveGreen}{$\checkmark\,+4.771~\mathrm{ppb}$}} \\
\parbox[t]{1.7in}{\raggedright Z boson (GeV)} & \shortstack[c]{91.18760000000 \\ \scriptsize$\pm 23.000~\mathrm{ppm}$} & \shortstack[c]{91.19003480815 \\ \scriptsize\textcolor{blue}{$?\,+26.701~\mathrm{ppm}$}} & \shortstack[c]{91.19003480853 \\ \scriptsize\textcolor{blue}{$?\,+26.701~\mathrm{ppm}$}} & \shortstack[c]{91.19003482252 \\ \scriptsize\textcolor{blue}{$?\,+26.701~\mathrm{ppm}$}} \\
\parbox[t]{1.7in}{\raggedright Weinberg mixing fraction} & \shortstack[c]{0.2232000000000 \\ \scriptsize$\pm 140.000~\mathrm{ppm}$} & \shortstack[c]{0.2232441517887 \\ \scriptsize\textcolor{blue}{$?\,+197.813~\mathrm{ppm}$}} & \shortstack[c]{0.2232441517887 \\ \scriptsize\textcolor{blue}{$?\,+197.813~\mathrm{ppm}$}} & \shortstack[c]{0.2232441517887 \\ \scriptsize\textcolor{blue}{$?\,+197.813~\mathrm{ppm}$}} \\
\parbox[t]{1.7in}{\raggedright SU(2) coupling} & \shortstack[c]{0.6533000000000 \\ \scriptsize$\pm 1500.000~\mathrm{ppm}$} & \shortstack[c]{0.6528252345176 \\ \scriptsize\textcolor{blue}{$?\,-726.719~\mathrm{ppm}$}} & \shortstack[c]{0.6528252345176 \\ \scriptsize\textcolor{blue}{$?\,-726.719~\mathrm{ppm}$}} & \shortstack[c]{0.6528252345176 \\ \scriptsize\textcolor{blue}{$?\,-726.719~\mathrm{ppm}$}} \\
\parbox[t]{1.7in}{\raggedright $U(1)_Y$ coupling (EW norm)} & \shortstack[c]{0.3500000000000 \\ \scriptsize$\pm 1000.000~\mathrm{ppm}$} & \shortstack[c]{0.3499811975896 \\ \scriptsize\textcolor{blue}{$?\,-53.721~\mathrm{ppm}$}} & \shortstack[c]{0.3499811975896 \\ \scriptsize\textcolor{blue}{$?\,-53.721~\mathrm{ppm}$}} & \shortstack[c]{0.3499811975896 \\ \scriptsize\textcolor{blue}{$?\,-53.721~\mathrm{ppm}$}} \\
\parbox[t]{1.7in}{\raggedright $U(1)_Y$ coupling (GUT norm)} & \shortstack[c]{0.4519000000000 \\ \scriptsize$\pm 1000.000~\mathrm{ppm}$} & \shortstack[c]{0.4518237832501 \\ \scriptsize\textcolor{blue}{$?\,-168.658~\mathrm{ppm}$}} & \shortstack[c]{0.4518237832501 \\ \scriptsize\textcolor{blue}{$?\,-168.658~\mathrm{ppm}$}} & \shortstack[c]{0.4518237832501 \\ \scriptsize\textcolor{blue}{$?\,-168.658~\mathrm{ppm}$}} \\
\parbox[t]{1.7in}{\raggedright SU(2) fine-structure} & \shortstack[c]{0.03394000000000 \\ \scriptsize$\pm 1500.000~\mathrm{ppm}$} & \shortstack[c]{0.03391438943683 \\ \scriptsize\textcolor{blue}{$?\,-754.583~\mathrm{ppm}$}} & \shortstack[c]{0.03391438943683 \\ \scriptsize\textcolor{blue}{$?\,-754.583~\mathrm{ppm}$}} & \shortstack[c]{0.03391438943683 \\ \scriptsize\textcolor{blue}{$?\,-754.583~\mathrm{ppm}$}} \\
\parbox[t]{1.7in}{\raggedright $U(1)_Y$ fine-structure (EW)} & \shortstack[c]{0.009745000000000 \\ \scriptsize$\pm 2000.000~\mathrm{ppm}$} & \shortstack[c]{0.009747192918717 \\ \scriptsize\textcolor{blue}{$?\,+225.030~\mathrm{ppm}$}} & \shortstack[c]{0.009747192918717 \\ \scriptsize\textcolor{blue}{$?\,+225.030~\mathrm{ppm}$}} & \shortstack[c]{0.009747192918717 \\ \scriptsize\textcolor{blue}{$?\,+225.030~\mathrm{ppm}$}} \\
\parbox[t]{1.7in}{\raggedright $U(1)_Y$ fine-structure (GUT)} & \shortstack[c]{0.01624100000000 \\ \scriptsize$\pm 2000.000~\mathrm{ppm}$} & \shortstack[c]{0.01624532153119 \\ \scriptsize\textcolor{blue}{$?\,+266.088~\mathrm{ppm}$}} & \shortstack[c]{0.01624532153119 \\ \scriptsize\textcolor{blue}{$?\,+266.088~\mathrm{ppm}$}} & \shortstack[c]{0.01624532153119 \\ \scriptsize\textcolor{blue}{$?\,+266.088~\mathrm{ppm}$}} \\
\parbox[t]{1.7in}{\raggedright EM fine-structure (tree)} & \shortstack[c]{0.007575408000000 \\ \scriptsize$\pm 1.000~\mathrm{ppb}$} & \shortstack[c]{0.007571189103257 \\ \scriptsize\textcolor{blue}{$?\,-556.920~\mathrm{ppm}$}} & \shortstack[c]{0.007571189103257 \\ \scriptsize\textcolor{blue}{$?\,-556.920~\mathrm{ppm}$}} & \shortstack[c]{0.007571189103257 \\ \scriptsize\textcolor{blue}{$?\,-556.920~\mathrm{ppm}$}} \\
\multicolumn{5}{l}{\textit{LEPTON YUKAWA COUPLINGS}} \\
\parbox[t]{1.7in}{\raggedright electron Yukawa} & \shortstack[c]{2.935028319017e-6 \\ \scriptsize$\pm 300.000~\mathrm{ppt}$} & \shortstack[c]{2.935028319245e-6 \\ \scriptsize\textcolor{OliveGreen}{$\checkmark\,+77.617~\mathrm{ppt}$}} & \shortstack[c]{2.935028319245e-6 \\ \scriptsize\textcolor{OliveGreen}{$\checkmark\,+77.617~\mathrm{ppt}$}} & \shortstack[c]{2.935028319245e-6 \\ \scriptsize\textcolor{OliveGreen}{$\checkmark\,+77.617~\mathrm{ppt}$}} \\
\parbox[t]{1.7in}{\raggedright muon Yukawa} & \shortstack[c]{0.0006068707660433 \\ \scriptsize$\pm 22.000~\mathrm{ppb}$} & \shortstack[c]{0.0006068707660878 \\ \scriptsize\textcolor{OliveGreen}{$\checkmark\,+73.361~\mathrm{ppt}$}} & \shortstack[c]{0.0006068707660878 \\ \scriptsize\textcolor{OliveGreen}{$\checkmark\,+73.361~\mathrm{ppt}$}} & \shortstack[c]{0.0006068707660878 \\ \scriptsize\textcolor{OliveGreen}{$\checkmark\,+73.361~\mathrm{ppt}$}} \\
\parbox[t]{1.7in}{\raggedright tau Yukawa} & \shortstack[c]{0.01020576347354 \\ \scriptsize$\pm 68.000~\mathrm{ppm}$} & \shortstack[c]{0.01020576197431 \\ \scriptsize\textcolor{OliveGreen}{$\checkmark\,-146.900~\mathrm{ppb}$}} & \shortstack[c]{0.01020576197431 \\ \scriptsize\textcolor{OliveGreen}{$\checkmark\,-146.900~\mathrm{ppb}$}} & \shortstack[c]{0.01020576197431 \\ \scriptsize\textcolor{OliveGreen}{$\checkmark\,-146.900~\mathrm{ppb}$}} \\
\multicolumn{5}{l}{\textit{QUARK YUKAWA COUPLINGS}} \\
\parbox[t]{1.7in}{\raggedright up Yukawa} & \shortstack[c]{1.240640742819e-5 \\ \scriptsize$\pm 170000.000~\mathrm{ppm}$} & \shortstack[c]{1.247610742106e-5 \\ \scriptsize\textcolor{OliveGreen}{$\checkmark\,+5618.064~\mathrm{ppm}$}} & \shortstack[c]{1.247610742106e-5 \\ \scriptsize\textcolor{OliveGreen}{$\checkmark\,+5618.064~\mathrm{ppm}$}} & \shortstack[c]{1.247610742106e-5 \\ \scriptsize\textcolor{OliveGreen}{$\checkmark\,+5618.064~\mathrm{ppm}$}} \\
\parbox[t]{1.7in}{\raggedright down Yukawa} & \shortstack[c]{2.682311235631e-5 \\ \scriptsize$\pm 70000.000~\mathrm{ppm}$} & \shortstack[c]{2.680285239466e-5 \\ \scriptsize\textcolor{OliveGreen}{$\checkmark\,-755.317~\mathrm{ppm}$}} & \shortstack[c]{2.680285239466e-5 \\ \scriptsize\textcolor{OliveGreen}{$\checkmark\,-755.317~\mathrm{ppm}$}} & \shortstack[c]{2.680285239466e-5 \\ \scriptsize\textcolor{OliveGreen}{$\checkmark\,-755.317~\mathrm{ppm}$}} \\
\parbox[t]{1.7in}{\raggedright strange Yukawa} & \shortstack[c]{0.0005364622471263 \\ \scriptsize$\pm 92000.000~\mathrm{ppm}$} & \shortstack[c]{0.0005356163642517 \\ \scriptsize\textcolor{OliveGreen}{$\checkmark\,-1576.780~\mathrm{ppm}$}} & \shortstack[c]{0.0005356163642517 \\ \scriptsize\textcolor{OliveGreen}{$\checkmark\,-1576.780~\mathrm{ppm}$}} & \shortstack[c]{0.0005356163642517 \\ \scriptsize\textcolor{OliveGreen}{$\checkmark\,-1576.780~\mathrm{ppm}$}} \\
\parbox[t]{1.7in}{\raggedright charm Yukawa} & \shortstack[c]{0.007311739192631 \\ \scriptsize$\pm 3600.000~\mathrm{ppm}$} & \shortstack[c]{0.007319192575631 \\ \scriptsize\textcolor{OliveGreen}{$\checkmark\,+1019.372~\mathrm{ppm}$}} & \shortstack[c]{0.007319192575631 \\ \scriptsize\textcolor{OliveGreen}{$\checkmark\,+1019.372~\mathrm{ppm}$}} & \shortstack[c]{0.007319192575631 \\ \scriptsize\textcolor{OliveGreen}{$\checkmark\,+1019.372~\mathrm{ppm}$}} \\
\parbox[t]{1.7in}{\raggedright bottom Yukawa} & \shortstack[c]{0.02402592697783 \\ \scriptsize$\pm 1700.000~\mathrm{ppm}$} & \shortstack[c]{0.02400251842267 \\ \scriptsize\textcolor{OliveGreen}{$\checkmark\,-974.304~\mathrm{ppm}$}} & \shortstack[c]{0.02400251842267 \\ \scriptsize\textcolor{OliveGreen}{$\checkmark\,-974.304~\mathrm{ppm}$}} & \shortstack[c]{0.02400251842267 \\ \scriptsize\textcolor{OliveGreen}{$\checkmark\,-974.304~\mathrm{ppm}$}} \\
\parbox[t]{1.7in}{\raggedright top Yukawa} & \shortstack[c]{0.9911915416122 \\ \scriptsize$\pm 1700.000~\mathrm{ppm}$} & \shortstack[c]{0.9911908712511 \\ \scriptsize\textcolor{OliveGreen}{$\checkmark\,-676.318~\mathrm{ppb}$}} & \shortstack[c]{0.9911908712511 \\ \scriptsize\textcolor{OliveGreen}{$\checkmark\,-676.318~\mathrm{ppb}$}} & \shortstack[c]{0.9911908712511 \\ \scriptsize\textcolor{OliveGreen}{$\checkmark\,-676.318~\mathrm{ppb}$}} \\
\multicolumn{5}{l}{\textit{HIGGS SELF-COUPLING}} \\
\parbox[t]{1.7in}{\raggedright Higgs self-coupling} & \shortstack[c]{0.1292805654113 \\ \scriptsize$\pm 1800.000~\mathrm{ppm}$} & \shortstack[c]{0.1292801138469 \\ \scriptsize\textcolor{OliveGreen}{$\checkmark\,-3.493~\mathrm{ppm}$}} & \shortstack[c]{0.1292801138469 \\ \scriptsize\textcolor{OliveGreen}{$\checkmark\,-3.493~\mathrm{ppm}$}} & \shortstack[c]{0.1292801138469 \\ \scriptsize\textcolor{OliveGreen}{$\checkmark\,-3.493~\mathrm{ppm}$}} \\
\multicolumn{5}{l}{\textit{STRONG COUPLING}} \\
\parbox[t]{1.7in}{\raggedright strong coupling} & \shortstack[c]{0.1180000000000 \\ \scriptsize$\pm 7600.000~\mathrm{ppm}$} & \shortstack[c]{0.1180478309232 \\ \scriptsize\textcolor{blue}{$?\,+405.347~\mathrm{ppm}$}} & \shortstack[c]{0.1180478309232 \\ \scriptsize\textcolor{blue}{$?\,+405.347~\mathrm{ppm}$}} & \shortstack[c]{0.1180478309232 \\ \scriptsize\textcolor{blue}{$?\,+405.347~\mathrm{ppm}$}} \\
\multicolumn{5}{l}{\textit{PMNS NEUTRINO MIXING}} \\
\parbox[t]{1.7in}{\raggedright PMNS solar mixing fraction} & \shortstack[c]{0.3030000000000 \\ \scriptsize$\pm 50000.000~\mathrm{ppm}$} & \shortstack[c]{0.3065788699420 \\ \scriptsize\textcolor{blue}{$?\,+11811.452~\mathrm{ppm}$}} & \shortstack[c]{0.3065788699420 \\ \scriptsize\textcolor{blue}{$?\,+11811.452~\mathrm{ppm}$}} & \shortstack[c]{0.3065788699420 \\ \scriptsize\textcolor{blue}{$?\,+11811.452~\mathrm{ppm}$}} \\
\parbox[t]{1.7in}{\raggedright PMNS atmospheric mixing fraction} & \shortstack[c]{0.5720000000000 \\ \scriptsize$\pm 30000.000~\mathrm{ppm}$} & \shortstack[c]{0.5714285714286 \\ \scriptsize\textcolor{blue}{$?\,-999.001~\mathrm{ppm}$}} & \shortstack[c]{0.5714285714286 \\ \scriptsize\textcolor{blue}{$?\,-999.001~\mathrm{ppm}$}} & \shortstack[c]{0.5714285714286 \\ \scriptsize\textcolor{blue}{$?\,-999.001~\mathrm{ppm}$}} \\
\parbox[t]{1.7in}{\raggedright PMNS reactor mixing fraction} & \shortstack[c]{0.02203000000000 \\ \scriptsize$\pm 30000.000~\mathrm{ppm}$} & \shortstack[c]{0.02197802197802 \\ \scriptsize\textcolor{blue}{$?\,-2359.420~\mathrm{ppm}$}} & \shortstack[c]{0.02197802197802 \\ \scriptsize\textcolor{blue}{$?\,-2359.420~\mathrm{ppm}$}} & \shortstack[c]{0.02197802197802 \\ \scriptsize\textcolor{blue}{$?\,-2359.420~\mathrm{ppm}$}} \\
\parbox[t]{1.7in}{\raggedright PMNS CP phase (rad)} & \shortstack[c]{3.438298626429 \\ \scriptsize$\pm 50000.000~\mathrm{ppm}$} & \shortstack[c]{3.455751918949 \\ \scriptsize\textcolor{blue}{$?\,+5076.142~\mathrm{ppm}$}} & \shortstack[c]{3.455751918949 \\ \scriptsize\textcolor{blue}{$?\,+5076.142~\mathrm{ppm}$}} & \shortstack[c]{3.455751918949 \\ \scriptsize\textcolor{blue}{$?\,+5076.142~\mathrm{ppm}$}} \\
\multicolumn{5}{l}{\textit{CKM QUARK MIXING}} \\
\parbox[t]{1.7in}{\raggedright Cabibbo mixing fraction} & \shortstack[c]{0.05064000000000 \\ \scriptsize$\pm 4000.000~\mathrm{ppm}$} & \shortstack[c]{0.05035971223022 \\ \scriptsize\textcolor{blue}{$?\,-5534.909~\mathrm{ppm}$}} & \shortstack[c]{0.05035971223022 \\ \scriptsize\textcolor{blue}{$?\,-5534.909~\mathrm{ppm}$}} & \shortstack[c]{0.05035971223022 \\ \scriptsize\textcolor{blue}{$?\,-5534.909~\mathrm{ppm}$}} \\
\parbox[t]{1.7in}{\raggedright CKM atmospheric mixing fraction} & \shortstack[c]{0.001682000000000 \\ \scriptsize$\pm 15000.000~\mathrm{ppm}$} & \shortstack[c]{0.001681796794019 \\ \scriptsize\textcolor{blue}{$?\,-120.812~\mathrm{ppm}$}} & \shortstack[c]{0.001681796794019 \\ \scriptsize\textcolor{blue}{$?\,-120.812~\mathrm{ppm}$}} & \shortstack[c]{0.001681796794019 \\ \scriptsize\textcolor{blue}{$?\,-120.812~\mathrm{ppm}$}} \\
\parbox[t]{1.7in}{\raggedright CKM reactor mixing fraction} & \shortstack[c]{1.460000000000e-5 \\ \scriptsize$\pm 30000.000~\mathrm{ppm}$} & \shortstack[c]{1.452605387890e-5 \\ \scriptsize\textcolor{blue}{$?\,-5064.803~\mathrm{ppm}$}} & \shortstack[c]{1.452605387890e-5 \\ \scriptsize\textcolor{blue}{$?\,-5064.803~\mathrm{ppm}$}} & \shortstack[c]{1.452605387890e-5 \\ \scriptsize\textcolor{blue}{$?\,-5064.803~\mathrm{ppm}$}} \\
\parbox[t]{1.7in}{\raggedright CKM CP phase (rad)} & \shortstack[c]{1.143190660056 \\ \scriptsize$\pm 50000.000~\mathrm{ppm}$} & \shortstack[c]{1.151917306316 \\ \scriptsize\textcolor{blue}{$?\,+7633.588~\mathrm{ppm}$}} & \shortstack[c]{1.151917306316 \\ \scriptsize\textcolor{blue}{$?\,+7633.588~\mathrm{ppm}$}} & \shortstack[c]{1.151917306316 \\ \scriptsize\textcolor{blue}{$?\,+7633.588~\mathrm{ppm}$}} \\
\multicolumn{5}{l}{\textit{QCD THETA ANGLE}} \\
\parbox[t]{1.7in}{\raggedright QCD theta angle (rad)} & 0.0 & 0.0 & 0.0 & 0.0 \\
\multicolumn{5}{l}{\textit{MASSLESS GAUGE BOSONS}} \\
\parbox[t]{1.7in}{\raggedright photon (MeV)} & 0.0 & 0.0 & 0.0 & 0.0 \\
\parbox[t]{1.7in}{\raggedright gluon 1 (MeV)} & 0.0 & 0.0 & 0.0 & 0.0 \\
\parbox[t]{1.7in}{\raggedright gluon 2 (MeV)} & 0.0 & 0.0 & 0.0 & 0.0 \\
\parbox[t]{1.7in}{\raggedright gluon 3 (MeV)} & 0.0 & 0.0 & 0.0 & 0.0 \\
\parbox[t]{1.7in}{\raggedright gluon 4 (MeV)} & 0.0 & 0.0 & 0.0 & 0.0 \\
\parbox[t]{1.7in}{\raggedright gluon 5 (MeV)} & 0.0 & 0.0 & 0.0 & 0.0 \\
\parbox[t]{1.7in}{\raggedright gluon 6 (MeV)} & 0.0 & 0.0 & 0.0 & 0.0 \\
\parbox[t]{1.7in}{\raggedright gluon 7 (MeV)} & 0.0 & 0.0 & 0.0 & 0.0 \\
\parbox[t]{1.7in}{\raggedright gluon 8 (MeV)} & 0.0 & 0.0 & 0.0 & 0.0 \\
\multicolumn{5}{l}{\textit{GRAVITATIONAL SECTOR}} \\
\parbox[t]{1.7in}{\raggedright Newton's G (m$^3$kg$^{-1}$s$^{-2}$)} & \shortstack[c]{6.674300000000e-11 \\ \scriptsize$\pm 22.000~\mathrm{ppm}$} & \shortstack[c]{6.674316910358e-11 \\ \scriptsize\textcolor{OliveGreen}{$\checkmark\,+2.534~\mathrm{ppm}$}} & \shortstack[c]{6.674316910301e-11 \\ \scriptsize\textcolor{OliveGreen}{$\checkmark\,+2.534~\mathrm{ppm}$}} & \shortstack[c]{6.674316908253e-11 \\ \scriptsize\textcolor{OliveGreen}{$\checkmark\,+2.533~\mathrm{ppm}$}} \\
\parbox[t]{1.7in}{\raggedright Planck mass (GeV)} & \shortstack[c]{1.220890000000e+22 \\ \scriptsize$\pm 11.000~\mathrm{ppm}$} & \shortstack[c]{1.220888581931e+22 \\ \scriptsize\textcolor{OliveGreen}{$\checkmark\,-1.162~\mathrm{ppm}$}} & \shortstack[c]{1.220888581936e+22 \\ \scriptsize\textcolor{OliveGreen}{$\checkmark\,-1.161~\mathrm{ppm}$}} & \shortstack[c]{1.220888582124e+22 \\ \scriptsize\textcolor{OliveGreen}{$\checkmark\,-1.161~\mathrm{ppm}$}} \\
\multicolumn{5}{l}{\textit{ANOMALOUS MAGNETIC MOMENTS}} \\
\parbox[t]{1.7in}{\raggedright electron anomalous moment} & \shortstack[c]{0.001159652181000 \\ \scriptsize$\pm 1.700~\mathrm{ppt}$} & \shortstack[c]{0.001160667181068 \\ \scriptsize\textcolor{blue}{$?\,+875.263~\mathrm{ppm}$}} & \shortstack[c]{0.001160667181068 \\ \scriptsize\textcolor{blue}{$?\,+875.263~\mathrm{ppm}$}} & \shortstack[c]{0.001160667181068 \\ \scriptsize\textcolor{blue}{$?\,+875.263~\mathrm{ppm}$}} \\
\parbox[t]{1.7in}{\raggedright muon anomalous moment} & \shortstack[c]{0.001165920620000 \\ \scriptsize$\pm 460.000~\mathrm{ppb}$} & \shortstack[c]{0.001166944225438 \\ \scriptsize\textcolor{blue}{$?\,+877.938~\mathrm{ppm}$}} & \shortstack[c]{0.001166944225438 \\ \scriptsize\textcolor{blue}{$?\,+877.938~\mathrm{ppm}$}} & \shortstack[c]{0.001166944225438 \\ \scriptsize\textcolor{blue}{$?\,+877.938~\mathrm{ppm}$}} \\
\parbox[t]{1.7in}{\raggedright tau anomalous moment} & \shortstack[c]{0.001177210000000 \\ \scriptsize$\pm 2100.000~\mathrm{ppm}$} & \shortstack[c]{0.001176416612358 \\ \scriptsize\textcolor{blue}{$?\,-673.956~\mathrm{ppm}$}} & \shortstack[c]{0.001176416612358 \\ \scriptsize\textcolor{blue}{$?\,-673.956~\mathrm{ppm}$}} & \shortstack[c]{0.001176416612358 \\ \scriptsize\textcolor{blue}{$?\,-673.956~\mathrm{ppm}$}} \\
\multicolumn{5}{l}{\textit{FERMI CONSTANT}} \\
\parbox[t]{1.7in}{\raggedright Fermi constant (GeV$^{-2}$)} & \shortstack[c]{1.166378700000e-5 \\ \scriptsize$\pm 500.000~\mathrm{ppb}$} & \shortstack[c]{1.166378707705e-5 \\ \scriptsize\textcolor{OliveGreen}{$\checkmark\,+6.606~\mathrm{ppb}$}} & \shortstack[c]{1.166378707695e-5 \\ \scriptsize\textcolor{OliveGreen}{$\checkmark\,+6.597~\mathrm{ppb}$}} & \shortstack[c]{1.166378707337e-5 \\ \scriptsize\textcolor{OliveGreen}{$\checkmark\,+6.291~\mathrm{ppb}$}} \\
\multicolumn{5}{l}{\textit{LEPTON LIFETIMES}} \\
\parbox[t]{1.7in}{\raggedright muon lifetime (s)} & \shortstack[c]{2.196981100000e-6 \\ \scriptsize$\pm 1.000~\mathrm{ppm}$} & \shortstack[c]{2.197000758647e-6 \\ \scriptsize\textcolor{blue}{$?\,+8.948~\mathrm{ppm}$}} & \shortstack[c]{2.197000758638e-6 \\ \scriptsize\textcolor{blue}{$?\,+8.948~\mathrm{ppm}$}} & \shortstack[c]{2.197000758301e-6 \\ \scriptsize\textcolor{blue}{$?\,+8.948~\mathrm{ppm}$}} \\
\parbox[t]{1.7in}{\raggedright tau lifetime (s)} & \shortstack[c]{2.903000000000e-13 \\ \scriptsize$\pm 400.000~\mathrm{ppm}$} & \shortstack[c]{2.975066653724e-13 \\ \scriptsize\textcolor{blue}{$?\,+24824.889~\mathrm{ppm}$}} & \shortstack[c]{2.975066653712e-13 \\ \scriptsize\textcolor{blue}{$?\,+24824.889~\mathrm{ppm}$}} & \shortstack[c]{2.975066653255e-13 \\ \scriptsize\textcolor{blue}{$?\,+24824.889~\mathrm{ppm}$}} \\
\multicolumn{5}{l}{\textit{W AND Z DECAY WIDTHS}} \\
\parbox[t]{1.7in}{\raggedright W total decay width (GeV)} & \shortstack[c]{2.085000000000 \\ \scriptsize$\pm 2000.000~\mathrm{ppm}$} & \shortstack[c]{2.097035187050 \\ \scriptsize\textcolor{blue}{$?\,+5772.272~\mathrm{ppm}$}} & \shortstack[c]{2.097035187059 \\ \scriptsize\textcolor{blue}{$?\,+5772.272~\mathrm{ppm}$}} & \shortstack[c]{2.097035187381 \\ \scriptsize\textcolor{blue}{$?\,+5772.272~\mathrm{ppm}$}} \\
\parbox[t]{1.7in}{\raggedright Z total decay width (GeV)} & \shortstack[c]{2.495500000000 \\ \scriptsize$\pm 90.000~\mathrm{ppm}$} & \shortstack[c]{2.508962908085 \\ \scriptsize\textcolor{blue}{$?\,+5394.874~\mathrm{ppm}$}} & \shortstack[c]{2.508962908096 \\ \scriptsize\textcolor{blue}{$?\,+5394.874~\mathrm{ppm}$}} & \shortstack[c]{2.508962908481 \\ \scriptsize\textcolor{blue}{$?\,+5394.874~\mathrm{ppm}$}} \\
\parbox[t]{1.7in}{\raggedright Z invisible width (3 neutrinos) (GeV)} & \shortstack[c]{0.4995000000000 \\ \scriptsize$\pm 3000.000~\mathrm{ppm}$} & \shortstack[c]{0.4976876518476 \\ \scriptsize\textcolor{blue}{$?\,-3628.325~\mathrm{ppm}$}} & \shortstack[c]{0.4976876518498 \\ \scriptsize\textcolor{blue}{$?\,-3628.325~\mathrm{ppm}$}} & \shortstack[c]{0.4976876519261 \\ \scriptsize\textcolor{blue}{$?\,-3628.324~\mathrm{ppm}$}} \\
\parbox[t]{1.7in}{\raggedright Z leptonic width (per lepton) (GeV)} & \shortstack[c]{0.08398000000000 \\ \scriptsize$\pm 500.000~\mathrm{ppm}$} & \shortstack[c]{0.08389802825078 \\ \scriptsize\textcolor{blue}{$?\,-976.087~\mathrm{ppm}$}} & \shortstack[c]{0.08389802825113 \\ \scriptsize\textcolor{blue}{$?\,-976.087~\mathrm{ppm}$}} & \shortstack[c]{0.08389802826400 \\ \scriptsize\textcolor{blue}{$?\,-976.086~\mathrm{ppm}$}} \\
\multicolumn{5}{l}{\textit{HIGGS DECAY WIDTHS}} \\
\parbox[t]{1.7in}{\raggedright Higgs width to bb (GeV)} & \shortstack[c]{0.002200000000000 \\ \scriptsize$\pm 140000.000~\mathrm{ppm}$} & \shortstack[c]{0.002303033296527 \\ \scriptsize\textcolor{blue}{$?\,+46833.317~\mathrm{ppm}$}} & \shortstack[c]{0.002303033296517 \\ \scriptsize\textcolor{blue}{$?\,+46833.317~\mathrm{ppm}$}} & \shortstack[c]{0.002303033296166 \\ \scriptsize\textcolor{blue}{$?\,+46833.316~\mathrm{ppm}$}} \\
\parbox[t]{1.7in}{\raggedright Higgs width to tau tau (GeV)} & \shortstack[c]{0.0002300000000000 \\ \scriptsize$\pm 100000.000~\mathrm{ppm}$} & \shortstack[c]{0.0002591191584134 \\ \scriptsize\textcolor{blue}{$?\,+126605.037~\mathrm{ppm}$}} & \shortstack[c]{0.0002591191584145 \\ \scriptsize\textcolor{blue}{$?\,+126605.037~\mathrm{ppm}$}} & \shortstack[c]{0.0002591191584543 \\ \scriptsize\textcolor{blue}{$?\,+126605.037~\mathrm{ppm}$}} \\
\parbox[t]{1.7in}{\raggedright Higgs width to cc (GeV)} & \shortstack[c]{0.0001200000000000 \\ \scriptsize$\pm 300000.000~\mathrm{ppm}$} & \shortstack[c]{0.0001140529647702 \\ \scriptsize\textcolor{blue}{$?\,-49558.627~\mathrm{ppm}$}} & \shortstack[c]{0.0001140529647697 \\ \scriptsize\textcolor{blue}{$?\,-49558.627~\mathrm{ppm}$}} & \shortstack[c]{0.0001140529647522 \\ \scriptsize\textcolor{blue}{$?\,-49558.627~\mathrm{ppm}$}} \\
\parbox[t]{1.7in}{\raggedright Higgs total fermionic width (GeV)} & \shortstack[c]{0.002510000000000 \\ \scriptsize$\pm 50000.000~\mathrm{ppm}$} & \shortstack[c]{0.002676205419710 \\ \scriptsize\textcolor{blue}{$?\,+66217.299~\mathrm{ppm}$}} & \shortstack[c]{0.002676205419701 \\ \scriptsize\textcolor{blue}{$?\,+66217.299~\mathrm{ppm}$}} & \shortstack[c]{0.002676205419372 \\ \scriptsize\textcolor{blue}{$?\,+66217.299~\mathrm{ppm}$}} \\
\multicolumn{5}{l}{\textit{COMPTON WAVELENGTHS AND ATOMIC CONSTANTS}} \\
\parbox[t]{1.7in}{\raggedright electron Compton wavelength (fm)} & \shortstack[c]{2426.310238670 \\ \scriptsize$\pm 1.500~\mathrm{ppb}$} & \shortstack[c]{2426.310237936 \\ \scriptsize\textcolor{OliveGreen}{$\checkmark\,-302.639~\mathrm{ppt}$}} & \shortstack[c]{2426.310237925 \\ \scriptsize\textcolor{OliveGreen}{$\checkmark\,-306.895~\mathrm{ppt}$}} & \shortstack[c]{2426.310237553 \\ \scriptsize\textcolor{OliveGreen}{$\checkmark\,-460.281~\mathrm{ppt}$}} \\
\parbox[t]{1.7in}{\raggedright muon Compton wavelength (fm)} & \shortstack[c]{11.73444110000 \\ \scriptsize$\pm 22.000~\mathrm{ppb}$} & \shortstack[c]{11.73444109942 \\ \scriptsize\textcolor{OliveGreen}{$\checkmark\,-49.815~\mathrm{ppt}$}} & \shortstack[c]{11.73444109937 \\ \scriptsize\textcolor{OliveGreen}{$\checkmark\,-54.070~\mathrm{ppt}$}} & \shortstack[c]{11.73444109757 \\ \scriptsize\textcolor{OliveGreen}{$\checkmark\,-207.457~\mathrm{ppt}$}} \\
\parbox[t]{1.7in}{\raggedright tau Compton wavelength (fm)} & \shortstack[c]{0.6977713400000 \\ \scriptsize$\pm 70.000~\mathrm{ppm}$} & \shortstack[c]{0.6977714429885 \\ \scriptsize\textcolor{OliveGreen}{$\checkmark\,+147.596~\mathrm{ppb}$}} & \shortstack[c]{0.6977714429855 \\ \scriptsize\textcolor{OliveGreen}{$\checkmark\,+147.592~\mathrm{ppb}$}} & \shortstack[c]{0.6977714428785 \\ \scriptsize\textcolor{OliveGreen}{$\checkmark\,+147.439~\mathrm{ppb}$}} \\
\parbox[t]{1.7in}{\raggedright Bohr radius (fm)} & \shortstack[c]{52917.72109000 \\ \scriptsize$\pm 1.500~\mathrm{ppb}$} & \shortstack[c]{52871.37492741 \\ \scriptsize\textcolor{blue}{$?\,-875.816~\mathrm{ppm}$}} & \shortstack[c]{52871.37492719 \\ \scriptsize\textcolor{blue}{$?\,-875.816~\mathrm{ppm}$}} & \shortstack[c]{52871.37491908 \\ \scriptsize\textcolor{blue}{$?\,-875.816~\mathrm{ppm}$}} \\
\parbox[t]{1.7in}{\raggedright Rydberg energy (eV)} & \shortstack[c]{13.60569312300 \\ \scriptsize$\pm 14.000~\mathrm{ppt}$} & \shortstack[c]{13.62955661536 \\ \scriptsize\textcolor{blue}{$?\,+1753.934~\mathrm{ppm}$}} & \shortstack[c]{13.62955661542 \\ \scriptsize\textcolor{blue}{$?\,+1753.934~\mathrm{ppm}$}} & \shortstack[c]{13.62955661751 \\ \scriptsize\textcolor{blue}{$?\,+1753.934~\mathrm{ppm}$}} \\
\multicolumn{5}{l}{\textit{EFFECTIVE NEUTRINO SPECIES}} \\
\parbox[t]{1.7in}{\raggedright effective neutrino species} & \shortstack[c]{3.046000000000 \\ \scriptsize$\pm 56000.000~\mathrm{ppm}$} & \shortstack[c]{3.000000000000 \\ \scriptsize\textcolor{blue}{$?\,-15101.773~\mathrm{ppm}$}} & \shortstack[c]{3.000000000000 \\ \scriptsize\textcolor{blue}{$?\,-15101.773~\mathrm{ppm}$}} & \shortstack[c]{3.000000000000 \\ \scriptsize\textcolor{blue}{$?\,-15101.773~\mathrm{ppm}$}} \\
\multicolumn{5}{l}{\textit{DARK MATTER ULTRALIGHT BAND (predicted-but-unobserved; falsification targets)}} \\
\parbox[t]{1.7in}{\raggedright DM ultralight 1 (heaviest) (eV)} & -- & \shortstack[c]{6.205011597344e-7 \\ \scriptsize\textcolor{blue}{$?\,\mathrm{predicted}$}} & \shortstack[c]{6.205011597370e-7 \\ \scriptsize\textcolor{blue}{$?\,\mathrm{predicted}$}} & \shortstack[c]{6.205011598322e-7 \\ \scriptsize\textcolor{blue}{$?\,\mathrm{predicted}$}} \\
\parbox[t]{1.7in}{\raggedright DM ultralight 2 (eV)} & -- & \shortstack[c]{2.232018560196e-9 \\ \scriptsize\textcolor{blue}{$?\,\mathrm{predicted}$}} & \shortstack[c]{2.232018560205e-9 \\ \scriptsize\textcolor{blue}{$?\,\mathrm{predicted}$}} & \shortstack[c]{2.232018560548e-9 \\ \scriptsize\textcolor{blue}{$?\,\mathrm{predicted}$}} \\
\parbox[t]{1.7in}{\raggedright DM ultralight 3 (eV)} & -- & \shortstack[c]{8.028843741711e-12 \\ \scriptsize\textcolor{blue}{$?\,\mathrm{predicted}$}} & \shortstack[c]{8.028843741745e-12 \\ \scriptsize\textcolor{blue}{$?\,\mathrm{predicted}$}} & \shortstack[c]{8.028843742977e-12 \\ \scriptsize\textcolor{blue}{$?\,\mathrm{predicted}$}} \\
\parbox[t]{1.7in}{\raggedright DM ultralight 4 (eV)} & -- & \shortstack[c]{2.888073288385e-14 \\ \scriptsize\textcolor{blue}{$?\,\mathrm{predicted}$}} & \shortstack[c]{2.888073288398e-14 \\ \scriptsize\textcolor{blue}{$?\,\mathrm{predicted}$}} & \shortstack[c]{2.888073288841e-14 \\ \scriptsize\textcolor{blue}{$?\,\mathrm{predicted}$}} \\
\parbox[t]{1.7in}{\raggedright DM ultralight 5 (eV)} & -- & \shortstack[c]{1.038875283592e-16 \\ \scriptsize\textcolor{blue}{$?\,\mathrm{predicted}$}} & \shortstack[c]{1.038875283596e-16 \\ \scriptsize\textcolor{blue}{$?\,\mathrm{predicted}$}} & \shortstack[c]{1.038875283756e-16 \\ \scriptsize\textcolor{blue}{$?\,\mathrm{predicted}$}} \\
\parbox[t]{1.7in}{\raggedright DM ultralight 6 (eV)} & -- & \shortstack[c]{3.736961451769e-19 \\ \scriptsize\textcolor{blue}{$?\,\mathrm{predicted}$}} & \shortstack[c]{3.736961451785e-19 \\ \scriptsize\textcolor{blue}{$?\,\mathrm{predicted}$}} & \shortstack[c]{3.736961452358e-19 \\ \scriptsize\textcolor{blue}{$?\,\mathrm{predicted}$}} \\
\parbox[t]{1.7in}{\raggedright DM ultralight 7 (lightest) (eV)} & -- & \shortstack[c]{1.344230738046e-21 \\ \scriptsize\textcolor{blue}{$?\,\mathrm{predicted}$}} & \shortstack[c]{1.344230738052e-21 \\ \scriptsize\textcolor{blue}{$?\,\mathrm{predicted}$}} & \shortstack[c]{1.344230738258e-21 \\ \scriptsize\textcolor{blue}{$?\,\mathrm{predicted}$}} \\
\multicolumn{5}{l}{\textit{DARK MATTER MeV--GeV BAND (predicted-but-unobserved; direct-detection targets)}} \\
\parbox[t]{1.7in}{\raggedright DM MeV--GeV mode A (MeV)} & -- & \shortstack[c]{0.8335357316677 \\ \scriptsize\textcolor{blue}{$?\,\mathrm{predicted}$}} & \shortstack[c]{0.8335357316712 \\ \scriptsize\textcolor{blue}{$?\,\mathrm{predicted}$}} & \shortstack[c]{0.8335357317991 \\ \scriptsize\textcolor{blue}{$?\,\mathrm{predicted}$}} \\
\parbox[t]{1.7in}{\raggedright DM MeV--GeV mode B (MeV)} & -- & \shortstack[c]{5.828295547031 \\ \scriptsize\textcolor{blue}{$?\,\mathrm{predicted}$}} & \shortstack[c]{5.828295547056 \\ \scriptsize\textcolor{blue}{$?\,\mathrm{predicted}$}} & \shortstack[c]{5.828295547950 \\ \scriptsize\textcolor{blue}{$?\,\mathrm{predicted}$}} \\
\parbox[t]{1.7in}{\raggedright DM MeV--GeV mode C (MeV)} & -- & \shortstack[c]{46.62636437625 \\ \scriptsize\textcolor{blue}{$?\,\mathrm{predicted}$}} & \shortstack[c]{46.62636437645 \\ \scriptsize\textcolor{blue}{$?\,\mathrm{predicted}$}} & \shortstack[c]{46.62636438360 \\ \scriptsize\textcolor{blue}{$?\,\mathrm{predicted}$}} \\
\parbox[t]{1.7in}{\raggedright DM MeV--GeV mode D (GeV)} & -- & \shortstack[c]{1.492043660040 \\ \scriptsize\textcolor{blue}{$?\,\mathrm{predicted}$}} & \shortstack[c]{1.492043660046 \\ \scriptsize\textcolor{blue}{$?\,\mathrm{predicted}$}} & \shortstack[c]{1.492043660275 \\ \scriptsize\textcolor{blue}{$?\,\mathrm{predicted}$}} \\
\parbox[t]{1.7in}{\raggedright DM MeV--GeV mode E (GeV)} & -- & \shortstack[c]{4.816169313316 \\ \scriptsize\textcolor{blue}{$?\,\mathrm{predicted}$}} & \shortstack[c]{4.816169313336 \\ \scriptsize\textcolor{blue}{$?\,\mathrm{predicted}$}} & \shortstack[c]{4.816169314075 \\ \scriptsize\textcolor{blue}{$?\,\mathrm{predicted}$}} \\
\parbox[t]{1.7in}{\raggedright DM MeV--GeV mode F (GeV)} & -- & \shortstack[c]{592.3496671013 \\ \scriptsize\textcolor{blue}{$?\,\mathrm{predicted}$}} & \shortstack[c]{592.3496671038 \\ \scriptsize\textcolor{blue}{$?\,\mathrm{predicted}$}} & \shortstack[c]{592.3496671946 \\ \scriptsize\textcolor{blue}{$?\,\mathrm{predicted}$}} \\
\multicolumn{5}{l}{\textit{HALF-SLOT CANDIDATES (spin-1/2 dark-mode candidates; slot occupancy open)}} \\
\parbox[t]{1.7in}{\raggedright DM half-slot candidate 1 (MeV)} & -- & \shortstack[c]{0.6552852638606 \\ \scriptsize\textcolor{violet}{$?\,\mathrm{slot~open}$}} & \shortstack[c]{0.6552852638634 \\ \scriptsize\textcolor{violet}{$?\,\mathrm{slot~open}$}} & \shortstack[c]{0.6552852639639 \\ \scriptsize\textcolor{violet}{$?\,\mathrm{slot~open}$}} \\
\parbox[t]{1.7in}{\raggedright DM half-slot candidate 2 (MeV)} & -- & \shortstack[c]{2.775832922367 \\ \scriptsize\textcolor{violet}{$?\,\mathrm{slot~open}$}} & \shortstack[c]{2.775832922379 \\ \scriptsize\textcolor{violet}{$?\,\mathrm{slot~open}$}} & \shortstack[c]{2.775832922805 \\ \scriptsize\textcolor{violet}{$?\,\mathrm{slot~open}$}} \\
\parbox[t]{1.7in}{\raggedright DM half-slot candidate 3 (MeV)} & -- & \shortstack[c]{7.267224937849 \\ \scriptsize\textcolor{violet}{$?\,\mathrm{slot~open}$}} & \shortstack[c]{7.267224937879 \\ \scriptsize\textcolor{violet}{$?\,\mathrm{slot~open}$}} & \shortstack[c]{7.267224938994 \\ \scriptsize\textcolor{violet}{$?\,\mathrm{slot~open}$}} \\
\parbox[t]{1.7in}{\raggedright DM half-slot candidate 4 (MeV)} & -- & \shortstack[c]{11.75861695333 \\ \scriptsize\textcolor{violet}{$?\,\mathrm{slot~open}$}} & \shortstack[c]{11.75861695338 \\ \scriptsize\textcolor{violet}{$?\,\mathrm{slot~open}$}} & \shortstack[c]{11.75861695518 \\ \scriptsize\textcolor{violet}{$?\,\mathrm{slot~open}$}} \\
\parbox[t]{1.7in}{\raggedright DM half-slot candidate 5 (MeV)} & -- & \shortstack[c]{30.78445884451 \\ \scriptsize\textcolor{violet}{$?\,\mathrm{slot~open}$}} & \shortstack[c]{30.78445884464 \\ \scriptsize\textcolor{violet}{$?\,\mathrm{slot~open}$}} & \shortstack[c]{30.78445884936 \\ \scriptsize\textcolor{violet}{$?\,\mathrm{slot~open}$}} \\
\parbox[t]{1.7in}{\raggedright DM half-slot candidate 6 (MeV)} & -- & \shortstack[c]{49.81030073569 \\ \scriptsize\textcolor{violet}{$?\,\mathrm{slot~open}$}} & \shortstack[c]{49.81030073590 \\ \scriptsize\textcolor{violet}{$?\,\mathrm{slot~open}$}} & \shortstack[c]{49.81030074354 \\ \scriptsize\textcolor{violet}{$?\,\mathrm{slot~open}$}} \\
\parbox[t]{1.7in}{\raggedright DM half-slot candidate 7 (MeV)} & -- & \shortstack[c]{130.4050603159 \\ \scriptsize\textcolor{violet}{$?\,\mathrm{slot~open}$}} & \shortstack[c]{130.4050603164 \\ \scriptsize\textcolor{violet}{$?\,\mathrm{slot~open}$}} & \shortstack[c]{130.4050603364 \\ \scriptsize\textcolor{violet}{$?\,\mathrm{slot~open}$}} \\
\parbox[t]{1.7in}{\raggedright DM half-slot candidate 8 (MeV)} & -- & \shortstack[c]{552.4047001080 \\ \scriptsize\textcolor{violet}{$?\,\mathrm{slot~open}$}} & \shortstack[c]{552.4047001104 \\ \scriptsize\textcolor{violet}{$?\,\mathrm{slot~open}$}} & \shortstack[c]{552.4047001951 \\ \scriptsize\textcolor{violet}{$?\,\mathrm{slot~open}$}} \\
\parbox[t]{1.7in}{\raggedright DM half-slot candidate 9 (MeV)} & -- & \shortstack[c]{0.6126241040672 \\ \scriptsize\textcolor{violet}{$?\,\mathrm{slot~open}$}} & \shortstack[c]{0.6126241040698 \\ \scriptsize\textcolor{violet}{$?\,\mathrm{slot~open}$}} & \shortstack[c]{0.6126241041638 \\ \scriptsize\textcolor{violet}{$?\,\mathrm{slot~open}$}} \\
\parbox[t]{1.7in}{\raggedright DM half-slot candidate 10 (MeV)} & -- & \shortstack[c]{1.732762633217 \\ \scriptsize\textcolor{violet}{$?\,\mathrm{slot~open}$}} & \shortstack[c]{1.732762633224 \\ \scriptsize\textcolor{violet}{$?\,\mathrm{slot~open}$}} & \shortstack[c]{1.732762633490 \\ \scriptsize\textcolor{violet}{$?\,\mathrm{slot~open}$}} \\
\parbox[t]{1.7in}{\raggedright DM half-slot candidate 11 (MeV)} & -- & \shortstack[c]{3.465525266434 \\ \scriptsize\textcolor{violet}{$?\,\mathrm{slot~open}$}} & \shortstack[c]{3.465525266449 \\ \scriptsize\textcolor{violet}{$?\,\mathrm{slot~open}$}} & \shortstack[c]{3.465525266980 \\ \scriptsize\textcolor{violet}{$?\,\mathrm{slot~open}$}} \\
\parbox[t]{1.7in}{\raggedright DM half-slot candidate 12 (MeV)} & -- & \shortstack[c]{9.801985665075 \\ \scriptsize\textcolor{violet}{$?\,\mathrm{slot~open}$}} & \shortstack[c]{9.801985665117 \\ \scriptsize\textcolor{violet}{$?\,\mathrm{slot~open}$}} & \shortstack[c]{9.801985666620 \\ \scriptsize\textcolor{violet}{$?\,\mathrm{slot~open}$}} \\
\parbox[t]{1.7in}{\raggedright DM half-slot candidate 13 (MeV)} & -- & \shortstack[c]{13.86210106574 \\ \scriptsize\textcolor{violet}{$?\,\mathrm{slot~open}$}} & \shortstack[c]{13.86210106580 \\ \scriptsize\textcolor{violet}{$?\,\mathrm{slot~open}$}} & \shortstack[c]{13.86210106792 \\ \scriptsize\textcolor{violet}{$?\,\mathrm{slot~open}$}} \\
\parbox[t]{1.7in}{\raggedright DM half-slot candidate 14 (MeV)} & -- & \shortstack[c]{27.72420213147 \\ \scriptsize\textcolor{violet}{$?\,\mathrm{slot~open}$}} & \shortstack[c]{27.72420213159 \\ \scriptsize\textcolor{violet}{$?\,\mathrm{slot~open}$}} & \shortstack[c]{27.72420213584 \\ \scriptsize\textcolor{violet}{$?\,\mathrm{slot~open}$}} \\
\multicolumn{5}{l}{\textit{RULED OUT BY FRAMEWORK}} \\
\parbox[t]{1.7in}{\raggedright 4th chiral generation} & \multicolumn{4}{l}{\parbox[t]{4.5in}{\raggedright\scriptsize\textcolor{red}{\textbf{ruled out:}} Framework predicts exactly three chiral fermion generations. \textbf{Falsified} if a 4th chiral lepton or quark generation is observed.}} \\
\parbox[t]{1.7in}{\raggedright Fundamental scalar beyond Higgs} & \multicolumn{4}{l}{\parbox[t]{4.5in}{\raggedright\scriptsize\textcolor{red}{\textbf{ruled out:}} Framework predicts the Higgs is the unique fundamental scalar. \textbf{Falsified} if a non-Higgs fundamental scalar is observed.}} \\
\parbox[t]{1.7in}{\raggedright Quantum-gravity scale below $\sim 10\,E_{\rm Planck}$} & \multicolumn{4}{l}{\parbox[t]{4.5in}{\raggedright\scriptsize\textcolor{red}{\textbf{ruled out:}} Framework predicts no quantum-gravity scale below $\sim 10\,E_{\rm Planck}$. \textbf{Falsified} if quantum-gravity signatures appear at sub-Planckian energies.}} \\
\parbox[t]{1.7in}{\raggedright Ultralight DM modes below $\sim 10^{-21}$~eV} & \multicolumn{4}{l}{\parbox[t]{4.5in}{\raggedright\scriptsize\textcolor{red}{\textbf{ruled out:}} Framework predicts no ultralight-DM modes below $\sim 10^{-21}$~eV (Lyman-$\alpha$ structure-formation cutoff, Rogers \& Peiris 2021). \textbf{Falsified} if direct detection finds a discrete-spectrum DM mode in this range.}} \\
\parbox[t]{1.7in}{\raggedright PQ axion (strong CP)} & \multicolumn{4}{l}{\parbox[t]{4.5in}{\raggedright\scriptsize\textcolor{red}{\textbf{ruled out:}} Framework predicts $\theta_{\rm QCD} = 0$ exactly with no PQ axion required. \textbf{Falsified} if $\theta_{\rm QCD}$ is measured nonzero.}} \\
\parbox[t]{1.7in}{\raggedright Dark photon (kinetic-mixing operator)} & \multicolumn{4}{l}{\parbox[t]{4.5in}{\raggedright\scriptsize\textcolor{red}{\textbf{ruled out:}} Framework predicts no 4D kinetic-mixing dark-photon operator. \textbf{Falsified} if dark-photon kinetic-mixing-style coupling is detected.}} \\
\end{longtable}
}


% Auto-generated by substrate_mass_spectrum.py
% Do not edit by hand; regenerate via:
%   python substrate_mass_spectrum.py --tex-output <DIR>
\begin{table*}[t]
\centering
\caption{Per-scheme within-scheme prediction table for the pole/direct scheme. Anchored on the Higgs vev; predictions for the remaining members computed as $m_X = (R_X / R_{\rm anchor}) \cdot m_{\rm anchor}$. $\Delta(\sigma_{\rm comb})$ folds target and propagated-anchor uncertainty in quadrature. $^{*}$ marks the row used as the anchor for this scheme. Note that $v_{\rm Higgs}$ is not a direct experimental observable: the value 246.21965~GeV is back-computed structurally from the measured Fermi constant $G_F$ via $v = (\sqrt{2}\,G_F)^{-1/2}$, so anchoring on the vev anchors on a structurally derived quantity rather than a raw measurement. Anchor: Higgs vev (observed = 246.2196500000 GeV, relative $\sigma$ = $24.4$ ppm).}
\label{tab:per-scheme-pole}
\begin{ruledtabular}
\begin{tabular}{lrrrrrr}
Particle & Observed & $\sigma_{\rm obs}$ (ppm) & Predicted & $\Delta$(pred$-$obs) & $\Delta$ (ppm) & $\Delta$ ($\sigma_{\rm comb}$) \\
\hline
electron & 0.5109989500 MeV & +0.00 & 0.5109989500 MeV & 0.0000000000 MeV & +0.0001 & +0.0000 \\
muon & 105.6583755000 MeV & +0.00 & 105.6583756290 MeV & 0.0000001290 MeV & +0.0012 & +0.0001 \\
tau & 1776.8600000000 MeV & +67.53 & 1776.8597389792 MeV & -0.0002610208 MeV & -0.1469 & -0.0020 \\
top & 172.5700000000 GeV & +1738.42 & 172.5698832877 GeV & -0.0001167123 GeV & -0.6763 & -0.0004 \\
W boson & 80.3692000000 GeV & +161.75 & 80.3692003770 GeV & 0.0000003770 GeV & +0.0047 & +0.0000 \\
Z boson & 91.1876000000 GeV & +23.03 & 91.1900348152 GeV & 0.0024348152 GeV & +26.7012 & +0.7964 \\
Higgs boson & 125.2000000000 GeV & +878.59 & 125.1997813441 GeV & -0.0002186559 GeV & -1.7465 & -0.0020 \\
Higgs vev$^{*}$ & 246.2196500000 GeV & +24.37 & 246.2196500000 GeV & -- & -- & -- \\
\end{tabular}
\end{ruledtabular}
\end{table*}


% Auto-generated by substrate_mass_spectrum.py
% Do not edit by hand; regenerate via:
%   python substrate_mass_spectrum.py --tex-output <DIR>
\begin{table*}[t]
\centering
\caption{Per-scheme within-scheme prediction table for the $\overline{\rm MS}$ scheme at $\mu = 2$~GeV (light quarks, single common scale -- fully within-scheme). $^{*}$ marks the row used as the anchor for this scheme. Anchor: down (observed = 4.6700000000 MeV, relative $\sigma$ = $7.07e+04$ ppm).}
\label{tab:per-scheme-msbar2gev}
\begin{ruledtabular}
\begin{tabular}{lrrrrrr}
Particle & Observed & $\sigma_{\rm obs}$ (ppm) & Predicted & $\Delta$(pred$-$obs) & $\Delta$ (ppm) & $\Delta$ ($\sigma_{\rm comb}$) \\
\hline
up & 2.1600000000 MeV & +175925.93 & 2.1737769100 MeV & 0.0137769100 MeV & +6378.1991 & +0.0336 \\
down$^{*}$ & 4.6700000000 MeV & +70663.81 & 4.6700000000 MeV & -- & -- & -- \\
strange & 93.4000000000 MeV & +92077.09 & 93.3232174033 MeV & -0.0767825967 MeV & -822.0835 & -0.0071 \\
\end{tabular}
\end{ruledtabular}
\end{table*}


% Auto-generated by substrate_mass_spectrum.py
% Do not edit by hand; regenerate via:
%   python substrate_mass_spectrum.py --tex-output <DIR>
\begin{table*}[t]
\centering
\caption{Per-scheme within-scheme prediction table for the $\overline{\rm MS}$ scheme at $\mu = m_q$ (heavy quarks, per-particle scale -- residual cross-scale within the group). $^{*}$ marks the row used as the anchor for this scheme. Anchor: bottom (observed = 4183.0000000000 MeV, relative $\sigma$ = $1.67e+03$ ppm).}
\label{tab:per-scheme-msbarmq}
\begin{ruledtabular}
\begin{tabular}{lrrrrrr}
Particle & Observed & $\sigma_{\rm obs}$ (ppm) & Predicted & $\Delta$(pred$-$obs) & $\Delta$ (ppm) & $\Delta$ ($\sigma_{\rm comb}$) \\
\hline
charm & 1273.0000000000 MeV & +3613.51 & 1275.5404247476 MeV & 2.5404247476 MeV & +1995.6204 & +0.5010 \\
bottom$^{*}$ & 4183.0000000000 MeV & +1673.44 & 4183.0000000000 MeV & -- & -- & -- \\
\end{tabular}
\end{ruledtabular}
\end{table*}


\begin{table*}[t]
\centering
\caption{Supplementary derived neutrino-sector and cosmological observables not covered by the auto-generated complete table. Mass-squared splittings, sum of neutrino masses ($\Sigma m_\nu$), kinematic effective neutrino mass ($m_\beta$), neutrinoless double-beta effective mass ($m_{\beta\beta}$), and the dark-matter relic-density estimate $\Omega_{\rm DM} h^2$. $e$-anchor values are computed from the predicted neutrino masses; $\mu$-anchor values are obtained by linear ($m$) or quadratic ($m^2$) rescaling under the new $h_{\rm eff}$ ratio.}
\label{tab:supplementary}
\begin{ruledtabular}
\begin{tabular}{lcccc}
Observable & Observed & $e$-anchor & $\mu$-anchor & Deviation \\
\hline
$\Delta m^2_{21}$ (eV$^2$) & $7.41^{+0.21}_{-0.20} \times 10^{-5}$ & $7.4193221245 \times 10^{-5}$ & $7.4193221075 \times 10^{-5}$ & $-90$~ppm \\
$\Delta m^2_{31}$ (eV$^2$) & $(2.511 \pm 0.027) \times 10^{-3}$ & $2.5211270648 \times 10^{-3}$ & $2.5211270590 \times 10^{-3}$ & $+2440$~ppm \\
$\Sigma m_\nu$ (meV)   & $< 72$~\cite{DESI2024}   & $59.2898456922$ & $59.2898456244$ & within \\
$m_\beta$ (meV)        & $< 450$~\cite{KATRIN2024} & $8.8276108075$ & $8.8276107974$ & within \\
$m_{\beta\beta}$ (meV) & $< 28$--$122$~\cite{KamLANDZen2024} & $3.9984876767$ & $3.9984876721$ & within \\
$\Omega_{\rm DM} h^2$ & $0.1200 \pm 0.001$~\cite{Planck2018} & $0.119981783712$ & $0.119981783369$ & $-151.8$~ppm \\
\end{tabular}
\end{ruledtabular}
\end{table*}

\paragraph{Probe-scale dependence and the deviation pattern.}
The deviation column of Table~\ref{tab:complete-table} is not flat; it
exhibits a structured pattern that tracks the experimental
\emph{probe scale} $Q$ at which each PDG value is extracted, and
the associated theoretical uncertainty in that extraction. Reading
the table from cleanest to noisiest:

\begin{itemize}
\item \emph{Charged leptons} ($Q \sim m_\ell$, pure-QED extraction).
$m_e$ is the framework anchor (zero by construction); $m_\mu$
agrees at $1.1$~ppb (the precision of muonium spectroscopy);
$m_\tau$ at $0.15$~ppm (the precision of $\tau$-pair production at
$e^+e^-$ thresholds). PDG-quoted experimental uncertainties on
these masses are themselves at $\lesssim 10^{-7}, 10^{-7}, 10^{-4}$
relative respectively, and the framework's deviations track them.
\item \emph{$Z$ pole electroweak observables}
($Q \sim M_Z$, on-shell or $\overline{\rm MS}$ extractions).
$m_Z$ agrees at $30$~ppm, $m_W$ at $4$~ppm, $\sin^2\theta_W^{\rm OS}$
at $198$~ppm, $\alpha_s^{\overline{\rm MS}}(M_Z)$ at $405$~ppm.
PDG experimental relative uncertainties at this scale are
$\sim 10^{-5}, 10^{-4}, 10^{-3}, 10^{-3}$ respectively;
again the framework's deviations track the experimental error
budgets.
\item \emph{Heavy quarks} ($Q \sim m_q$,
$\overline{\rm MS}(m_q)$ or pole/direct extraction).
$m_c$ at $1019$~ppm, $m_b$ at $974$~ppm, $m_t$ at $1293$~ppm,
all within the $\sim 10^{-3}$ relative PDG uncertainty band that
itself reflects lattice-QCD systematics, scheme-conversion
ambiguities (especially the renormalon-induced
pole-vs-$\overline{\rm MS}$ split for the top), and Coulombic-bound
state corrections at the $b$ and $c$ scales.
\item \emph{Light quarks} ($Q \sim 2$~GeV,
$\overline{\rm MS}(2~\text{GeV})$ from chiral perturbation theory
plus lattice). $m_u$ at $5618$~ppm, $m_d$ at $755$~ppm, $m_s$ at
$1577$~ppm. Here the PDG-quoted uncertainties are
$\sim 10^{-2}$~relative for $m_u, m_d$ (FLAG-averaged lattice
results) and $\sim 5 \times 10^{-3}$ for $m_s$. The framework's
$\sim 0.5\%$--$0.6\%$ deviations on $m_u, m_s$ are within this
band; the $0.08\%$ on $m_d$ is well within. The
$0.6\%$ on $m_u$ in particular is comparable to FLAG's own
$\overline{\rm MS}(2~\text{GeV})$ chiral-extrapolation systematic.
\item \emph{Mixing angles and CKM/PMNS phases}
($Q$ varying with extraction; PDG uncertainties percent-level).
All within the percent-level PDG uncertainties.
\item \emph{Cosmological / gravitational sector}
($\Omega_{\rm DM} h^2$ from CMB at horizon scales; $G$ at
torsion-balance laboratory scales). Agreement at $152$~ppm and
$2.5$~ppm respectively, within experimental error.
\end{itemize}

\noindent The framework's deviations therefore correlate with the
experimental probe scale and extraction precision rather than with
any framework-internal hierarchy: cleaner extractions yield
sharper agreement, and noisier extractions yield deviations within
their own uncertainty bands. We do \emph{not} claim that this
correlation alone establishes the framework --- the residual
deviations could in principle be hiding a systematic offset of a
single multiplicative scale that happens to be inside all PDG
error bars --- but it does establish that the framework is not
\emph{tuned} to one probe scale at the expense of others. Future
PDG updates that shrink the experimental uncertainty in any row
will tighten the falsification window in that row.

\section{Predictions and falsification}
\label{sec:predictions}

The framework makes pre-registered predictions across multiple sectors,
falsifiable by:
\begin{enumerate}
\item Inverted neutrino mass ordering in future oscillation data.
\item Cosmological convergence to $\Sigma m_\nu > 70$~meV in future CMB and large-scale-structure analyses.
\item $\delta_{\rm CP}$ outside $198\degree \pm 5\degree$ in future long-baseline oscillation data.
\item $0\nu\beta\beta$ at effective Majorana mass significantly above
4~meV (assuming Majorana nature) in next-generation
neutrinoless-double-beta-decay searches.
\item Neutron EDM above the current upper bound, or a confirmed nonzero
value above $\sim 10^{-30}\,e \cdot \text{cm}$ in next-generation EDM
searches~\cite{nEDM2020}.
\item Newton's $G$ convergence inconsistent with $6.67431691 \times 10^{-11}$
at ppm precision.
\item Dark-matter detection at masses inconsistent with the predicted
tower (orders 5--11), or coupling significantly different from
$g_{a\gamma\gamma} = 1.78 \times 10^{-17}$~GeV$^{-1}$.
\item Precision EW measurements converging to $\sin^2\theta_W$
inconsistent with $2/9 + \Delta_\varphi/h$ at ppm precision.
\end{enumerate}

\section{Open items and further-exploration program}
\label{sec:open-items}

We make this section explicit and visible. The numerical demonstration
of the present paper is internally complete in the sense that, given
the starting assumptions of Table~\ref{tab:starting-assumptions}, no
fitted continuous parameter is introduced and every reported value is
a direct structural computation. The items below are open methodological
questions that affect either the empirical CY3 identification of
Sec.~\ref{sec:cy3-substrate-identification} or the substrate-physical
support of the starting assumptions themselves.

\subsection{Pipeline state and in-flight remediation}
\label{subsec:pipeline-state}

We disclose the present state of the CY3 discrimination pipeline
explicitly, separating closed results from indicative-but-unconverged
results from in-flight work. The full engineering plan to upgrade the
indicative results to converged-and-cross-checked status is the
project's \texttt{REMEDIATION\_PLAN\_OPTION\_B.md},
which specifies a seven-phase program with 111 positive and negative
acceptance tests organised into five CI tiers.

\paragraph{Closed results (do not depend on in-flight work).}
\begin{itemize}
\item \emph{Catalogue uniqueness sweep at Hodge $(3,3)$}, modulo the
gCICY non-manifest CAS residue. The Kreuzer--Skarke direct query
returns \texttt{\#NF: 0} at output limits up to $L=10{,}000$; the
$\chi=0$ self-mirror diagonal sweep $h \in \{1, \ldots, 14\}$
confirms the bin is empty for $h \le 13$ and populated from $h = 14$.
The CICY, Schoen-class, Constantin--Lukas, Borisov--Caldararu
Pfaffian, and gCICY-manifest classifications are programmatically
checked.
\item \emph{Tian--Yau Yukawa-channel structural exclusion} at the
searched bidegree range in the bicubic ambient at $h^{1,1}=2$. No
rank-3 SU(3) stable polystable monad bundle on TY/$\mathbb{Z}_3$ in
the searched range satisfies the heterotic-SM constraint suite; the
leading rank-7 monad survivor $V_{\min 2}$ is excluded by
$h^0(V_{\min 2}) = \mathbb{C}^2 \neq 0$.
\item \emph{$E_8$ Coxeter / McKay / closed-form algebraic identities}
($k_v = 1132/30$, $\delta_H = 1/59$, $\delta_W = 1/1121$,
$\alpha_\tau = 1/278$, $\omega_{\rm fix} = 123/248$, $\alpha_s = 849/7192$,
the gateway formula $F_e = (\dim-2)/\dim \cdot 5/(\dim \cdot h^M) /
(1 - 1/(4(\dim-1)))$, etc.), all closed in rational arithmetic.
\item \emph{Mass-spectrum prediction pipeline} (\texttt{substrate\_mass\_spectrum.py})
producing the residuals reported in
Tabs.~\ref{tab:complete-table}--\ref{tab:complete-table-2}.
\end{itemize}

\paragraph{Indicative, not converged.}
\begin{itemize}
\item \emph{Chain-channel Bayes factor $\ln \mathrm{BF} = -14.76$}.
The current chain-channel discriminators score the
metric-Laplacian Galerkin residual against the framework's
sub-Coxeter chain at $k_\mathrm{test} \in \{2,3,4,5\}$. Both
\texttt{p7\_11\_quark\_chain.json} and
\texttt{p7\_11\_lepton\_chain.json} report
\texttt{converged: false} with consecutive-degree relative deltas
of order 11--35\% on the Schoen-side residual. The figure is
reported as supporting evidence; the converged replacement is
Phase 2 of the remediation plan.
\item \emph{$\sigma$-channel matched-basis defense}. The current
basis-matching at $n_b = 27$ is implemented as a lex-sorted
truncation of TY's 87-dimensional invariant basis to Schoen's
27-dimensional basis. A geometric basis projection
(eigenvalue-ranked principal components of the Laplacian, or
harmonic-form ranking) is the in-flight Phase 3.3 upgrade; the
$\sigma$-channel exclusion from the model-comparison BF is
conditional on its survival under the upgrade.
\item \emph{$\sigma$ functional definition}. The production code
computes a weighted L$^2$ variance per Donaldson 2009 §3.4; a uniform
pin of the L$^1$ vs L$^2$ functional choice across code, paper, and
methodology table is in-flight Phase 3.1.
\end{itemize}

\paragraph{In-flight work (Option B remediation).}
The phases below are scheduled per the remediation plan. Each phase
has explicit acceptance criteria gating the next.
\begin{enumerate}
\item \emph{Phase 1 — pipeline integrity}: anomaly-cancellation
regression test for the Schoen monad with $W=(27,27)$;
slope-stability proof on a stated Kähler-cone region;
structure-group-of-$V$ consistency (SU(3) vs SU(4) vs SU(5));
bundle-provenance forward-confirmation; replacement or
ensemble-invariance of the LCG-random monad coefficients in
\texttt{zero\_modes.rs}; explicit Wilson-line representation
specification.
\item \emph{Phase 2 — real chain-channel Yukawa likelihood}: wire
the existing \texttt{yukawa\_pipeline.rs} (publication-grade
end-to-end Yukawa pipeline already implemented in the rust\_solver
tree but not currently plumbed into the BF combiner) into
\texttt{bayes\_factor\_multichannel.rs::from\_yukawa\_pipeline\_prediction};
parallel TY implementation; bootstrap with $\ge 100$ seeds; convergence
guarantees with \texttt{Result<\ldots, NotConvergedError>} short-circuit;
Higgs-vev fitting/normalisation; two-loop SM RG running with
threshold corrections.
\item \emph{Phase 3 — $\sigma$-channel functional pin and
matched-basis fix}: pin L$^1$ vs L$^2$; replace lex truncation with
geometric basis projection; Donaldson convention robustness sweep
across $\ge 3$ published variants.
\item \emph{Phase 4 — dark-matter $\Omega_{\rm DM} h^2$
re-derivation or relabelling}: the current pipeline uses the observed
Planck $0.12$ as a multiplicative prefactor; either re-derive the
prefactor structurally or honestly relabel as
\texttt{Omega\_DM\_h2\_consistency\_estimate} with the Planck
empirical anchor disclosed.
\item \emph{Phase 5 — structural-rule generators} for the
per-particle Coxeter chain indices $k_X$, the per-quark sign
assignments $\chi_q, \lambda_q$, and the rational closures used in
$\sin^2\theta_W$ and $\alpha_s(M_Z)$; or honest documentation as
encoded structural-rule-consistent choices in a data file.
\item \emph{Phase 6 — convergence, reproducibility, external
verification}: programmatic gCICY $(3,3)$ cycle 2.6 in a
Sage/Macaulay2-equipped container; adjacent-Hodge-bin sweep at
$(3,4), (4,3)$; rational-search baseline for closure rules;
Zenodo-archived reproducibility kit
(\texttt{cargo run --release --bin reproduce\_paper\_results});
external independent re-implementation in the AKLP PyTorch suite or
equivalent; pre-registration of falsifiable predictions.
\item \emph{Phase 7 — paper/pipeline alignment audit}: pre-submission
line-by-line audit against a 111-test acceptance matrix.
\end{enumerate}
The acceptance criteria, decision points, risk register, and test
matrix for each phase are detailed in
\texttt{REMEDIATION\_PLAN\_OPTION\_B.md}. Estimated total effort: 6--12
person-weeks; submission-readiness gate is Phase 7 sign-off.

\subsection{Open methodological items}
\label{subsec:open-method}

\begin{enumerate}
\item \emph{gCICY non-manifest residue at Hodge $(3,3)$.} The
catalogue uniqueness sweep
(Sec.~\ref{sec:cy3-substrate-identification}) closes the polytope
catalogue~\cite{KreuzerSkarke2002} and the published gCICY
classifications~\cite{AndersonApruzziGaoGrayLee2015,LarforsLukas2020,ConstantinLukasManuwal2016}
at $(h^{1,1},h^{2,1})=(3,3)$, but a future codim~$\ge (3,1)$
extension producing a $(3,3)$ entry is not pre-empted. \emph{Required
work:} structural enumeration over the gCICY-extension constructions
at parity with the polytope sweep, applying the same four-principle
filter and reporting non-Schoen survivors if any. Until this is done
the catalogue uniqueness statement at $(3,3)$ is conditional on the
gCICY non-manifest residue.

\item \emph{Donaldson convention robustness on the $\sigma$-channel
discrimination.} The $|t| = 6.92$ Schoen-vs-TY $\sigma$-channel
result reported in Sec.~\ref{sec:cy3-substrate-identification}
depends on the specific choice of Donaldson h-inversion procedure and
the Kähler-form projection convention used in the metric pipeline.
The paired-bootstrap sign-reversal defense addresses the basis-size
artifact at matched basis size; it does not address convention-choice
artifacts. \emph{Required work:} robustness sweep over $\geq 3$
published convention variants for both the h-inversion handling and
the Kähler-form projection, with a stress-test report on whether the
$|t| = 6.92$ result survives or degrades. The $\sigma$-channel
discrimination should be re-stated with an explicit convention
sensitivity envelope.

\item \emph{Schoen-uniqueness sweep at adjacent Hodge bins.} The
catalogue uniqueness statement is at the single Hodge bin $(3,3)$.
\emph{Required work:} extension of the four-cycle sweep to bins
adjacent to $(3,3)$ (in particular $(3,4)$ and $(4,3)$) at the same
fidelity, to confirm that the empty-bin result at $(3,3)$ is not
isolated to one bin under the four-principle filter.

\item \emph{Bundle provenance forward-confirmation.} An audit
performed during the development of this paper retracted some
previously asserted bundle provenance to nominally cited
references; bundles currently in use have been re-pointed to
primary sources. \emph{Required work:} a comprehensive
forward-confirmation that, for every bundle data set used in the
present paper, the citation chain terminates at a primary source
where the bundle is constructed and its stability/topological
constraints are stated.
\end{enumerate}

\subsection{Open framework items deferred to companion paper}
\label{subsec:open-framework}

The following items are starting commitments of the present paper and
are listed here for transparency. Their substrate-physical support
--- i.e.\ the argument that one would expect these and not other
substrate-physical structures --- is the subject of the companion
paper~\cite{UnreasonableEffectiveness} and of the foundational
monograph~\cite{ExistBook}, and is not load-bearing for the numerical
demonstration of this paper.

\begin{enumerate}
\item \emph{Heterotic $E_8 \times E_8$ doubling.} The Coldea
\textit{et al.}~\cite{Coldea2010} CoNb$_2$O$_6$ measurement empirically
anchors one $E_8$ factor. The doubling to $E_8 \times E_8$ is a
heterotic-string commitment~\cite{GHMR1985} adopted here; an
independent empirical anchor for the second $E_8$ factor is not
claimed.

\item \emph{McKay-icosahedral seed.} The selection of the binary
icosahedral McKay-dual of $E_8$ as the framework's harmonic seed,
and the resulting $\varphi$-chain / $\sqrt{2}$-chain decomposition,
is a structural commitment.

\item \emph{Gateway-formula structural integers.} The factor of $4$
in the sub-harmonic loss $988/987$, the gateway-depth $M=14$, and
the boundary capacity $h + m_{\max} = 59$ used in the dark-sector
coupling family admit dual-sector / two-orientation readings within
the framework but are not derived in the present paper from a
specific group-theoretic counting rule on the Schoen quotient.

\item \emph{Equivariant-cohomology identity for $\theta_{\rm QCD}=0$.}
The vanishing of the strong-CP angle is presented in
Sec.~\ref{sec:strongCP} contingent on a $\mathbb{Z}_3 \times
\mathbb{Z}_3$ quotient-projected twisted-trace cancellation. The
identity itself is adopted as input from the framework; an explicit
proof on the Schoen $X/(\mathbb{Z}_3 \times \mathbb{Z}_3)$ quotient
is part of the open program.

\item \emph{The $\omega_{\rm fix} = 123/248$ identity is algebraic,
not eigenvalue.} $\omega_{\rm fix} = (\dim - 2)/(2 \cdot \dim)$
appears in the gateway formula as an exact $E_8$-algebraic
coefficient. It is \emph{not} an eigenvalue of any operator and is
not predicted to be empirically measurable as such. We have
internally tested and retracted the alternative hypothesis under
which it would be an eigenvalue; the algebraic reading is the
correct one and is the only role $\omega_{\rm fix}$ plays in the
construction.

\item \emph{Hodge-pair self-conjugacy $h^{1,1} = h^{2,1}$.} The
substrate-physical argument that the framework selects
self-conjugate Hodge structures (and hence $(3,3)$ as the
mathematical-minimum admissible bin compatible with three
generations and Wilson-line breaking) lives in the companion
paper~\cite{UnreasonableEffectiveness}. The present paper takes the
$(3,3)$ outcome as the input to Sec.~\ref{sec:cy3-substrate-identification}.
\end{enumerate}

\subsection{What this paper does not claim}
\label{subsec:not-claimed}

For the avoidance of overreach we state plainly:

\begin{itemize}
\item This paper does not claim that the starting assumptions of
Table~\ref{tab:starting-assumptions} are forced from first principles
in the sense of being derivable from inside the framework's
macroscopic 4D regime. From a 4D observer's vantage they are evaluated
\emph{empirically} --- by the predictions they generate.

\item The 26-to-zero-free-parameter reduction reported here is
\emph{conditional} on the starting assumptions; without those
assumptions the present construction makes no statement.

\item Schoen uniqueness in full generality across the Calabi--Yau
threefold landscape is not claimed; the catalogue uniqueness
statement is at the searched Hodge bin and is conditional on the
gCICY residue (item~1 of Sec.~\ref{subsec:open-method}).

\item Internal consistency of the present construction at the
\SI{0.0298}{ppm} ($\sim 30$~ppb) two-anchor level
(Sec.~\ref{sec:anchor}) is a self-consistency statement of the
structural pipeline; it is not itself a derivation of the starting
assumptions.

\item No claim is made that the substrate-physical readings of $c$,
$\hbar$, and the macroscopic emergence of GR/QM/Newton/Maxwell from
substrate elasticity, deferred to the companion
paper~\cite{UnreasonableEffectiveness}, are required for the
numerical demonstration of this paper. They are interpretive context;
the numerical results stand or fall on their own.
\end{itemize}

\subsection{Further-exploration program}
\label{subsec:further-exploration}

The most informative empirical and theoretical extensions of the
present construction, in priority order:

\begin{enumerate}
\item Direct dark-sector searches at the predicted dominant order-5
slot $m_5 = 602.645$~neV with $g_{a\gamma\gamma} = 1.78 \times 10^{-17}$
GeV$^{-1}$. A null result at the predicted mass and coupling, with
sensitivity exceeding the prediction, falsifies the dark-sector
construction.

\item Improved precision measurements on $m_W$, $m_H$, $m_t$,
$\sin^2\theta_W$, $m_\beta$ (KATRIN), and $\Sigma m_\nu$
(cosmological) at thresholds that either tighten the present
ppm/ppb-level residuals or break them.

\item Independent reproduction of the Schoen-vs-TY $\sigma$-channel
discrimination pipeline by a third party, using independent metric
codes and convention sweeps.

\item Closure of the gCICY non-manifest residue
(Sec.~\ref{subsec:open-method} item~1) and the adjacent-Hodge-bin
sweep (item~3).

\item Heterotic-anomaly index or $\mathbb{Z}_3 \times \mathbb{Z}_3$
quotient-trace derivation of the gateway-formula structural integers
(Sec.~\ref{subsec:open-framework} item~3).

\item Explicit Schoen-side proof of the
$\mathbb{Z}_3 \times \mathbb{Z}_3$ quotient-projected twisted-trace
cancellation underlying $\theta_{\rm QCD} = 0$
(Sec.~\ref{subsec:open-framework} item~4).

\item Peer-reviewed development of the foundational arguments
relegated to the companion paper~\cite{UnreasonableEffectiveness}.
\end{enumerate}

\section{Discussion}
\label{sec:discussion}

\subsection{Parameter-count semantics}

Conditional on the starting assumptions of
Table~\ref{tab:starting-assumptions}, the construction presented here
contains \emph{no fitted continuous parameter}: the 26 continuous
parameters of the SM are computed structurally from the algebraic
data of the starting assumptions, and the residuals reported in
Sec.~\ref{sec:results} are predictions versus measurement, not fits
versus measurement. We are explicit about the scope of this
statement. (i)~Without the starting commitments
(Table~\ref{tab:starting-assumptions}) the construction has no
content; the reduction is not a model-independent claim, but a
zero-free-parameter consequence of those commitments.
(ii)~Two sectors carry residual conditionality even within the
framework: the dark-matter abundance closes through a
substrate-to-misalignment dictionary whose substrate-physical
derivation lives in the companion paper, and the strong-CP angle
closes through the quotient-projected twisted-trace cancellation
described in Sec.~\ref{sec:strongCP} and listed in
Sec.~\ref{subsec:open-framework}. (iii)~The argument that the starting
commitments themselves are minimal and substrate-physically expected
is the subject of the companion paper~\cite{UnreasonableEffectiveness}
and is not load-bearing here.

\subsection{Common geometric origin}

Notable unifying features include both CP phases as rational multiples
of $\pi$ ($\delta_{\rm CP} = \pi + \pi/10$, $\delta_{\rm CKM} = 11\pi/30$),
the dark-sector coupling family graded by $E_8$ exponents
($\delta_H = 1/59$, $\delta_W = 1/(59 \cdot 19)$), the quotient-based
conditional route to a vanishing strong CP angle, the unified
neutrino-dark-matter tower, and the strong coupling as the rational
$849/7192$ from forced $E_8$ integers.

\subsection{Comparison to prior work}

The closest empirical precedent is Coldea \textit{et al.}~\cite{Coldea2010}'s
confirmation of the $E_8$ mass spectrum in CoNb$_2$O$_6$. The closest
theoretical precedent is Koide's mass formula~\cite{Koide1981}, which
relates the three charged-lepton masses at $10^{-5}$ precision with no
derivation. The present framework matches Koide-level precision on
leptons and extends the structural construction to all 26 SM parameters
plus dark matter and gravity, while leaving the compactification-level
explanation of some closure factors for follow-up work.

The geometric strong-CP mechanism, if the quotient-projected twisted
trace vanishes, avoids invoking Peccei--Quinn~\cite{PecceiQuinn1977}
and so contrasts with the dominant axion-mediated approach.

\subsection{Open directions}

The complete enumeration of open methodological items, deferred
framework commitments, and the further-exploration program is given
in Sec.~\ref{sec:open-items}. Among the items most likely to sharpen
or break the present construction:
\begin{enumerate}
\item Direct dark-sector searches at the predicted dominant order-5
slot (Sec.~\ref{subsec:further-exploration} item~1).
\item Closure of the gCICY non-manifest residue and the
adjacent-Hodge-bin sweep
(Sec.~\ref{subsec:open-method} items~1 and~3).
\item Donaldson-convention robustness sweep on the
$\sigma$-channel discrimination
(Sec.~\ref{subsec:open-method} item~2).
\item Substrate-native cosmological evolution for the dark-matter
tower, replacing the misalignment-style abundance estimate with a
full continuum calculation.
\item Cosmological inflationary connection: $n_s = 58/60$ from dual-$E_8$
Coxeter structure.
\end{enumerate}

\subsection{Conclusion}

Conditional on the starting assumptions of
Table~\ref{tab:starting-assumptions} --- heterotic
$E_8 \times E_8$ on the Schoen Calabi--Yau quotient
$X/(\mathbb{Z}_3 \times \mathbb{Z}_3)$, the McKay-icosahedral seed,
the phenomenological $\varphi$/$\sqrt{2}$ chain split for the lepton
and quark/EW-boson sectors, the $E_8$ Coxeter spectrum, and the
substrate-tension reading of the Higgs vacuum --- the present paper computes the SM
mass spectrum, mixing parameters, gauge couplings, the strong-CP
angle, the cosmological dark-matter abundance, and Newton's $G$
\emph{directly}, with no fitted continuous parameter and no
per-particle scheme choice. The Schoen-versus-Tian--Yau
$\sigma$-channel discrimination (Sec.~\ref{sec:cy3-substrate-identification})
and the four-cycle catalogue uniqueness sweep at Hodge $(3,3)$
provide the empirical CY3 identification, with explicit hedges on
the gCICY residue and the Donaldson convention sensitivity
(Sec.~\ref{sec:open-items}). The numerical residuals against PDG and
Planck~2018 measurements are predictions, not fits. From a 4D
observer's vantage the starting assumptions cannot be \emph{proven};
this paper evaluates them by empirical consistency, and reports the
result. The companion paper~\cite{UnreasonableEffectiveness}
addresses the separate question of whether those starting assumptions
are minimal and substrate-physically expected; that argument is
independent of the present demonstration and is not relied upon here.

\appendix

\section{Coxeter step assignments}
\label{app:assignments}

\begin{table}[h]
\centering
\caption{Coxeter step assignments. The $\varphi$-chain hosts integer
$E_8$ exponent positions; the $\sqrt{2}$-chain hosts integer and
half-integer positions. Cycle 2 indicates positions beyond $h = 30$.}
\begin{tabular}{lcl}
\toprule
Particle & $k$ & Chain type / sector \\
\midrule
Electron      & 0       & $\varphi$ gateway, dark-sector \\
Muon          & 11      & $\varphi$ exponent, visible \\
Tau           & 17      & $\varphi$ exponent, visible \\
Up            & 4       & $\sqrt{2}$ int \\
Down          & 6.5     & $\sqrt{2}$ half \\
Strange       & 15      & $\sqrt{2}$ midpoint \\
Charm         & 22.5    & $\sqrt{2}$ half \\
Bottom        & 26      & $\sqrt{2}$ int \\
Top           & 36.5    & $\sqrt{2}$ half, cycle 2 \\
$W$ boson     & 34.5    & cycle 2, with $\delta_W$ \\
$Z$ boson     & via $m_W/\cos\theta_W$ & derived \\
Higgs vev     & $1132/30$ & cycle 2 \\
Higgs boson   & $k_v - 2$ & breathing mode, with $\delta_H$ \\
\bottomrule
\end{tabular}
\end{table}

\section{Complete structural correction list}
\label{app:corrections}

\begin{table*}[t]
\centering
\caption{All structural corrections, built from framework-fixed
$E_8$ invariants $\dim = 248$, $h = 30$, $M = 14$, rank $= 8$,
$m_{\max} = 29$, together with the framework-selected McKay order $5$
and golden ratio $\varphi$.}
\begin{tabular}{lll}
\toprule
Symbol & Definition & Role \\
\midrule
$\Delta_\varphi$       & $1/(h + \varphi^2)$              & enters $\sin^2\theta_W$ \\
$\Delta_2$             & $1/2^5 = 1/32$                   & 2-fold structure \\
$\alpha_\tau$          & $1/(\dim + h) = 1/278$           & tau base \\
$\alpha_\tau^{\rm full}$ & $1/(\dim + h + \varphi^2)$    & tau closed-form \\
$\alpha_\mu$           & $2 \alpha_\tau^2$                & muon \\
$\alpha_v$             & $(\text{rank}\,h/m_{\max})\alpha_\tau^2/(1+\alpha_\tau/(2\varphi))$ & Higgs vev (closure $2\varphi$) \\
$\alpha_H$             & $((\text{rank}-2)h/m_{\max})\alpha_\tau^2/(1+\alpha_\tau/h_{SU(3)})$ & Higgs boson (closure $h_{SU(3)}$) \\
$\alpha_W$             & $(2/e_2)\alpha_\tau^2/(1+\alpha_\tau/h_{SU(3)})$ & W boson sub-harmonic \\
$\alpha_t$             & $(h/h_{SU(3)})^2 \alpha_\tau^2/(1+\alpha_\tau/h_{SU(3)})$ & Top quark sub-harmonic \\
$\delta_H$             & $1/(h + m_{\max}) = 1/59$        & Higgs dark coupling \\
$\delta_W$             & $1/(59 \cdot 19) = 1/1121$       & W dark coupling \\
$\varepsilon$          & $1/(2\dim) = 1/496$              & per-step quark loss \\
$\omega_{\rm fix}$     & $1/2 - 1/\dim = 123/248$         & gateway normalization \\
$k_v$                  & $1132/30$                        & Higgs vev position \\
$f_5$                  & $\frac{\text{rank}(h-\text{rank})}{\dim} \frac{\sqrt{v M_{\rm Pl}}}{\alpha_\tau}$ & DM decay constant \\
$\theta_5$             & $\arctan \varphi$                & DM misalignment \\
\bottomrule
\end{tabular}
\end{table*}

\section{Numerical precision}
\label{app:numerics}

All numerical evaluations use mpmath at $\geq 50$-digit decimal
precision. Key exact rational values:
\begin{itemize}
\item $h_{\rm eff}^{(e)}/m_e = 124/123$;
\item $30^{14} = 4{,}782{,}969 \times 10^{14}$;
\item $\sin^2 \theta_{12}^{\rm CKM} = 14/278 = 7/139$;
\item $\delta_{\rm CKM} = 11\pi/30$;
\item $\delta_{\rm CP} = 11\pi/10$;
\item $\sin^2 \theta_W = 2/9 + 1/[h(h+\varphi^2)]$;
\item $\alpha_s(M_Z) = 849/7192$;
\item $\delta_H = 1/59$, $\delta_W = 1/1121$.
\end{itemize}

\begin{thebibliography}{99}

\bibitem{PDG2024} S.~Navas \textit{et al.} (Particle Data Group),
``Review of Particle Physics,'' Phys.\ Rev.\ D \textbf{110}, 030001 (2024).
Open-access PDG portal (free, complete review):
\href{https://pdg.lbl.gov/2024/}{pdg.lbl.gov/2024}.
DOI: 10.1103/PhysRevD.110.030001.

\bibitem{Bertone2018} G.~Bertone and T.~M.~P.~Tait,
``A new era in the search for dark matter,''
Nature \textbf{562}, 51--56 (2018).
Open-access preprint (free): \href{https://arxiv.org/abs/1810.01668}{arXiv:1810.01668}.
DOI: 10.1038/s41586-018-0542-z.

\bibitem{Coldea2010} R.~Coldea, D.~A.~Tennant, E.~M.~Wheeler, E.~Wawrzynska,
D.~Prabhakaran, M.~Telling, K.~Habicht, P.~Smeibidl, and K.~Kiefer,
``Quantum Criticality in an Ising Chain: Experimental Evidence for Emergent $E_8$ Symmetry,''
Science \textbf{327}, 177--180 (2010).
Open-access repository copy (free, same content):
\href{https://arxiv.org/abs/1103.3694}{arXiv:1103.3694}.
PubMed (free abstract): \href{https://pubmed.ncbi.nlm.nih.gov/20056884/}{PMID 20056884}.
DOI: 10.1126/science.1180085.

\bibitem{Zamolodchikov1989} A.~B.~Zamolodchikov,
``Integrals of motion and S-matrix of the (scaled) $T = T_c$ Ising model with magnetic field,''
Int.\ J.\ Mod.\ Phys.\ A \textbf{4}, 4235--4248 (1989).
Inspire-HEP record (free abstract, cross-references):
\href{https://inspirehep.net/literature/276727}{inspirehep.net/literature/276727}.
DOI: 10.1142/S0217751X8900176X.

\bibitem{McKay1980} J.~McKay,
``Graphs, singularities, and finite groups,''
Proc.\ Symp.\ Pure Math.\ \textbf{37}, 183--186 (Amer.\ Math.\ Soc., 1980).
Open scan of the original conference paper (free):
\href{https://www-users.cse.umn.edu/~reiner/ReadingSeminars/McKaySeminar2016/McKaysOriginalPaper.pdf}{McKaysOriginalPaper.pdf}.
DOI of the original: 10.1090/pspum/037/604577.

\bibitem{Humphreys1990} J.~E.~Humphreys,
\textit{Reflection Groups and Coxeter Groups},
Cambridge Studies in Advanced Mathematics \textbf{29}
(Cambridge University Press, 1990).
ISBN 978-0-521-43613-7.
Identity $|\Phi| = h \cdot \text{rank}$ for irreducible finite
Coxeter groups: \S 3.18, Eq.\ (3) (consequence of the exponents
satisfying $\sum_j (e_j + 1) = |\Phi|/2 + \text{rank}$ with
mean exponent $h/2$).
Cambridge product page (free abstract):
\href{https://www.cambridge.org/core/books/reflection-groups-and-coxeter-groups/0D40E18E1B0CC9CD32EF1B14D9E89B79}{cambridge.org/core/books/reflection-groups-and-coxeter-groups}.
DOI: 10.1017/CBO9780511623646.

\bibitem{Anderson1962} P.~W.~Anderson,
``Plasmons, Gauge Invariance, and Mass,''
Phys.\ Rev.\ \textbf{130}, 439--442 (1963).
APS abstract page (free abstract):
\href{https://journals.aps.org/pr/abstract/10.1103/PhysRev.130.439}{journals.aps.org/pr/abstract/10.1103/PhysRev.130.439}.
NASA ADS (free abstract + cross-references):
\href{https://ui.adsabs.harvard.edu/abs/1963PhRv..130..439A/abstract}{ui.adsabs.harvard.edu/abs/1963PhRv..130..439A}.
DOI: 10.1103/PhysRev.130.439.

\bibitem{Volovik2003} G.~E.~Volovik,
\textit{The Universe in a Helium Droplet},
International Series of Monographs on Physics \textbf{117}
(Oxford University Press, 2003).
Oxford Academic landing page (free metadata):
\href{https://academic.oup.com/book/11557}{academic.oup.com/book/11557}.
DOI of book: 10.1093/acprof:oso/9780199564842.001.0001.

\bibitem{Volovik2001} G.~E.~Volovik,
``Superfluid analogies of cosmological phenomena,''
Phys.\ Rep.\ \textbf{351}, 195--348 (2001).
Open-access preprint (free): \href{https://arxiv.org/abs/gr-qc/0005091}{arXiv:gr-qc/0005091}.
DOI: 10.1016/S0370-1573(00)00139-3.

\bibitem{Schoen1988} C.~Schoen,
``On fiber products of rational elliptic surfaces with section,''
Math.\ Z.\ \textbf{197}, 177--199 (1988).
EuDML page (free metadata):
\href{https://eudml.org/doc/183725}{eudml.org/doc/183725}.
DOI: 10.1007/BF01215188.

\bibitem{GHMR1985} D.~J.~Gross, J.~A.~Harvey, E.~Martinec, and R.~Rohm,
``Heterotic String,''
Phys.\ Rev.\ Lett.\ \textbf{54}, 502--505 (1985).
Inspire-HEP record (free abstract + cross-references):
\href{https://inspirehep.net/literature/15974}{inspirehep.net/literature/15974}.
NASA ADS:
\href{https://ui.adsabs.harvard.edu/abs/1985PhRvL..54..502G/abstract}{ui.adsabs.harvard.edu/abs/1985PhRvL..54..502G}.
DOI: 10.1103/PhysRevLett.54.502.

\bibitem{BraunHeOvrutPantev2005} V.~Braun, Y.-H.~He, B.~A.~Ovrut, and T.~Pantev,
``A standard model from the $E_8 \times E_8$ heterotic superstring,''
JHEP \textbf{06}, 039 (2005).
Open-access preprint (free): \href{https://arxiv.org/abs/hep-th/0502155}{arXiv:hep-th/0502155}.
DOI: 10.1088/1126-6708/2005/06/039.

\bibitem{DonagiOvrut2005} R.~Donagi, Y.-H.~He, B.~A.~Ovrut, and R.~Reinbacher,
``The spectra of heterotic standard model vacua,''
JHEP \textbf{06}, 070 (2005).
Open-access preprint (free): \href{https://arxiv.org/abs/hep-th/0411156}{arXiv:hep-th/0411156}.
DOI: 10.1088/1126-6708/2005/06/070.

\bibitem{BHOPVectorBundle2006} V.~Braun, Y.-H.~He, B.~A.~Ovrut, and T.~Pantev,
``Vector Bundle Extensions, Sheaf Cohomology, and the Heterotic Standard Model,''
Adv.\ Theor.\ Math.\ Phys.\ \textbf{10}, 525--589 (2006).
Open-access preprint (free): \href{https://arxiv.org/abs/hep-th/0505041}{arXiv:hep-th/0505041}.
\S 6 introduces the rank-4 $SU(4)$ extension construction used as
the framework's bundle commitment; Eq.~99 records the heterotic
Bianchi anomaly-cancellation identity with the positive-tension
$5$-brane wrapping $\tau_1^2$.

\bibitem{PatiSalam1974} J.~C.~Pati and A.~Salam,
``Lepton number as the fourth color,''
Phys.\ Rev.\ D \textbf{10}, 275--289 (1974);
erratum \emph{ibid.}\ \textbf{11}, 703 (1975).
APS abstract page (free abstract):
\href{https://journals.aps.org/prd/abstract/10.1103/PhysRevD.10.275}{journals.aps.org/prd/abstract/10.1103/PhysRevD.10.275}.
DOI: 10.1103/PhysRevD.10.275. Source for the substrate-physical
``lepton number as the fourth colour'' reading of the rank-4
$SU(4)$ internal-symmetry structure committed in this paper.

\bibitem{TianYau1987} G.~Tian and S.-T.~Yau,
``Three-dimensional algebraic manifolds with $c_1 = 0$ and $\chi = -6$,''
in \emph{Mathematical Aspects of String Theory} (San Diego, CA, 1986),
Adv.\ Ser.\ Math.\ Phys.\ Vol.~1, World Scientific, Singapore (1987).

\bibitem{AndersonGrayLukasPalti2012} L.~B.~Anderson, J.~Gray, A.~Lukas, and E.~Palti,
``Heterotic line bundle standard models,''
JHEP \textbf{06}, 113 (2012).
Open-access preprint (free): \href{https://arxiv.org/abs/1202.1757}{arXiv:1202.1757}.
DOI: 10.1007/JHEP06(2012)113.

\bibitem{AndersonHeLukas2008} L.~B.~Anderson, Y.-H.~He, and A.~Lukas,
``Heterotic compactification, an algorithmic approach,''
JHEP \textbf{07}, 049 (2007).
Open-access preprint (free): \href{https://arxiv.org/abs/hep-th/0702210}{arXiv:hep-th/0702210}.
DOI: 10.1088/1126-6708/2007/07/049.

\bibitem{KreuzerSkarke2002} M.~Kreuzer and H.~Skarke,
``Complete classification of reflexive polyhedra in four dimensions,''
Adv.\ Theor.\ Math.\ Phys.\ \textbf{4}, 1209--1230 (2000).
Open-access preprint (free): \href{https://arxiv.org/abs/hep-th/0002240}{arXiv:hep-th/0002240}.
Catalogue front-end: \url{http://hep.itp.tuwien.ac.at/~kreuzer/CY/}.

\bibitem{AndersonApruzziGaoGrayLee2015} L.~B.~Anderson, F.~Apruzzi, X.~Gao, J.~Gray, and S.-J.~Lee,
``A new construction of Calabi--Yau manifolds: generalized CICYs,''
Nucl.\ Phys.\ B \textbf{906}, 441--496 (2016).
Open-access preprint (free): \href{https://arxiv.org/abs/1507.03235}{arXiv:1507.03235}.
DOI: 10.1016/j.nuclphysb.2016.03.016.

\bibitem{LarforsLukas2020} M.~Larfors and R.~Schneider,
``Line bundle cohomologies on CICYs with Picard number two,''
Fortsch.\ Phys.\ \textbf{67}, 1900083 (2019).
Open-access preprint (free): \href{https://arxiv.org/abs/2003.10130}{arXiv:2003.10130}.

\bibitem{ConstantinLukasManuwal2016} A.~Constantin, A.~Lukas, and C.~Mishra,
``The family of quintic Calabi--Yau threefolds, free quotients, and their cohomology,''
Fortsch.\ Phys.\ \textbf{64}, 463--509 (2016).
Open-access preprint (free): \href{https://arxiv.org/abs/1606.04032}{arXiv:1606.04032}.

\bibitem{HuybrechtsLehn2010} D.~Huybrechts and M.~Lehn,
\emph{The Geometry of Moduli Spaces of Sheaves}, 2nd ed.,
Cambridge Mathematical Library, Cambridge University Press (2010).
DOI: 10.1017/CBO9780511711985.

\bibitem{KassRaftery1995} R.~E.~Kass and A.~E.~Raftery,
``Bayes factors,''
J. Amer. Statist. Assoc. \textbf{90}, 773--795 (1995).
DOI: 10.1080/01621459.1995.10476572.

\bibitem{Donaldson2009} S.~K.~Donaldson,
\emph{Some numerical results in complex differential geometry},
Pure Appl.\ Math.\ Q.\ \textbf{5} (2009), 571--618.
DOI: \href{https://doi.org/10.4310/PAMQ.2009.v5.n2.a2}{10.4310/PAMQ.2009.v5.n2.a2}.
arXiv:math/0512625.
See §3 eq.~(3.4) for the weighted L$^2$ Monge--Ampère residual functional
$\sigma^2 = \langle (\eta/\kappa - 1)^2 \rangle_{w}$ with $\kappa = \langle \eta \rangle_{w}$
that defines the production $\sigma$-discriminator used in this work.

\bibitem{Jeffreys1961} H.~Jeffreys,
\emph{Theory of Probability}, 3rd ed.\
(Oxford University Press, 1961).
Source of the ``decisive evidence'' threshold $\log_{10} \mathrm{BF} > 2$ on the Bayes-factor strength scale.

\bibitem{DonaldsonStability} S.~K.~Donaldson,
``Anti self-dual Yang--Mills connections over complex algebraic surfaces and stable vector bundles,''
Proc.\ London Math.\ Soc.\ \textbf{50}, 1--26 (1985).
DOI: 10.1112/plms/s3-50.1.1.

\bibitem{UhlenbeckYau} K.~Uhlenbeck and S.-T.~Yau,
``On the existence of Hermitian--Yang--Mills connections in stable vector bundles,''
Comm.\ Pure Appl.\ Math.\ \textbf{39}, S257--S293 (1986).
DOI: 10.1002/cpa.3160390714.

\bibitem{KimCarosi2010} J.~E.~Kim and G.~Carosi,
``Axions and the strong CP problem,''
Rev.\ Mod.\ Phys.\ \textbf{82}, 557--601 (2010).
Open-access preprint (free): \href{https://arxiv.org/abs/0807.3125}{arXiv:0807.3125}.
DOI: 10.1103/RevModPhys.82.557;
Erratum: Rev.\ Mod.\ Phys.\ \textbf{91}, 049902 (2019).

\bibitem{PecceiQuinn1977} R.~D.~Peccei and H.~R.~Quinn,
``CP Conservation in the Presence of Pseudoparticles,''
Phys.\ Rev.\ Lett.\ \textbf{38}, 1440--1443 (1977).
NASA ADS (free abstract + cross-references):
\href{https://ui.adsabs.harvard.edu/abs/1977PhRvL..38.1440P/abstract}{ui.adsabs.harvard.edu/abs/1977PhRvL..38.1440P}.
Inspire-HEP: \href{https://inspirehep.net/literature/120612}{inspirehep.net/literature/120612}.
DOI: 10.1103/PhysRevLett.38.1440.

\bibitem{nEDM2020} C.~Abel \textit{et al.} (nEDM Collaboration at PSI),
``Measurement of the Permanent Electric Dipole Moment of the Neutron,''
Phys.\ Rev.\ Lett.\ \textbf{124}, 081803 (2020).
Open-access preprint (free): \href{https://arxiv.org/abs/2001.11966}{arXiv:2001.11966}.
DOI: 10.1103/PhysRevLett.124.081803.
Reports $|d_n| < 1.8 \times 10^{-26}\,e\cdot\text{cm}$ at 90\% C.L.

\bibitem{CODATA2022} P.~J.~Mohr, D.~B.~Newell, B.~N.~Taylor, and E.~Tiesinga,
``CODATA recommended values of the fundamental physical constants: 2022,''
Rev.\ Mod.\ Phys.\ \textbf{97}, 025002 (2025).
Open-access preprint (free): \href{https://arxiv.org/abs/2409.03787}{arXiv:2409.03787}.
NIST free interactive constant database:
\href{https://physics.nist.gov/cuu/Constants/}{physics.nist.gov/cuu/Constants}.
DOI: 10.1103/RevModPhys.97.025002.

\bibitem{GReview2017} C.~Rothleitner and S.~Schlamminger,
``Invited Review Article: Measurements of the Newtonian constant of gravitation, $G$,''
Rev.\ Sci.\ Instrum.\ \textbf{88}, 111101 (2017).
NIST author manuscript (free):
\href{https://www.nist.gov/publications/invited-review-article-measurements-newtonian-constant-gravitation-g}{nist.gov/publications/invited-review-article-measurements-newtonian-constant-gravitation-g}.
DOI: 10.1063/1.4994619.

\bibitem{Esteban2020} I.~Esteban, M.~C.~Gonzalez-Garcia, M.~Maltoni,
T.~Schwetz, and A.~Zhou,
``The fate of hints: updated global analysis of three-flavor neutrino oscillations,''
JHEP \textbf{09}, 178 (2020).
JHEP is open access; published version free at:
doi:10.1007/JHEP09(2020)178.
Preprint: \href{https://arxiv.org/abs/2007.14792}{arXiv:2007.14792}.

\bibitem{NuFIT53} NuFIT Collaboration (I.~Esteban, M.~C.~Gonzalez-Garcia,
M.~Maltoni, T.~Schwetz, A.~Zhou),
``NuFIT 5.3: Three-neutrino fit based on data available in March 2024,''
\url{http://www.nu-fit.org/?q=node/278} (free, accessed April 2026).

\bibitem{DESI2024} A.~G.~Adame \textit{et al.} (DESI Collaboration),
``DESI 2024 VI: Cosmological Constraints from the Measurements of Baryon Acoustic Oscillations,''
JCAP \textbf{02}, 021 (2025).
JCAP is open access; published version free at:
doi:10.1088/1475-7516/2025/02/021.
Preprint: \href{https://arxiv.org/abs/2404.03002}{arXiv:2404.03002}.

\bibitem{KATRIN2024} M.~Aker \textit{et al.} (KATRIN Collaboration),
``Direct neutrino-mass measurement based on 259 days of KATRIN data,''
Science \textbf{388}, 180--185 (2025).
Open-access preprint (free): \href{https://arxiv.org/abs/2406.13516}{arXiv:2406.13516}.
DOI: 10.1126/science.adq9592.
Reports $m_\beta < 0.45$~eV at 90\% C.L.

\bibitem{KamLANDZen2024} S.~Abe \textit{et al.} (KamLAND-Zen Collaboration),
``Search for Majorana Neutrinos with the Complete KamLAND-Zen Dataset,''
Open-access preprint (free): \href{https://arxiv.org/abs/2406.11438}{arXiv:2406.11438} (2024).
Reports $T_{1/2}^{0\nu} > 3.8 \times 10^{26}$~yr at 90\% C.L.,
$\langle m_{\beta\beta}\rangle < 28$--$122$~meV depending on nuclear matrix element.

\bibitem{Planck2018} N.~Aghanim \textit{et al.} (Planck Collaboration),
``Planck 2018 results. VI. Cosmological parameters,''
Astron.\ Astrophys.\ \textbf{641}, A6 (2020).
A\&A is open access; published version free at:
\href{https://www.aanda.org/articles/aa/full_html/2020/09/aa33910-18/aa33910-18.html}{aanda.org/articles/aa/full\_html/2020/09/aa33910-18}.
Preprint: \href{https://arxiv.org/abs/1807.06209}{arXiv:1807.06209}.
DOI: 10.1051/0004-6361/201833910;
Erratum: A\&A \textbf{652}, C4 (2021).

\bibitem{RogersPeiris2021} K.~K.~Rogers and H.~V.~Peiris,
``Strong Bound on Canonical Ultralight Axion Dark Matter from the Lyman-Alpha Forest,''
Phys.\ Rev.\ Lett.\ \textbf{126}, 071302 (2021).
Open-access preprint (free): \href{https://arxiv.org/abs/2007.12705}{arXiv:2007.12705}.
DOI: 10.1103/PhysRevLett.126.071302.
Reports $m_a > 2 \times 10^{-20}$~eV at 95\% C.L.

\bibitem{Koide1981} Y.~Koide,
``Fermion-boson two-body model of quarks and leptons and Cabibbo mixing,''
Lett.\ Nuovo Cimento \textbf{34}, 201--205 (1982).
Inspire-HEP record (free abstract):
\href{https://inspirehep.net/literature/176414}{inspirehep.net/literature/176414}.
DOI: 10.1007/BF02817096;
expanded in
Y.~Koide, Phys.\ Rev.\ D \textbf{28}, 252 (1983),
\href{https://journals.aps.org/prd/abstract/10.1103/PhysRevD.28.252}{APS abstract};
DOI: 10.1103/PhysRevD.28.252;
Phys.\ Lett.\ B \textbf{120}, 161--165 (1983),
\href{https://www.sciencedirect.com/science/article/abs/pii/0370269383906445}{publisher abstract}.
DOI: 10.1016/0370-2693(83)90644-5.

\bibitem{ExistBook} E.~Gable, \textit{Exist: The Universal Ontology}
(2026). Foundational substrate-physical derivations of $c$, $\hbar$,
general relativity, quantum mechanics, Newton's laws, Maxwell's
equations, and the heterotic $E_8 \times E_8$ commitment.

\bibitem{UnreasonableEffectiveness} E.~Gable,
``The unreasonable effectiveness of mathematics, solved:
substrate-physical derivation of the first-principles starting
assumptions,'' companion paper (in preparation, 2026).
Argues that the starting assumptions of
Table~\ref{tab:starting-assumptions} are minimal and substrate-physically
expected. Independent of, and not load-bearing for, the numerical
demonstration of the present paper.

\bibitem{ArnoldCatastrophe} V.~I.~Arnold,
\textit{Catastrophe Theory}, 3rd ed.\
(Springer-Verlag, Berlin, 1992).
Free expository synonym (open access):
\href{https://encyclopediaofmath.org/wiki/Catastrophe_theory}{encyclopediaofmath.org/wiki/Catastrophe\_theory} (Springer/EMS, free).
Springer landing page (free metadata):
\href{https://link.springer.com/book/10.1007/978-3-642-58124-3}{link.springer.com/book/10.1007/978-3-642-58124-3}.
DOI of book: 10.1007/978-3-642-58124-3.
The simply-laced ADE classification of elementary catastrophes
shares its label set with the McKay correspondence used in
Sec.~\ref{sec:framework}, but the two structures are mathematically distinct.

\bibitem{Alcubierre1994} M.~Alcubierre,
``The warp drive: hyper-fast travel within general relativity,''
Classical Quantum Gravity \textbf{11}, L73--L77 (1994).
Open-access preprint (free): \href{https://arxiv.org/abs/gr-qc/0009013}{arXiv:gr-qc/0009013}.
DOI: 10.1088/0264-9381/11/5/001.

\bibitem{MICROSCOPE2022} P.~Touboul \textit{et al.} (MICROSCOPE Collaboration),
``MICROSCOPE Mission: Final Results of the Test of the Equivalence Principle,''
Phys.\ Rev.\ Lett.\ \textbf{129}, 121102 (2022).
Open-access preprint (free): \href{https://arxiv.org/abs/2209.15487}{arXiv:2209.15487}.
DOI: 10.1103/PhysRevLett.129.121102.
Reports E\"otv\"os parameter $\eta(\text{Ti},\text{Pt}) = (-1.5 \pm 2.3_{\rm stat} \pm 1.5_{\rm sys}) \times 10^{-15}$.

\end{thebibliography}

\end{document}
